\documentclass[11pt]{article}
% -------------
% Preambule pro DRILY (cvičné úlohy ve stylu statnic) - MFF UK, UI / Strojové učení
% Sazba: pdflatex (UTF-8, čeština). Samostatné PDF na úroveň (1/2/3).
% -------------
\usepackage[T1]{fontenc}
\usepackage[utf8]{inputenc}
\usepackage{lmodern}
\usepackage[czech]{babel}
\usepackage[a4paper,margin=2.3cm]{geometry}

\usepackage{amsmath,amssymb,mathtools}
\usepackage{textcomp}
\usepackage[nopatch=all]{microtype}
\usepackage{xcolor}
\usepackage{enumitem}
\usepackage{titlesec}
\usepackage[most]{tcolorbox}
\usepackage{booktabs}
\usepackage{listings}
\usepackage{hyperref}

% - seznamy -
\setlist{leftmargin=1.4em, itemsep=2pt, topsep=3pt, parsep=0pt}

% - barvy -
\definecolor{mffblue}{RGB}{27,54,93}
\definecolor{mffaccent}{RGB}{0,90,156}
\definecolor{ulohagray}{RGB}{70,70,75}
\definecolor{resgreen}{RGB}{30,110,60}
\definecolor{gapred}{RGB}{150,30,30}
\definecolor{calorange}{RGB}{180,110,0}
\definecolor{codebg}{RGB}{245,245,245}

\titleformat{\section}{\Large\bfseries\color{mffblue}}{\thesection}{0.6em}{}
\titleformat{\subsection}{\large\bfseries\color{mffaccent}}{\thesubsection}{0.6em}{}
\titlespacing*{\section}{0pt}{1.4em}{0.6em}
\setcounter{secnumdepth}{2}
\setcounter{tocdepth}{1}

\hypersetup{colorlinks=true, linkcolor=mffaccent, urlcolor=mffaccent,
  pdftitle={Cvičné úlohy ke státnicím - Informatika / UI / Strojové učení},
  pdfauthor={Viktor Číhal}, bookmarksnumbered=true}

\renewcommand{\arraystretch}{1.2}
\setlength{\parindent}{0pt}
\setlength{\parskip}{5pt plus 2pt minus 1pt}
\emergencystretch=3em
\hbadness=10000

% - kód (C#); komentáře držíme bez diakritiky kvuli listings+utf8 -
\lstdefinestyle{cs}{
  basicstyle=\ttfamily\footnotesize, backgroundcolor=\color{codebg},
  keywordstyle=\color{mffaccent}\bfseries, commentstyle=\color{resgreen!70!black}\itshape,
  stringstyle=\color{gapred}, showstringspaces=false, columns=fullflexible,
  breaklines=true, frame=single, framerule=0pt, rulecolor=\color{codebg},
  xleftmargin=6pt, xrightmargin=6pt, aboveskip=4pt, belowskip=4pt,
  literate={ě}{{\v e}}1 {š}{{\v s}}1 {č}{{\v c}}1 {ř}{{\v r}}1 {ž}{{\v z}}1
           {á}{{\'a}}1 {í}{{\'\i}}1 {é}{{\'e}}1 {ý}{{\'y}}1 {ů}{{\r u}}1 {ú}{{\'u}}1,
  language=[Sharp]C}
\lstset{style=cs}

% =====================================================================
% Prostředí pro ÚLOHU a ŘEŠENÍ ve stylu statnic.
% =====================================================================

% čítač úloh
\newcounter{uloha}
\renewcommand{\theuloha}{\arabic{uloha}}

% \uloha{okruh}{téma} -> nadpis úlohy ve stylu zadání statnic
% okruh: "společná matematika" / "společná informatika" / "specializace UI-SU"
\newcommand{\ulohahead}[2]{%
  \refstepcounter{uloha}%
  \par\medskip\noindent
  {\large\bfseries\color{ulohagray}Úloha \theuloha.}\quad
  {\small\colorbox{ulohagray!12}{\textbf{#1}}}\quad{\bfseries #2}\par
  \vspace{2pt}}

% zadání: lehce rámovaný blok jako na zkoušce
\newtcolorbox{zadani}{enhanced, breakable, sharp corners=south,
  colback=ulohagray!5, colframe=ulohagray, boxrule=0pt, leftrule=3pt,
  left=8pt, right=8pt, top=4pt, bottom=5pt, before skip=4pt, after skip=4pt}

% podotázky se značkou bodu (kolik bodu daná část typicky nese)
\newcommand{\bod}[1]{\tikz[baseline=-0.55ex]\node[circle,fill=mffaccent,text=white,
  inner sep=1.3pt,font=\bfseries\scriptsize]{#1};}
% bez tikz (rychlejší): kroužek přes textový symbol
\renewcommand{\bod}[1]{{\footnotesize\colorbox{mffaccent}{\textcolor{white}{\bfseries #1\,b}}}}

% řešení
\newtcolorbox{reseni}{enhanced, breakable, colback=resgreen!6, colframe=resgreen,
  boxrule=0.6pt, arc=1.5pt, left=7pt, right=7pt, top=3pt, bottom=4pt,
  before skip=4pt, after skip=8pt, fonttitle=\bfseries\small, coltitle=resgreen,
  title={Řešení}}

% kalibrační poznámka: jak tato úloha sytí podlahu (common>=5 / spec>=7)
\newcommand{\kalibrace}[1]{\par\smallskip\noindent
  {\footnotesize\textcolor{calorange}{\textbf{Kalibrace:}}~\textit{#1}}\par\smallskip}

% past / častá chyba
\newcommand{\past}[1]{\par\smallskip\noindent
  \textbf{\color{gapred}Pozor (častá chyba):}~#1\par}

% info box
\newtcolorbox{infobox}[1][]{enhanced, breakable, colback=mffaccent!5, colframe=mffaccent,
  boxrule=0.7pt, left=6pt, right=6pt, top=4pt, bottom=4pt, arc=2pt,
  before skip=6pt, after skip=6pt, fonttitle=\bfseries, #1}


\begin{document}

\begin{center}
  {\LARGE\bfseries\color{mffblue} Cvičné úlohy ke státnicím \\[2pt]
  Úroveň 3 \textnormal{(maximální pokrytí, realisticky $\sim$97-98\,\%)}}\\[4pt]
  {\small Informatika - Umělá inteligence, zaměření Strojové učení (UI-SU) \quad|\quad MFF UK}\\[2pt]
  {\footnotesize\itshape Kumulativní: úroveň 1 + 2 + rozšíření úrovně 3 (označené [úroveň 3]) + zkoušky nanečisto.}
\end{center}

\begin{infobox}[title={Co přidává úroveň 3 a poctivý strop}]
\footnotesize
Úroveň 2 ($\approx95\,\%$) pokryla nejčastější témata numericky. Úroveň 3 cílí na \emph{zbytkové} riziko: vzácné, ale legální položky požadavku a chyby z formátu/času. Přidává:
\begin{itemize}\setlength{\itemsep}{1pt}
  \item \textbf{Zbylý ocas požadavku} (aby žádný draw nebyl nepokrytý): algebraické struktury (grupy, tělesa), vektorové prostory a lineární zobrazení, posloupnosti a řady, uspořádání (Dilworth), konkrétní rozdělení + nerovnosti, odhady (MLE, intervaly); běh/JIT/linkování, správa zdroju a výjimky, generika, garbage collection; multiclass softmax + křížová validace, POMDP, predikátová rezoluce s unifikací, plánování regresí.
  \item \textbf{Dvě zkoušky nanečisto} (8 otázek se správnou strukturou) a \textbf{tvrdý specializační dril} - trénink formátu, časování a strategie parciálních bodu, bodováno podle reálných podlah.
\end{itemize}
\textbf{Poctivý strop:} z 9 zkoušek dat a velkého oficiálního okruhu je realistická hranice pilného studenta $\sim$\textbf{97,5\,\%} (98\,\% optimisticky); $>99\,\%$ \emph{nelze} z konečné sady poctivě tvrdit. Zbytek je nepředvídatelné frázování a lidská chyba - proti té pomáhají nanečisto-zkoušky.\\
\textbf{Jak na to:} nejdřív zvládni úrovně 1-2; úroveň 3 ber jako pojistku ocasu + ostrý trénink. U zkoušek nanečisto se časuj a poctivě se oboduj podle obou podlah.
\end{infobox}

\tableofcontents
\bigskip

\section{Společná matematika}
\ulohahead{společná matematika}{Vlastní čísla, diagonalizace a symetrické matice}
\begin{zadani}
Je dána reálná matice
\[
A=\begin{pmatrix} 2 & 1\\ 1 & 2 \end{pmatrix}.
\]
\begin{enumerate}[label=(\alph*)]
  \item \bod{1} Definujte vlastní číslo a vlastní vektor matice a charakteristický polynom. Uveďte vztah mezi vlastními čísly a determinantem (resp. regularitou) matice.
  \item \bod{2} Spočtěte vlastní čísla matice $A$ a ke každému najděte vlastní vektor.
  \item \bod{3} Rozhodněte, zda je $A$ diagonalizovatelná. Pokud ano, napište matice $P,D$ tak, že $A=PDP^{-1}$. Vysvětlete, proč je každá reálná symetrická matice (ortogonálně) diagonalizovatelná.
\end{enumerate}
\end{zadani}

\begin{reseni}
\textbf{(a) Definice.} Číslo $\lambda\in\mathbb{R}$ (obecně $\mathbb{C}$) je \emph{vlastní číslo} matice $A$, pokud existuje nenulový vektor $v$ s $Av=\lambda v$; takový $v$ je \emph{vlastní vektor}. Vlastní čísla jsou kořeny \emph{charakteristického polynomu} $p_A(\lambda)=\det(A-\lambda I)$. Platí $\det A=\prod_i \lambda_i$ (součin vlastních čísel s násobností), takže $A$ je regulární právě tehdy, když $0$ není vlastní číslo.

\textbf{(b) Výpočet.}
\[
\det(A-\lambda I)=\det\!\begin{pmatrix}2-\lambda & 1\\ 1 & 2-\lambda\end{pmatrix}=(2-\lambda)^2-1=\lambda^2-4\lambda+3=(\lambda-1)(\lambda-3).
\]
Vlastní čísla: $\lambda_1=3$, $\lambda_2=1$.
\begin{itemize}
  \item $\lambda_1=3$: $(A-3I)=\begin{pmatrix}-1&1\\1&-1\end{pmatrix}$, řešení $v_1=(1,1)^{\!\top}$.
  \item $\lambda_2=1$: $(A-I)=\begin{pmatrix}1&1\\1&1\end{pmatrix}$, řešení $v_2=(1,-1)^{\!\top}$.
\end{itemize}

\textbf{(c) Diagonalizace.} $A$ má dvě různá vlastní čísla, příslušné vlastní vektory jsou nezávislé, takže $A$ je diagonalizovatelná:
\[
P=\begin{pmatrix}1&1\\1&-1\end{pmatrix},\qquad D=\begin{pmatrix}3&0\\0&1\end{pmatrix},\qquad A=PDP^{-1}.
\]
$A$ je navíc symetrická. \emph{Spektrální věta}: každá reálná symetrická matice má reálná vlastní čísla a vlastní vektory příslušné různým vlastním číslům jsou na sebe kolmé, takže je lze znormovat na ortonormální bázi. Pak $A=QDQ^{\!\top}$ s ortogonální $Q$ ($Q^{-1}=Q^{\!\top}$). Zde $v_1\cdot v_2=1\cdot1+1\cdot(-1)=0$, takže $Q=\tfrac{1}{\sqrt2}\begin{psmallmatrix}1&1\\1&-1\end{psmallmatrix}$.
\end{reseni}

\kalibrace{Lineární algebra je nejčastější okruh společné matematiky (~65\,\% zkoušek). Část (a)+(b) = 2 body bezpečně sytí podlahu společných okruhů (cíl $\ge 5/12$). Definice z (a) je nejlevnější bod, nikdy ji nevynechej.}
\past{U diagonalizovatelnosti se neptáme jen ,,jdou spočítat vlastní čísla'', ale zda existuje báze z vlastních vektoru (tj. zda je algebraická = geometrická násobnost). Různá vlastní čísla to zaručí; u násobných je nutno ověřit dimenzi vlastního podprostoru.}

\ulohahead{společná matematika}{Soustavy rovnic, determinant, hodnost, skalární součin a ortogonalizace}
\begin{zadani}
\begin{enumerate}[label=(\alph*)]
  \item \bod{1} Definujte hodnost matice, regulární a inverzní matici a determinant. Uveďte vztah mezi regularitou, determinantem a vlastními čísly.
  \item \bod{2} Spočtěte determinant matice $B=\begin{psmallmatrix}1&2&0\\0&1&3\\1&0&1\end{psmallmatrix}$ a rozhodněte o její regularitě. Popište, jak Gaussova eliminace určí množinu řešení soustavy.
  \item \bod{3} Definujte skalární součin a kolmost. Gram-Schmidtovou ortogonalizací zpracujte vektory $v_1=(1,1,0)$, $v_2=(1,0,1)$ a výsledek znormujte.
\end{enumerate}
\end{zadani}

\begin{reseni}
\textbf{(a)} \emph{Hodnost} $\operatorname{rank}A$ je dimenze řádkového (= sloupcového) prostoru, tj.\ počet nenulových řádku v odstupňovaném tvaru. Matice $A\in\mathbb{R}^{n\times n}$ je \emph{regulární}, má-li hodnost $n$; pak existuje \emph{inverzní} $A^{-1}$ s $AA^{-1}=I$. \emph{Determinant} $\det A$ je multilineární antisymetrická funkce sloupcu s $\det I=1$; platí $\det(AB)=\det A\det B$, $\det A^{\!\top}=\det A$. Ekvivalentně: $A$ regulární $\iff \det A\neq 0 \iff 0$ není vlastní číslo $\iff \operatorname{rank}A=n$.

\textbf{(b)} Laplaceuv rozvoj podle prvního řádku:
\[
\det B = 1\cdot\det\!\begin{psmallmatrix}1&3\\0&1\end{psmallmatrix} -2\cdot\det\!\begin{psmallmatrix}0&3\\1&1\end{psmallmatrix} +0 = 1\cdot 1 -2\cdot(0-3) = 1+6 = 7 \neq 0,
\]
takže $B$ je regulární. \emph{Gaussova eliminace}: soustavu $Ax=b$ převedeme řádkovými úpravami na odstupňovaný tvar. Pak: pokud vznikne řádek $0=c$ s $c\neq0$, soustava nemá řešení; jinak je počet volných proměnných $=$ (počet neznámých) $-\operatorname{rank}A$ a řešení je jednoznačné právě tehdy, když $\operatorname{rank}A=$ počtu neznámých.

\textbf{(c)} \emph{Skalární součin} na reálném vektorovém prostoru $V$ je zobrazení $\langle\cdot,\cdot\rangle:V\times V\to\mathbb{R}$, které je \emph{symetrické} ($\langle x,y\rangle=\langle y,x\rangle$), \emph{bilineární} (lineární v každé složce) a \emph{pozitivně definitní} ($\langle x,x\rangle\ge0$, s rovností jen pro $x=0$). (V komplexním případě se požaduje hermitovskost/sesquilinearita.) Standardní příklad v $\mathbb{R}^n$ je $\langle x,y\rangle=\sum_i x_iy_i$. Indukovaná \emph{norma} $\lVert x\rVert=\sqrt{\langle x,x\rangle}$; vektory jsou \emph{kolmé}, je-li $\langle x,y\rangle=0$. Gram-Schmidt:
\[
u_1=v_1=(1,1,0),\quad u_2=v_2-\frac{\langle v_2,u_1\rangle}{\langle u_1,u_1\rangle}u_1=(1,0,1)-\tfrac12(1,1,0)=\big(\tfrac12,-\tfrac12,1\big).
\]
Kontrola: $\langle u_1,u_2\rangle=\tfrac12-\tfrac12+0=0$. Znormování: $e_1=\tfrac1{\sqrt2}(1,1,0)$, $\lVert u_2\rVert=\sqrt{\tfrac14+\tfrac14+1}=\sqrt{3/2}$, takže $e_2=\tfrac1{\sqrt6}(1,-1,2)$.
\end{reseni}

\kalibrace{Lineární algebra je nejčastější okruh společné matematiky; tato úloha jistí ,,výpočetní'' polovinu (soustavy/determinant/ortogonalizace) vedle úlohy o vlastních číslech. Definice z (a) = jistý 1 bod.}
\past{Determinant počítej rozvojem podle řádku/sloupce s nejvíce nulami. U Gram-Schmidta nezapomeň odečítat \emph{projekci} (s dělením $\langle u_1,u_1\rangle$), ne jen vektor.}

\ulohahead{společná matematika}{Limity, spojitost, derivace a Tayloruv polynom}
\begin{zadani}
\begin{enumerate}[label=(\alph*)]
  \item \bod{1} Definujte limitu funkce v bodě a spojitost funkce. Zformulujte větu o nabývání mezihodnot a větu o nabývání maxima na uzavřeném intervalu.
  \item \bod{2} Spočtěte $\displaystyle\lim_{x\to0}\frac{\sin x - x}{x^3}$. Pro $f(x)=x^3-3x$ najděte lokální extrémy a intervaly monotonie.
  \item \bod{3} Napište Tayloruv polynom (limitní formu) a pomocí něj zdůvodněte výsledek limity z (b). Určete konvexitu/konkavitu $f$.
\end{enumerate}
\end{zadani}

\begin{reseni}
\textbf{(a)} $\lim_{x\to a}f(x)=L$, pokud $\forall\varepsilon>0\,\exists\delta>0:\,0<|x-a|<\delta\Rightarrow|f(x)-L|<\varepsilon$. Funkce je \emph{spojitá} v $a$, je-li $\lim_{x\to a}f(x)=f(a)$. \emph{Věta o mezihodnotě (Bolzano):} je-li $f$ spojitá na $[a,b]$ a $y$ leží mezi $f(a),f(b)$, existuje $c\in[a,b]$ s $f(c)=y$. \emph{Weierstrass:} spojitá funkce na uzavřeném omezeném intervalu nabývá maxima i minima.

\textbf{(b)} Limita je typu $0/0$. l'Hospital třikrát, nebo lépe Taylor: $\sin x = x-\tfrac{x^3}{6}+o(x^3)$, takže
\[
\frac{\sin x-x}{x^3}=\frac{-x^3/6+o(x^3)}{x^3}\to-\tfrac16.
\]
Pro $f(x)=x^3-3x$: $f'(x)=3x^2-3=3(x-1)(x+1)$, nulové body $x=\pm1$. $f'>0$ na $(-\infty,-1)\cup(1,\infty)$ (rostoucí), $f'<0$ na $(-1,1)$ (klesající). V $x=-1$ lokální maximum, v $x=1$ lokální minimum.

\textbf{(c)} \emph{Tayloruv polynom} řádu $n$ v bodě $a$: $T_n(x)=\sum_{k=0}^n\frac{f^{(k)}(a)}{k!}(x-a)^k$, přičemž $f(x)=T_n(x)+o((x-a)^n)$ (limitní/Peanova forma). Rozvoj $\sin x$ kolem $0$ dává přímo čitatel $-x^3/6+o(x^3)$, odkud limita $-1/6$. Konvexita: $f''(x)=6x$, takže $f$ je konkávní na $(-\infty,0)$ a konvexní na $(0,\infty)$; v $x=0$ je inflexní bod.
\end{reseni}

\kalibrace{Analýza je druhý nejčastější okruh matematiky. Definice limity/spojitosti + jeden spočtený limit/extrém spolehlivě dají 2 body do podlahy společných okruhů.}
\past{l'Hospital lze použít jen na typy $\tfrac00$ nebo $\tfrac{\infty}{\infty}$ a po ověření, že limita podílu derivací existuje. Taylor je často rychlejší a méně chybový.}

\ulohahead{společná matematika}{Integrály, primitivní funkce a řady}
\begin{zadani}
\begin{enumerate}[label=(\alph*)]
  \item \bod{1} Definujte primitivní funkci a Riemannuv integrál a uveďte jejich vztah (Newtonova-Leibnizova formule). Definujte součet geometrické řady a uveďte vzorec pro $\sum_{n=0}^\infty ar^n$.
  \item \bod{2} Spočtěte $\displaystyle\int x e^{x}\,dx$ metodou per partes a určete $\displaystyle\int_0^1 x e^{x}\,dx$.
  \item \bod{3} Napište vzorec pro obsah pod grafem a pro objem rotačního tělesa. Vysvětlete odhad součtu řady integrálem (integrální kritérium).
\end{enumerate}
\end{zadani}

\begin{reseni}
\textbf{(a)} $F$ je \emph{primitivní funkce} k $f$ na intervalu, je-li $F'=f$. \emph{Riemannuv integrál} $\int_a^b f$ je společná limita horních a dolních součtu přes zjemňující se dělení (existuje, je-li $f$ např.\ spojitá). \emph{Newton-Leibniz:} je-li $F$ primitivní k spojité $f$, pak $\int_a^b f = F(b)-F(a)$. \emph{Geometrická řada:} $\sum_{n=0}^\infty ar^n=\frac{a}{1-r}$ pro $|r|<1$ (jinak diverguje); částečný součet $\sum_{n=0}^{N}ar^n=a\frac{1-r^{N+1}}{1-r}$.

\textbf{(b)} Per partes $u=x,\ dv=e^x dx$, tedy $du=dx,\ v=e^x$:
\[
\int x e^x\,dx = x e^x-\int e^x\,dx = (x-1)e^x + C.
\]
Určitě: $\int_0^1 x e^x\,dx=\big[(x-1)e^x\big]_0^1=(0)\cdot e-(-1)\cdot 1 = 1.$

\textbf{(c)} \emph{Obsah} pod nezápornou $f$ na $[a,b]$: $S=\int_a^b f(x)\,dx$. \emph{Objem} tělesa vzniklého rotací grafu kolem osy $x$: $V=\pi\int_a^b f(x)^2\,dx$. \emph{Integrální kritérium:} je-li $f$ kladná klesající, pak $\sum_{n=1}^\infty f(n)$ konverguje právě tehdy, když konverguje $\int_1^\infty f(x)\,dx$, a platí odhad $\int_1^{N+1}f \le \sum_{n=1}^{N} f(n) \le f(1)+\int_1^{N} f$.
\end{reseni}

\kalibrace{Integrál + geometrická řada jsou opakované matematické podotázky (geometrická řada byla i samostatná otázka). Vzorec per partes a součet geometrické řady umět zpaměti.}
\past{Per partes: vol $u$ tak, aby se derivováním zjednodušilo (polynom), $dv$ tak, aby šlo integrovat. Geometrická řada konverguje jen pro $|r|<1$ - vždy ověř.}

\ulohahead{společná matematika}{Teorie grafu: souvislost, stromy, barevnost, rovinnost, toky}
\begin{zadani}
\begin{enumerate}[label=(\alph*)]
  \item \bod{1} Definujte graf, stupeň vrcholu, souvislost a komponentu, strom (uveďte aspoň dvě ekvivalentní charakteristiky) a bipartitní graf.
  \item \bod{2} Definujte dobré obarvení a barevnost $\chi(G)$. Uveďte vztah barevnosti a klikovosti. Určete $\chi(C_5)$ (kružnice délky 5) a zdůvodněte.
  \item \bod{3} Zformulujte Eulerovu formuli pro rovinný graf a odvoďte horní mez počtu hran. Definujte síť, tok a řez a uveďte větu o max.\ toku a min.\ řezu.
\end{enumerate}
\end{zadani}

\begin{reseni}
\textbf{(a)} \emph{Graf} $G=(V,E)$, $E\subseteq\binom{V}{2}$. \emph{Stupeň} $\deg(v)$ = počet hran u $v$ (princip podání rukou: $\sum_v\deg v=2|E|$). $G$ je \emph{souvislý}, existuje-li cesta mezi každou dvojicí vrcholu; \emph{komponenta} = maximální souvislý podgraf. \emph{Strom} = souvislý graf bez kružnice; ekvivalentně: souvislý s $|E|=|V|-1$; nebo: mezi každými dvěma vrcholy existuje právě jedna cesta. \emph{Bipartitní} graf: $V$ lze rozdělit na dvě části tak, že každá hrana vede mezi nimi (ekvivalentně: neobsahuje lichou kružnici).

\textbf{(b)} \emph{Dobré obarvení} $k$ barvami: zobrazení $V\to\{1,\dots,k\}$ tak, že sousední vrcholy mají různou barvu. \emph{Barevnost} $\chi(G)$ = nejmenší takové $k$. Platí $\chi(G)\ge\omega(G)$ (klikovost; klika potřebuje tolik barev, kolik má vrcholu), ale rovnost obecně neplatí. $C_5$: jako lichá kružnice není bipartitní, takže $\chi>2$; třemi barvami obarvit lze, tedy $\chi(C_5)=3$.

\textbf{(c)} \emph{Eulerova formule:} pro souvislý rovinný graf s $V$ vrcholy, $E$ hranami a $F$ stěnami platí $V-E+F=2$. U jednoduchého grafu s $V\ge3$ je každá stěna ohraničena $\ge3$ hranami a každá hrana sousedí se 2 stěnami, takže $2E\ge3F$; dosazením $F=2-V+E$ dostaneme $E\le 3V-6$. (U grafu s mosty je třeba argument doplnit, závěr ale platí.) \emph{Síť} = orientovaný graf s kapacitami $c$, zdrojem $s$ a stokem $t$; \emph{tok} $f$ splňuje $0\le f\le c$ a Kirchhoffuv zákon (zachování v každém vnitřním vrcholu); \emph{řez} = rozdělení vrcholu na $S\ni s$ a $T\ni t$, jeho kapacita = součet kapacit hran z $S$ do $T$. \emph{Věta:} velikost maximálního toku se rovná kapacitě minimálního řezu (max-flow = min-cut).
\end{reseni}

\kalibrace{Grafy jsou stabilní okruh matematiky (barevnost, eulerovské/rovinné grafy, souvislost se v zadáních opakují). Definice + jedno odvození (např.\ $E\le3V-6$) = 2 body.}
\past{$\chi\ge\omega$ je jen nerovnost - existují grafy bez trojúhelníku s libovolně velkou barevností. Eulerova formule platí pro \emph{souvislé} rovinné grafy.}

\ulohahead{společná matematika}{Pravděpodobnost a statistika: Bayes, střední hodnota, rozptyl, CLT, testování}
\begin{zadani}
\begin{enumerate}[label=(\alph*)]
  \item \bod{1} Definujte pravděpodobnostní prostor, podmíněnou pravděpodobnost a nezávislost jevu. Zformulujte Bayesovu větu.
  \item \bod{2} Definujte střední hodnotu a rozptyl náhodné veličiny. Uveďte linearitu střední hodnoty a vzorec pro rozptyl součtu. Spočtěte $\mathbb{E}X$ a $\operatorname{var}X$ pro hod férovou kostkou.
  \item \bod{3} Zformulujte zákon velkých čísel a centrální limitní větu. Popište testování hypotéz: hladina významnosti, chyba 1.\ a 2.\ druhu.
\end{enumerate}
\end{zadani}

\begin{reseni}
\textbf{(a)} \emph{Pravděpodobnostní prostor} $(\Omega,\mathcal{F},P)$: množina elementárních jevu $\Omega$, $\sigma$-algebra jevu $\mathcal{F}$, míra $P$ s $P(\Omega)=1$. \emph{Podmíněná} $P(A\mid B)=\frac{P(A\cap B)}{P(B)}$ pro $P(B)>0$. Jevy \emph{nezávislé}, je-li $P(A\cap B)=P(A)P(B)$. \emph{Bayes:} $P(A\mid B)=\dfrac{P(B\mid A)P(A)}{P(B)}$, kde dle \emph{úplné pravděpodobnosti} $P(B)=\sum_i P(B\mid A_i)P(A_i)$ pro rozklad $\{A_i\}$ množiny $\Omega$ (po dvou disjunktní, pokrývající, $P(A_i)>0$).

\textbf{(b)} $\mathbb{E}X=\sum_x x\,P(X=x)$ (diskrétně), resp.\ $\mathbb{E}X=\int_{-\infty}^{\infty}x\,f_X(x)\,dx$ (spojitě). \emph{Linearita}: $\mathbb{E}(aX+bY)=a\mathbb{E}X+b\mathbb{E}Y$ (vždy, i pro závislé). $\operatorname{var}X=\mathbb{E}(X-\mathbb{E}X)^2=\mathbb{E}X^2-(\mathbb{E}X)^2$; $\operatorname{var}(aX)=a^2\operatorname{var}X$. Obecně $\operatorname{var}(X+Y)=\operatorname{var}X+\operatorname{var}Y+2\operatorname{cov}(X,Y)$, takže pro \emph{nezávislé} (kde $\operatorname{cov}=0$) je $\operatorname{var}(X+Y)=\operatorname{var}X+\operatorname{var}Y$. Kostka: $\mathbb{E}X=\tfrac{1+\dots+6}{6}=3{,}5$, $\mathbb{E}X^2=\tfrac{1+4+9+16+25+36}{6}=\tfrac{91}{6}$, $\operatorname{var}X=\tfrac{91}{6}-\tfrac{49}{4}=\tfrac{35}{12}\approx2{,}92$.

\textbf{(c)} \emph{ZVČ:} prumer $\bar X_n$ nezávislých stejně rozdělených veličin konverguje k $\mathbb{E}X$. \emph{CLT:} pro nezávislé stejně rozdělené se střední hodnotou $\mu$ a rozptylem $\sigma^2$ platí $\frac{\sqrt n(\bar X_n-\mu)}{\sigma}\xrightarrow{d}\mathcal{N}(0,1)$, tj.\ normovaný prumer má asymptoticky normální rozdělení. \emph{Testování:} nulová hypotéza $H_0$ vs.\ alternativa; \emph{hladina významnosti} $\alpha$ = max.\ přípustná pravděpodobnost \emph{chyby 1.\ druhu} (zamítnout platné $H_0$); \emph{chyba 2.\ druhu} $\beta$ = nezamítnout neplatné $H_0$; síla testu $=1-\beta$.
\end{reseni}

\kalibrace{Pravděpodobnost a teorie informace se v matematice objevují středně často; definice Bayes/EX/var + jeden výpočet je bezpečná záloha. Tyto pojmy navíc podpírají ML specializaci (regrese, naivní pravděpodobnostní úvahy).}
\past{Linearita střední hodnoty platí \emph{i pro závislé} veličiny; aditivita rozptylu jen pro \emph{nezávislé} (jinak $+2\operatorname{cov}$). Hladina významnosti řídí chybu 1.\ druhu, ne 2.}

\ulohahead{společná matematika}{Logika: normální tvary, sémantika, rezoluce}
\begin{zadani}
\begin{enumerate}[label=(\alph*)]
  \item \bod{1} Vysvětlete základní pojmy syntaxe výrokové a predikátové logiky (formule, otevřená/uzavřená formule). Definujte CNF, DNF a PNF.
  \item \bod{2} Definujte model, splnitelnost, tautologii a důsledek. Převeďte $(p\to q)$ na CNF a rozhodněte, zda je $(p\to q)\wedge p\to q$ tautologie.
  \item \bod{3} Popište rezoluci jako dokazovací metodu. Rezolucí ukažte, že množina klauzulí $\{p\vee q,\ \neg p,\ \neg q\}$ je nesplnitelná.
\end{enumerate}
\end{zadani}

\begin{reseni}
\textbf{(a)} \emph{Formule} se tvoří z prvovýroku (resp.\ atomických formulí s predikáty, termy a kvantifikátory) logickými spojkami. Ve výrokové logice není kvantifikace; v predikátové je formule \emph{otevřená}, má-li volnou proměnnou, a \emph{uzavřená} (sentence), nemá-li žádnou. \emph{CNF} = konjunkce disjunkcí literálu (klauzulí); \emph{DNF} = disjunkce konjunkcí literálu; \emph{PNF} (prenexní) = všechny kvantifikátory na začátku, za nimi otevřené jádro.

\textbf{(b)} \emph{Model} teorie/formule = ohodnocení (struktura), v němž je pravdivá. Formule je \emph{splnitelná}, má-li model; \emph{tautologie}, je-li pravdivá v každém ohodnocení; $\varphi$ je \emph{důsledek} teorie $T$ ($T\models\varphi$), platí-li ve všech modelech $T$. Převod: $p\to q\equiv\neg p\vee q$ (už je CNF, jediná klauzule). Výraz $((p\to q)\wedge p)\to q$ je \emph{modus ponens}; pravdivostní tabulkou (4 řádky) vyjde vždy $1$, je to tautologie.

\textbf{(c)} \emph{Rezoluce}: pracuje s klauzulemi (CNF); rezolventa dvojice klauzulí $C_1\vee\ell$ a $C_2\vee\neg\ell$ je $C_1\vee C_2$. Množina je nesplnitelná, lze-li z ní odvodit prázdnou klauzuli $\square$. Zde:
\[
\{p\vee q\},\ \{\neg p\}\ \Rightarrow\ \{q\};\qquad \{q\},\ \{\neg q\}\ \Rightarrow\ \square.
\]
Odvodili jsme prázdnou klauzuli, takže $\{p\vee q,\neg p,\neg q\}$ je nesplnitelná. (Doplňkově: věty o kompaktnosti a úplnosti zaručují, že rezoluce je úplná dokazovací metoda - detail je nad rámec úrovně 1.)
\end{reseni}

\kalibrace{Logika je vzácnější matematický okruh, ale definice (CNF/DNF, model, splnitelnost) jsou levné body a rezoluce/DPLL se prolínají i se specializací ,,Základy UI''. Tady stačí 2-bodová záloha.}
\past{Důsledek ($\models$) je sémantický pojem (přes modely), dokazatelnost ($\vdash$) syntaktický (přes odvození); věta o úplnosti je spojuje. Nezaměňuj splnitelnost (existuje model) s tautologií (všechny modely).}

\ulohahead{společná matematika}{Diskrétní matematika: relace, kombinatorika, inkluze-exkluze}
\begin{zadani}
\begin{enumerate}[label=(\alph*)]
  \item \bod{1} Definujte vlastnosti binární relace (reflexivita, symetrie, antisymetrie, tranzitivita), ekvivalenci a rozkladové třídy a částečné uspořádání (řetězec, antiřetězec).
  \item \bod{2} Kolik je prostých a kolik všech funkcí z $m$-prvkové do $n$-prvkové množiny? Zformulujte binomickou větu a princip inkluze a exkluze.
  \item \bod{3} Pomocí inkluze-exkluze odvoďte počet surjekcí z $m$-prvkové na $n$-prvkovou množinu a počet permutací bez pevného bodu (problém šatnářky).
\end{enumerate}
\end{zadani}

\begin{reseni}
\textbf{(a)} Relace $R\subseteq X\times X$ je \emph{reflexivní} ($\forall x: xRx$), \emph{symetrická} ($xRy\Rightarrow yRx$), \emph{antisymetrická} ($xRy\wedge yRx\Rightarrow x=y$), \emph{tranzitivní} ($xRy\wedge yRz\Rightarrow xRz$). \emph{Ekvivalence} = reflexivní + symetrická + tranzitivní; její \emph{rozkladové třídy} tvoří rozklad $X$. \emph{Částečné uspořádání} = reflexivní + antisymetrické + tranzitivní; \emph{řetězec} = po dvou porovnatelné prvky, \emph{antiřetězec} = po dvou neporovnatelné.

\textbf{(b)} Všech funkcí $X\to Y$ je $n^m$; \emph{prostých} (pro $m\le n$) je $n(n-1)\cdots(n-m+1)=\frac{n!}{(n-m)!}$. \emph{Binomická věta:} $(x+y)^n=\sum_{k=0}^n\binom{n}{k}x^k y^{n-k}$. \emph{Inkluze-exkluze:} $\big|\bigcup_{i=1}^n A_i\big|=\sum_i|A_i|-\sum_{i<j}|A_i\cap A_j|+\cdots+(-1)^{n+1}|A_1\cap\cdots\cap A_n|$.

\textbf{(c)} \emph{Surjekce} z $m$ na $n$: od všech $n^m$ funkcí odečteme ty, které vynechají aspoň jeden prvek cíle. Označíme-li $A_i$ funkce vynechávající $i$-tý prvek, dostaneme inkluzí-exkluzí
\[
\#\text{surjekcí}=\sum_{i=0}^{n}(-1)^i\binom{n}{i}(n-i)^m.
\]
\emph{Problém šatnářky} (derangement): permutace $\{1,\dots,n\}$ bez pevného bodu; $A_i$ = permutace s $\pi(i)=i$, $|A_{i_1}\cap\cdots\cap A_{i_k}|=(n-k)!$, takže
\[
D_n=\sum_{k=0}^n(-1)^k\binom{n}{k}(n-k)!=n!\sum_{k=0}^n\frac{(-1)^k}{k!}\approx \frac{n!}{e}.
\]
(Hallova věta o SRR a párování v bipartitním grafu je příbuzné téma - viz vyšší úroveň.)
\end{reseni}

\kalibrace{Diskrétní matematika dává levné definiční body (relace, ekvivalence) a jeden klasický I-E výpočet. Inkluze-exkluze je opakovaný typ.}
\past{Antisymetrie není ,,opak symetrie''. Počet surjekcí není $\binom{n}{?}$ trik - jde přes inkluzi-exkluzi; nezapomeň člen $i=0$ (= $n^m$).}

\ulohahead{společná matematika \textnormal{\footnotesize[úroveň 2]}}{Pozitivně definitní matice, Sylvester a Choleského rozklad}
\begin{zadani}
Matice $A=\begin{psmallmatrix}4&2\\2&3\end{psmallmatrix}$.
\begin{enumerate}[label=(\alph*)]
  \item \bod{1} Definujte pozitivně definitní matici a uveďte Sylvestrovo kritérium a vztah k vlastním číslum.
  \item \bod{2} Rozhodněte, zda je $A$ pozitivně definitní.
  \item \bod{3} Spočtěte Choleského rozklad $A=LL^{\!\top}$.
\end{enumerate}
\end{zadani}

\begin{reseni}
\textbf{(a)} Symetrická $A$ je \emph{pozitivně definitní}, pokud $x^{\!\top}Ax>0$ pro každé $x\neq0$. Ekvivalentně: všechna vlastní čísla jsou kladná; nebo (\emph{Sylvestrovo kritérium}) všechny hlavní rohové minory jsou kladné. Pak existuje jednoznačný \emph{Choleského rozklad} $A=LL^{\!\top}$ s dolní trojúhelníkovou $L$ s kladnou diagonálou.

\textbf{(b)} Rohové minory: $4>0$ a $\det A=4\cdot3-2\cdot2=8>0$. Obě kladné $\Rightarrow$ $A$ je pozitivně definitní. (Kontrola přes vlastní čísla: stopa $7$, determinant $8$, $\lambda=\tfrac{7\pm\sqrt{17}}{2}\approx5{,}56;\,1{,}44>0$.)

\textbf{(c)} $A=LL^{\!\top}$, $L=\begin{psmallmatrix}l_{11}&0\\l_{21}&l_{22}\end{psmallmatrix}$: $l_{11}=\sqrt{4}=2$; $l_{21}=a_{21}/l_{11}=2/2=1$; $l_{22}=\sqrt{a_{22}-l_{21}^2}=\sqrt{3-1}=\sqrt2$. Tedy
\[
L=\begin{psmallmatrix}2&0\\1&\sqrt2\end{psmallmatrix},\qquad LL^{\!\top}=\begin{psmallmatrix}4&2\\2&3\end{psmallmatrix}=A.
\]
\end{reseni}

\begin{jadro}
PD $\iff$ vlastní čísla $>0$ $\iff$ všechny rohové minory $>0$ (Sylvester). Cholesky: $l_{11}=\sqrt{a_{11}}$, $l_{21}=a_{21}/l_{11}$, $l_{22}=\sqrt{a_{22}-l_{21}^2}$.
\end{jadro}

\kalibrace{Pozitivní definitnost + Cholesky je doslova v požadavku a opakuje se v lineární algebře (skalární součin, pozitivně definitní matice). Sylvester + jeden rozklad = 2-3 body.}
\past{Sylvester vyžaduje \emph{rohové} (leading) minory, ne libovolné. Pozitivní semidefinitní připouští nulové vlastní číslo (minory $\ge0$).}

\ulohahead{společná matematika \textnormal{\footnotesize[úroveň 2]}}{Určité integrály: obsah, objem, substituce}
\begin{zadani}
\begin{enumerate}[label=(\alph*)]
  \item \bod{1} Spočtěte $\displaystyle\int_0^1 (3x^2+2x)\,dx$.
  \item \bod{2} Spočtěte obsah plochy mezi křivkami $y=x$ a $y=x^2$ na $[0,1]$ a objem tělesa vzniklého rotací $y=x$ kolem osy $x$ na $[0,1]$.
  \item \bod{3} Spočtěte $\displaystyle\int_0^1 x\,e^{x^2}\,dx$ substitucí.
\end{enumerate}
\end{zadani}

\begin{reseni}
\textbf{(a)} $\int_0^1(3x^2+2x)\,dx=\big[x^3+x^2\big]_0^1=1+1=2.$

\textbf{(b)} Na $[0,1]$ je $x\ge x^2$, takže obsah $=\int_0^1 (x-x^2)\,dx=\big[\tfrac{x^2}{2}-\tfrac{x^3}{3}\big]_0^1=\tfrac12-\tfrac13=\tfrac16$. Objem rotace $y=x$ kolem osy $x$: $V=\pi\int_0^1 x^2\,dx=\pi\big[\tfrac{x^3}{3}\big]_0^1=\tfrac{\pi}{3}$.

\textbf{(c)} Substituce $u=x^2$, $du=2x\,dx$, tedy $x\,dx=\tfrac12 du$; meze $x{:}0\to1$ dají $u{:}0\to1$:
\[
\int_0^1 x e^{x^2}\,dx=\tfrac12\int_0^1 e^{u}\,du=\tfrac12\big[e^u\big]_0^1=\tfrac{e-1}{2}\approx0{,}859.
\]
\end{reseni}

\begin{jadro}
Obsah $=\int_a^b (\text{horní}-\text{dolní})\,dx$; objem rotace kolem osy $x$ $=\pi\int_a^b f(x)^2 dx$. Substituce: zaveď $u$, přepiš $dx$ i meze. $\int_0^1 x e^{x^2}=\frac{e-1}{2}$.
\end{jadro}

\kalibrace{Obsah/objem a substituce jsou opakované výpočetní podotázky analýzy (~45\,\%). Jeden bezchybný integrál = 2 body do podlahy společných okruhu.}
\past{U obsahu mezi křivkami integruj \emph{rozdíl} (horní minus dolní) a hlídej, která je nahoře. U substituce přepiš i meze, jinak musíš vracet zpět do $x$.}

\ulohahead{společná matematika \textnormal{\footnotesize[úroveň 2]}}{Pravděpodobnost numericky: Bayes a centrální limitní věta}
\begin{zadani}
\begin{enumerate}[label=(\alph*)]
  \item \bod{1} Nemoc má prevalenci $P(D)=0{,}01$. Test má $P(+\mid D)=0{,}99$ a $P(+\mid \neg D)=0{,}05$. Spočtěte $P(D\mid +)$.
  \item \bod{2} Interpretujte výsledek (proč je tak nízký).
  \item \bod{3} Pomocí CLT odhadněte $P(\text{součet }100\text{ hodu kostkou}>380)$. (Pro jednu kostku $\mathbb{E}X=3{,}5$, $\operatorname{var}X=\tfrac{35}{12}$.)
\end{enumerate}
\end{zadani}

\begin{reseni}
\textbf{(a)} Bayes:
\[
P(D\mid+)=\frac{P(+\mid D)P(D)}{P(+\mid D)P(D)+P(+\mid\neg D)P(\neg D)}=\frac{0{,}99\cdot0{,}01}{0{,}99\cdot0{,}01+0{,}05\cdot0{,}99}=\frac{0{,}0099}{0{,}0594}\approx0{,}167.
\]

\textbf{(b)} I při velmi přesném testu je posterior jen $\approx16{,}7\,\%$, protože nemoc je vzácná: falešně pozitivních ze zdravé většiny ($0{,}05\cdot0{,}99\approx0{,}0495$) je víc než pravdivě pozitivních ($0{,}0099$). Nízká \emph{prevalence} (base rate) sráží posterior.

\textbf{(c)} Součet $S=\sum_{i=1}^{100}X_i$: $\mathbb{E}S=100\cdot3{,}5=350$, $\operatorname{var}S=100\cdot\tfrac{35}{12}\approx291{,}7$, $\sigma_S\approx17{,}08$. Dle CLT je $S$ přibližně normální, takže $z=\frac{380-350}{17{,}08}\approx1{,}76$ a $P(S>380)\approx P(Z>1{,}76)\approx0{,}039$.
\end{reseni}

\begin{jadro}
Bayes: posterior $\propto$ věrohodnost $\times$ prior; u vzácných jevu pozor na base rate. CLT: součet/průměr nezávislých je přibližně $\mathcal{N}$; $\mathbb{E}S=n\mu$, $\operatorname{var}S=n\sigma^2$, standardizuj $z=\frac{x-\mathbb{E}S}{\sigma_S}$.
\end{jadro}

\kalibrace{Bayes numericky a CLT aplikace jsou opakované (pravděpodobnost ~25\,\%, CLT se zkoušela 2025-léto). Dosazení do vzorce = 2 body.}
\past{Rozptyl součtu je $n\sigma^2$, směrodatná odchylka $\sqrt n\,\sigma$ (ne $n\sigma$). U Bayese nezapomeň na jmenovatel přes úplnou pravděpodobnost.}

\ulohahead{společná matematika \textnormal{\footnotesize[úroveň 2]}}{Eulerova formule (důkaz) a Mengerova věta}
\begin{zadani}
\begin{enumerate}[label=(\alph*)]
  \item \bod{1} Zformulujte Eulerovu formuli pro souvislý rovinný graf a naznačte její důkaz indukcí.
  \item \bod{2} Odvoďte $E\le 3V-6$ a pomocí toho ukažte, že $K_5$ není rovinný.
  \item \bod{3} Zformulujte hranovou a vrcholovou verzi Mengerovy věty.
\end{enumerate}
\end{zadani}

\begin{reseni}
\textbf{(a)} Pro souvislý rovinný graf platí $V-E+F=2$. \emph{Důkaz indukcí podle počtu hran}: kostra (strom) má $E=V-1$ a $F=1$, takže $V-(V-1)+1=2$ platí. Přidáním hrany, která uzavře kružnici, vznikne nová stěna: $E$ i $F$ vzrostou o 1, takže $V-E+F$ se nezmění. Indukcí platí pro každý souvislý rovinný graf.

\textbf{(b)} U jednoduchého grafu s $V\ge3$ ohraničuje každou stěnu $\ge3$ hran a každá hrana sousedí se 2 stěnami, tedy $2E\ge3F$. Dosazením $F=2-V+E$: $2E\ge3(2-V+E)=6-3V+3E$, odkud $E\le3V-6$. Pro $K_5$: $V=5$, $E=\binom52=10$, ale $3V-6=9<10$, takže $K_5$ nesplňuje nutnou podmínku rovinnosti $\Rightarrow$ není rovinný.

\textbf{(c)} \emph{Mengerova věta (hranová):} minimální počet hran, jejichž odebrání rozpojí vrcholy $u,v$, se rovná maximálnímu počtu \emph{hranově disjunktních} cest mezi $u,v$. \emph{Vrcholová verze:} minimální počet vrcholu (různých od $u,v$), jejichž odebrání rozpojí $u,v$, se rovná maximálnímu počtu \emph{vnitřně vrcholově disjunktních} cest mezi $u,v$.
\end{reseni}

\begin{jadro}
$V-E+F=2$ (souvislý rovinný), důkaz indukcí přes hrany (kružnicová hrana přidá stěnu). Důsledek $E\le3V-6$ vylučuje $K_5$. Menger: min.\ řez $=$ max.\ počet disjunktních cest (hranová/vrcholová verze).
\end{jadro}

\kalibrace{Eulerova formule, rovinnost a Menger jsou opakovaný grafový okruh (~45\,\%); úroveň 2 přidává \emph{důkaz}, který zkouška může chtít. Náčrt důkazu = 3.\ bod.}
\past{$E\le3V-6$ platí pro jednoduché grafy $V\ge3$ (bipartitní: $E\le2V-4$). Mengerova věta je grafový protějšek max-flow/min-cut.}

\ulohahead{společná matematika \textnormal{\footnotesize[úroveň 2]}}{Tablo metoda a věty o kompaktnosti a úplnosti}
\begin{zadani}
\begin{enumerate}[label=(\alph*)]
  \item \bod{1} Popište princip tablo metody (důkaz sporem, rozkladová pravidla, uzavřená větev).
  \item \bod{2} Tablo metodou ukažte, že $((p\to q)\wedge p)\to q$ je tautologie.
  \item \bod{3} Zformulujte větu o kompaktnosti a větu o úplnosti a vysvětlete jejich význam.
\end{enumerate}
\end{zadani}

\begin{reseni}
\textbf{(a)} \emph{Tablo} dokazuje tautologii sporem: předpokládá, že formule \emph{neplatí}, a rozkládá ji \emph{rozkladovými pravidly} na složky (konjunktivní pravidla větev prodlouží, disjunktivní ji rozvětví). Větev se \emph{uzavře}, obsahuje-li formuli i její negaci. Uzavřou-li se všechny větve, je výchozí předpoklad sporný, takže formule je tautologie.

\textbf{(b)} Předpokládáme nepravdivost: $\neg\big(((p\to q)\wedge p)\to q\big)$, což znamená $(p\to q)\wedge p$ pravdivé a $q$ nepravdivé. Z konjunkce: $p\to q$ pravdivé, $p$ pravdivé, $\neg q$. Z $p\to q$ (disjunkce $\neg p\vee q$) se větev rozdělí: větev s $\neg p$ se uzavře (máme $p$ i $\neg p$); větev s $q$ se uzavře (máme $q$ i $\neg q$). Všechny větve uzavřené $\Rightarrow$ tautologie.

\textbf{(c)} \emph{Věta o kompaktnosti}: teorie má model právě tehdy, když má model každá její konečná podmnožina. \emph{Věta o úplnosti}: $T\models\varphi$ (sémantický důsledek) právě tehdy, když $T\vdash\varphi$ (existuje formální důkaz). Význam: úplnost spojuje pravdivost ve všech modelech s odvoditelností v kalkulu (dokazatelnost ,,dohoní'' platnost); kompaktnost umožňuje přenášet vlastnosti z konečných částí na celek (časté u nekonečných teorií).
\end{reseni}

\begin{jadro}
Tablo = důkaz sporem: rozkládej negaci cíle, uzavři větev při $X$ i $\neg X$; všechny uzavřené $\Rightarrow$ tautologie. Kompaktnost: model existuje $\iff$ pro každou konečnou podmnožinu. Úplnost: $\models\,\iff\,\vdash$.
\end{jadro}

\kalibrace{Logika (znění vět o kompaktnosti/úplnosti, tablo/rezoluce) je v požadavku a je levná na definice. Úroveň 2 přidává jeden vyřešený tablo důkaz.}
\past{Tablo dokazuje sporem - začínáš \emph{negací} dokazované formule. Úplnost ($\models\iff\vdash$) nezaměňuj s korektností (jen $\vdash\Rightarrow\models$).}

\ulohahead{společná matematika \textnormal{\footnotesize[úroveň 3]}}{Algebraické struktury: grupy, podgrupy, tělesa, konečná tělesa}
\begin{zadani}
\begin{enumerate}[label=(\alph*)]
  \item \bod{1} Definujte grupu a podgrupu. Co je permutace a grupa permutací $S_n$?
  \item \bod{2} Definujte těleso. Uveďte příklady (i konečné) a vysvětlite charakteristiku tělesa.
  \item \bod{3} Zformulujte Lagrangeovu větu a uveďte, kdy je $\mathbb{Z}_n$ těleso.
\end{enumerate}
\end{zadani}

\begin{reseni}
\textbf{(a)} \emph{Grupa} $(G,\cdot)$ je množina s binární operací, která je asociativní ($a(bc)=(ab)c$), má \emph{neutrální prvek} $e$ ($ae=ea=a$) a ke každému $a$ \emph{inverzní} $a^{-1}$ ($aa^{-1}=e$). Je-li navíc komutativní, je \emph{abelovská}. \emph{Podgrupa} $H\le G$ je podmnožina, která je sama grupou vuči téže operaci (uzavřená na operaci i inverzy, obsahuje $e$). \emph{Permutace} množiny je bijekce na sebe; všechny permutace $\{1,\dots,n\}$ se skládáním tvoří grupu $S_n$ řádu $n!$ (obecně nekomutativní).

\textbf{(b)} \emph{Těleso} $(T,+,\cdot)$ je množina se dvěma operacemi, kde $(T,+)$ je abelovská grupa (s neutrálním $0$), $(T\setminus\{0\},\cdot)$ je abelovská grupa (s neutrálním $1$) a platí distributivita. Příklady: $\mathbb{Q},\mathbb{R},\mathbb{C}$ (nekonečná) a $\mathbb{Z}_p$ pro prvočíslo $p$ (konečné). \emph{Charakteristika} je nejmenší $k>0$ s $\underbrace{1+\dots+1}_{k}=0$ (nebo $0$, pokud takové neexistuje); u konečného tělesa je to vždy prvočíslo $p$ a těleso má $p^m$ prvku.

\textbf{(c)} \emph{Lagrangeova věta}: v konečné grupě řád každé podgrupy dělí řád grupy ($|H|\mid|G|$); důsledek: řád každého prvku dělí $|G|$. $\mathbb{Z}_n$ (zbytky modulo $n$) je \emph{těleso právě tehdy, když je $n$ prvočíslo} - jen tehdy má každý nenulový prvek multiplikativní inverz (jinak existují dělitelé nuly).
\end{reseni}

\begin{jadro}
Grupa = asociativní operace + neutrál + inverzy. Těleso = $(T,+)$ a $(T\setminus\{0\},\cdot)$ abelovské grupy + distributivita. $\mathbb{Z}_p$ těleso $\iff p$ prvočíslo. Lagrange: $|H|\mid|G|$.
\end{jadro}

\kalibrace{Algebraické struktury (grupy, tělesa, konečná tělesa) jsou v požadavku, ale v moderních zkouškách vzácné - úroveň 3 je pokrývá, aby žádný draw společné matematiky nebyl nepokrytý.}
\past{Těleso vyžaduje inverzy i pro \emph{násobení} (kromě 0) - proto $\mathbb{Z}_4$ není těleso (2 nemá inverz). Charakteristika konečného tělesa je vždy prvočíslo.}

\ulohahead{společná matematika \textnormal{\footnotesize[úroveň 3]}}{Vektorové prostory a lineární zobrazení}
\begin{zadani}
\begin{enumerate}[label=(\alph*)]
  \item \bod{1} Definujte vektorový prostor, lineární nezávislost, bázi, dimenzi a souřadnice vektoru.
  \item \bod{2} Zformulujte Steinitzovu větu o výměně. Definujte lineární zobrazení, jeho obraz a jádro.
  \item \bod{3} Zformulujte větu o dimenzi (rank-nullity) a vysvětlete izomorfismus vektorových prostoru.
\end{enumerate}
\end{zadani}

\begin{reseni}
\textbf{(a)} \emph{Vektorový prostor} nad tělesem $T$ je množina $V$ se sčítáním (abelovská grupa) a násobením skalárem $T\times V\to V$ splňujícím distributivitu, asociativitu a $1\cdot v=v$. Vektory $v_1,\dots,v_k$ jsou \emph{lineárně nezávislé}, pokud $\sum c_iv_i=0$ implikuje všechna $c_i=0$. \emph{Báze} = lineárně nezávislý systém generátoru (každý vektor je jejich jednoznačnou lineární kombinací); \emph{dimenze} = počet prvku báze. \emph{Souřadnice} vektoru v bázi jsou koeficienty této kombinace.

\textbf{(b)} \emph{Steinitzova věta o výměně}: je-li $\{v_1,\dots,v_k\}$ lineárně nezávislá a $\{w_1,\dots,w_m\}$ systém generátoru, pak $k\le m$ a lze $k$ generátoru vyměnit za $v_i$ tak, že vznikne opět systém generátoru (důsledek: všechny báze mají stejnou velikost). \emph{Lineární zobrazení} $f:V\to W$ splňuje $f(au+bv)=af(u)+bf(v)$. \emph{Obraz} $\operatorname{Im}f=\{f(v):v\in V\}$ (podprostor $W$), \emph{jádro} $\operatorname{Ker}f=\{v:f(v)=0\}$ (podprostor $V$).

\textbf{(c)} \emph{Věta o dimenzi (rank-nullity)}: pro lineární $f:V\to W$ s konečnou dimenzí $V$ platí $\dim\operatorname{Im}f+\dim\operatorname{Ker}f=\dim V$ (hodnost + defekt). \emph{Izomorfismus} je bijektivní lineární zobrazení; dva prostory nad týmž tělesem jsou izomorfní právě tehdy, mají-li stejnou (konečnou) dimenzi - každý $n$-rozměrný prostor je izomorfní $T^n$ (přes souřadnice v bázi).
\end{reseni}

\begin{jadro}
Báze = nezávislý systém generátoru; dimenze = její velikost (všechny báze stejně velké, Steinitz). $f$ lineární: $\operatorname{Im}f\le W$, $\operatorname{Ker}f\le V$, $\dim\operatorname{Im}f+\dim\operatorname{Ker}f=\dim V$. Izomorfní $\iff$ stejná dimenze.
\end{jadro}

\kalibrace{Vektorové prostory a lineární zobrazení (báze, dimenze, jádro/obraz, rank-nullity) jsou jádro lineární algebry; úroveň 3 doplňuje abstraktní část vedle výpočetních úloh z úrovní 1-2.}
\past{Báze není ,,libovolná generující množina'' - musí být i nezávislá. Rank-nullity sčítá dimenze obrazu a jádra do dimenze \emph{definičního oboru}, ne cíle.}

\ulohahead{společná matematika \textnormal{\footnotesize[úroveň 3]}}{Posloupnosti, limity a řady}
\begin{zadani}
\begin{enumerate}[label=(\alph*)]
  \item \bod{1} Definujte limitu posloupnosti a konvergenci. Uveďte aritmetiku limit a větu o dvou policajtech.
  \item \bod{2} Spočtěte $\displaystyle\lim_{n\to\infty}\frac{2n^2+1}{n^2+n}$ a $\displaystyle\lim_{n\to\infty}\sqrt[n]{n}$.
  \item \bod{3} Definujte součet řady a integrálním kritériem rozhodněte o konvergenci. Proč harmonická řada diverguje a geometrická (pro $|q|<1$) konverguje?
\end{enumerate}
\end{zadani}

\begin{reseni}
\textbf{(a)} $\lim_{n\to\infty}a_n=L$, pokud $\forall\varepsilon>0\,\exists n_0\,\forall n\ge n_0:|a_n-L|<\varepsilon$ (posloupnost je pak \emph{konvergentní}). \emph{Aritmetika limit}: limita součtu/součinu/podílu je součet/součin/podíl limit (u podílu při nenulové limitě jmenovatele). \emph{Věta o dvou policajtech}: je-li $a_n\le b_n\le c_n$ a $\lim a_n=\lim c_n=L$, pak $\lim b_n=L$.

\textbf{(b)} $\lim\frac{2n^2+1}{n^2+n}=\lim\frac{2+1/n^2}{1+1/n}=\frac{2}{1}=2$ (vydělením $n^2$). $\lim\sqrt[n]{n}=1$ (např.\ $\sqrt[n]{n}=e^{\frac{\ln n}{n}}$ a $\frac{\ln n}{n}\to0$).

\textbf{(c)} \emph{Součet řady} $\sum a_n$ je limita částečných součtu $s_N=\sum_{n=1}^N a_n$ (pokud existuje). \emph{Integrální kritérium} (pro kladnou klesající $f$): $\sum_{n\ge1} f(n)$ konverguje právě tehdy, když konverguje $\int_1^\infty f(x)\,dx$. \emph{Harmonická} $\sum\frac1n$ diverguje ($\int_1^\infty\frac{dx}{x}=\infty$; částečné součty rostou jako $\ln N$); \emph{geometrická} $\sum_{n=0}^\infty q^n$ pro $|q|<1$ konverguje k $\frac{1}{1-q}$ (částečný součet $\sum_{n=0}^N q^n=\frac{1-q^{N+1}}{1-q}$; součet od $n=1$ je $\frac{q}{1-q}$). \emph{(Nad rámec požadavku existují i srovnávací, podílové a odmocninové kritérium.)}
\end{reseni}

\begin{jadro}
Limita posloupnosti: $\forall\varepsilon\,\exists n_0:\,|a_n-L|<\varepsilon$. Racionální podíl polynomu: vyděl nejvyšší mocninou. Řada: integrální kritérium ($\sum f(n)$ konverguje $\iff \int f$ konverguje); harmonická diverguje, geometrická ($|q|<1$) $\to\frac1{1-q}$.
\end{jadro}

\kalibrace{Posloupnosti a řady jsou v požadavku analýzy a opakovaně se zkouší (geometrická řada byla samostatná otázka). Definice limity + jeden výpočet = 2 body.}
\past{Harmonická řada diverguje, ač $a_n\to0$ - samotné $a_n\to0$ je nutná, ne postačující podmínka konvergence.}

\ulohahead{společná matematika \textnormal{\footnotesize[úroveň 3]}}{Částečná uspořádání: výška, šířka a věta o dlouhém a širokém}
\begin{zadani}
\begin{enumerate}[label=(\alph*)]
  \item \bod{1} Definujte částečné uspořádání, řetězec a antiřetězec a minimální/maximální prvek.
  \item \bod{2} Definujte výšku a šířku částečně uspořádané množiny.
  \item \bod{3} Zformulujte větu o dlouhém a širokém (Dilworth/Mirsky) a ilustrujte na příkladu (dělitelnost na $\{1,2,3,4,6,12\}$).
\end{enumerate}
\end{zadani}

\begin{reseni}
\textbf{(a)} \emph{Částečné uspořádání} $\le$ je reflexivní, antisymetrická a tranzitivní relace. \emph{Řetězec} = podmnožina po dvou \emph{porovnatelných} prvku; \emph{antiřetězec} = podmnožina po dvou \emph{neporovnatelných} prvku. \emph{Minimální} prvek nemá nic ostře pod sebou (\emph{nejmenší} je pod vším); \emph{maximální} nemá nic nad sebou.

\textbf{(b)} \emph{Výška} = velikost nejdelšího řetězce. \emph{Šířka} = velikost největšího antiřetězce.

\textbf{(c)} \emph{Dilworthova věta}: minimální počet řetězcu, na něž lze rozložit uspořádanou množinu, se rovná \emph{šířce} (velikosti největšího antiřetězce). \emph{Duálně (Mirsky)}: minimální počet antiřetězcu pokrývajících množinu se rovná \emph{výšce}. Příklad (dělitelnost na $\{1,2,3,4,6,12\}$): nejdelší řetězec $1\mid2\mid4\mid12$ (resp.\ $1\mid2\mid6\mid12$) má 4 prvky, takže \emph{výška} $=4$. Největší antiřetězec je $\{4,6\}$ (vzájemně nedělitelné), případně $\{2,3\}$ - \emph{šířka} $=2$, takže množinu lze pokrýt 2 řetězci. Z Mirskyho plyne \emph{věta o dlouhém a širokém}: pro každou konečnou uspořádanou množinu je výška$\,\cdot\,$šířka $\ge|X|$ (pokrytí výškou antiřetězcu, každý nejvýše šířka); zde $4\cdot2=8\ge6$.
\end{reseni}

\begin{jadro}
Řetězec = porovnatelné prvky, antiřetězec = neporovnatelné. Výška = nejdelší řetězec, šířka = největší antiřetězec. Dilworth: min.\ počet řetězcu na pokrytí $=$ šířka; Mirsky: min.\ počet antiřetězcu $=$ výška.
\end{jadro}

\kalibrace{Výška/šířka a věta o dlouhém a širokém jsou explicitně v požadavku diskrétní matematiky, ale vzácné; úroveň 3 je pokrývá pro úplnost.}
\past{Dilworth páruje \emph{řetězce} se \emph{šířkou} (antiřetězcem), Mirsky \emph{antiřetězce} s \emph{výškou} - nezaměň směr. Minimální prvek nemusí být nejmenší.}

\ulohahead{společná matematika \textnormal{\footnotesize[úroveň 3]}}{Konkrétní rozdělení a nerovnosti}
\begin{zadani}
\begin{enumerate}[label=(\alph*)]
  \item \bod{1} Definujte binomické, geometrické, Poissonovo, normální a exponenciální rozdělení a uveďte jejich střední hodnotu.
  \item \bod{2} Zformulujte Markovovu nerovnost a použijte ji. \emph{(Volitelně: Čebyševova nerovnost.)}
  \item \bod{3} Zformulujte zákon velkých čísel a vysvětlete vztah binomické $\to$ Poissonovo/normální.
\end{enumerate}
\end{zadani}

\begin{reseni}
\textbf{(a)} \emph{Binomické} $\mathrm{Bi}(n,p)$: počet úspěchu v $n$ nezávislých pokusech, $P(X=k)=\binom{n}{k}p^k(1-p)^{n-k}$, $\mathbb{E}X=np$. \emph{Geometrické}: počet pokusu do prvního úspěchu, $P(X=k)=(1-p)^{k-1}p$, $\mathbb{E}X=1/p$. \emph{Poissonovo} $\mathrm{Po}(\lambda)$: $P(X=k)=e^{-\lambda}\lambda^k/k!$, $\mathbb{E}X=\lambda$. \emph{Normální} $\mathcal{N}(\mu,\sigma^2)$: hustota $\frac{1}{\sigma\sqrt{2\pi}}e^{-(x-\mu)^2/(2\sigma^2)}$, $\mathbb{E}X=\mu$. \emph{Exponenciální} $\mathrm{Exp}(\lambda)$: hustota $\lambda e^{-\lambda x}$ pro $x\ge0$, $\mathbb{E}X=1/\lambda$.

\textbf{(b)} \emph{Markovova}: pro $X\ge0$ a $a>0$ je $P(X\ge a)\le\frac{\mathbb{E}X}{a}$. Příklad: má-li $X\ge0$ střední hodnotu $2$, pak $P(X\ge10)\le 2/10=0{,}2$. \emph{(Nad rámec požadavku, příbuzná Čebyševova:} $P(|X-\mathbb{E}X|\ge k)\le\frac{\operatorname{var}X}{k^2}$.\emph{)}

\textbf{(c)} \emph{Zákon velkých čísel} (slabý): pro nezávislé stejně rozdělené se střední hodnotou $\mu$ platí $\bar X_n\xrightarrow{P}\mu$, tj.\ $\forall\varepsilon>0:P(|\bar X_n-\mu|>\varepsilon)\to0$ (silný: konvergence skoro jistě). Binomické $\to$ \emph{Poissonovo} pro velké $n$ a malé $p$ s $np=\lambda$ (zákon vzácných jevu); binomické $\to$ \emph{normální} pro velké $n$ (de Moivre-Laplace, speciální případ CLT) s $\mu=np$, $\sigma^2=np(1-p)$.
\end{reseni}

\begin{jadro}
$\mathrm{Bi}(n,p)$: $\mathbb{E}=np$; geom.: $1/p$; $\mathrm{Po}(\lambda)$: $\lambda$; $\mathcal{N}(\mu,\sigma^2)$: $\mu$; $\mathrm{Exp}(\lambda)$: $1/\lambda$. Markov $P(X\ge a)\le\mathbb{E}X/a$ (Čebyšev $\sigma^2/k^2$ navíc). ZVČ: prumer $\to$ střední hodnota.
\end{jadro}

\kalibrace{Práce s konkrétními rozděleními a nerovnostmi je v požadavku pravděpodobnosti; úroveň 3 je shrnuje pro úplnost a podporuje statistické testy ze specializace.}
\past{Geometrické: $\mathbb{E}=1/p$, ne $p$. Markov potřebuje $X\ge0$. Binomické $\to$ Poisson při \emph{malém} $p$, $\to$ normální při velkém $n$.}

\ulohahead{společná matematika \textnormal{\footnotesize[úroveň 3]}}{Bodové a intervalové odhady}
\begin{zadani}
\begin{enumerate}[label=(\alph*)]
  \item \bod{1} Definujte bodový odhad, nestrannost a konzistenci odhadu.
  \item \bod{2} Popište metodu maximální věrohodnosti (MLE) a odvoďte MLE parametru $p$ z $n$ Bernoulliho pokusu s $k$ úspěchy.
  \item \bod{3} Popište intervalový odhad střední hodnoty založený na normální aproximaci.
\end{enumerate}
\end{zadani}

\begin{reseni}
\textbf{(a)} \emph{Bodový odhad} $\hat\theta=\hat\theta(X_1,\dots,X_n)$ je statistika odhadující neznámý parametr $\theta$. Je \emph{nestranný}, je-li $\mathbb{E}[\hat\theta]=\theta$ (žádné systematické vychýlení). Je \emph{konzistentní}, pokud $\hat\theta\xrightarrow{P}\theta$ s rostoucím $n$ (např.\ prumer $\bar X$ je nestranný a konzistentní odhad $\mathbb{E}X$).

\textbf{(b)} \emph{MLE}: zvol $\theta$ maximalizující věrohodnost $L(\theta)=\prod_i P(x_i\mid\theta)$ (typicky maximalizujeme $\ln L$). Pro $n$ Bernoulliho pokusu s $k$ úspěchy: $L(p)=p^k(1-p)^{n-k}$, $\ln L=k\ln p+(n-k)\ln(1-p)$; derivace $\frac{k}{p}-\frac{n-k}{1-p}=0$ dává $\hat p=\frac{k}{n}$ (relativní četnost); hraničně $k{=}0\Rightarrow\hat p{=}0$ a $k{=}n\Rightarrow\hat p{=}1$.

\textbf{(c)} \emph{Intervalový odhad} střední hodnoty: pro velké $n$ je $\bar X\approx\mathcal{N}(\mu,\sigma^2/n)$ (CLT), takže $100(1-\alpha)\,\%$ konfidenční interval je
\[
\bar X\pm z_{1-\alpha/2}\,\frac{\sigma}{\sqrt n}
\]
($z_{1-\alpha/2}$ je kvantil normálního rozdělení, např.\ $1{,}96$ pro 95\,\%); neznámé $\sigma$ se nahradí výběrovou směrodatnou odchylkou $s$ (pro malé $n$ se použije $t$-kvantil).
\end{reseni}

\begin{jadro}
Nestranný: $\mathbb{E}[\hat\theta]=\theta$. MLE: maximalizuj $\ln L(\theta)=\sum\ln P(x_i\mid\theta)$; pro Bernoulli $\hat p=k/n$. CI střední hodnoty: $\bar X\pm z_{1-\alpha/2}\,\sigma/\sqrt n$.
\end{jadro}

\kalibrace{Bodové i intervalové odhady jsou explicitně v požadavku statistiky a codex je označil za vysokovýnosové pro úplnost; podporují i ML (maximální věrohodnost).}
\past{Nestrannost je o $\mathbb{E}[\hat\theta]=\theta$, ne o rozptylu. Šířka konfidenčního intervalu klesá jako $1/\sqrt n$ (čtyřnásobek dat = poloviční interval).}


\section{Společná informatika}
\ulohahead{společná informatika}{Regulární jazyky: automaty, výrazy, uzávěry, neregularita}
\begin{zadani}
\begin{enumerate}[label=(\alph*)]
  \item \bod{1} Definujte deterministický (DFA) a nedeterministický (NFA) konečný automat, regulární gramatiku a regulární výraz. Zformulujte Kleeneho větu.
  \item \bod{2} Sestrojte DFA pro jazyk slov nad $\{a,b\}$ obsahujících podslovo $ab$. Uveďte uzávěrové vlastnosti regulárních jazyku (sjednocení, průnik, doplněk).
  \item \bod{3} Dokažte, že jazyk $L=\{a^n b^n\mid n\ge0\}$ není regulární.
\end{enumerate}
\end{zadani}

\begin{reseni}
\textbf{(a)} \emph{DFA} $=(Q,\Sigma,\delta,q_0,F)$ s $\delta:Q\times\Sigma\to Q$; přijímá $w$, skončí-li ve stavu z $F$. \emph{NFA} má $\delta:Q\times\Sigma\to 2^Q$ (a případně $\varepsilon$-přechody); přijímá, existuje-li přijímající běh. \emph{Regulární (pravá lineární) gramatika} je čtveřice $G=(N,\Sigma,P,S)$: neterminály $N$, terminály $\Sigma$, počáteční neterminál $S\in N$ a pravidla $P$ tvaru $A\to aB$ nebo $A\to a$ (případně $S\to\varepsilon$). \emph{Regulární výraz} se tvoří z $\emptyset,\varepsilon,$ symbolu operacemi $\cup,\cdot,{}^*$. \emph{Kleeneho věta:} třídy jazyku popsatelných DFA, NFA, regulární gramatikou a regulárním výrazem jsou totožné (= regulární jazyky).

\textbf{(b)} DFA se 3 stavy, $F=\{q_2\}$:
\[
\begin{array}{c|cc}
 & a & b\\\hline
q_0 & q_1 & q_0\\
q_1 & q_1 & q_2\\
q_2 & q_2 & q_2
\end{array}
\]
($q_0$ = zatím nic, $q_1$ = naposledy $a$, $q_2$ = už viděno $ab$). \emph{Uzávěry:} regulární jazyky jsou uzavřené na sjednocení a průnik (součinový automat na $Q_1\times Q_2$), na doplněk (u DFA prohodíme $F\leftrightarrow Q\setminus F$), dále na zřetězení a iteraci.

\textbf{(c)} Sporem přes \emph{pumping lemma}: nechť $L$ regulární s konstantou $p$. Vezmi $w=a^p b^p\in L$, $|w|\ge p$. Lemma dá rozklad $w=xyz$, $|xy|\le p$, $|y|\ge1$. Protože $|xy|\le p$, leží $y$ celé v úvodních $a$, tedy $y=a^k$, $k\ge1$. Pak $xy^2z=a^{p+k}b^p\notin L$ (více $a$ než $b$), což je spor. Tedy $L$ není regulární.
\end{reseni}

\kalibrace{Automaty a jazyky jsou nejčastější okruh společné informatiky (~78\,\% zkoušek, skoro vždy jako Q1). Konstrukce DFA + definice = spolehlivé 2 body; tato úloha je proto v jádru úrovně 1.}
\past{Pumping lemma dokazuje \emph{ne}regularitu (sporem); neslouží k důkazu regularity. Vždy zdůvodni, proč $y$ padne do bloku $a$ (díky $|xy|\le p$).}

\ulohahead{společná informatika}{Bezkontextové jazyky, Turinguv stroj a nerozhodnutelnost}
\begin{zadani}
\begin{enumerate}[label=(\alph*)]
  \item \bod{1} Definujte bezkontextovou gramatiku (CFG), jazyk jí generovaný a zásobníkový automat (PDA). Co je Chomského hierarchie (typy 0-3)?
  \item \bod{2} Sestrojte CFG pro $\{a^n b^n\mid n\ge0\}$ a popište princip převodu CFG na PDA.
  \item \bod{3} Definujte Turinguv stroj a rekurzivně spočetné jazyky. Vysvětlete existenci algoritmicky nerozhodnutelného problému (problém zastavení).
\end{enumerate}
\end{zadani}

\begin{reseni}
\textbf{(a)} \emph{CFG} $G=(N,\Sigma,P,S)$ s pravidly $A\to\alpha$, $A\in N$, $\alpha\in(N\cup\Sigma)^*$; \emph{generovaný jazyk} = slova nad $\Sigma$ odvoditelná z $S$. \emph{PDA} = konečný automat se zásobníkem; třída jazyku přijímaných PDA = bezkontextové jazyky (CFL). \emph{Chomského hierarchie:} typ 3 regulární $\subset$ typ 2 bezkontextové $\subset$ typ 1 kontextové $\subset$ typ 0 rekurzivně spočetné, s odpovídajícími stroji (KA, PDA, lineárně omezený automat, Turinguv stroj).

\textbf{(b)} CFG: $S\to aSb \mid \varepsilon$. \emph{Převod CFG$\to$PDA} (přijímání prázdným zásobníkem): na zásobník dej $S$; v každém kroku buď nahraď nejvyšší neterminál pravou stranou pravidla ($\varepsilon$-přechod), nebo paruj nejvyšší terminál se vstupním symbolem (a oba odeber). Slovo je přijato, vyprázdní-li se zásobník při dočtení vstupu. PDA tak simuluje levou derivaci.

\textbf{(c)} \emph{Turinguv stroj} = konečný řídicí automat s nekonečnou páskou, hlavou (čte/zapisuje/posouvá), přechodovou funkcí; přijímá vstup, dostane-li se do přijímajícího stavu. Jazyk je \emph{rekurzivně spočetný} (typ 0), přijímá-li ho nějaký TS (může cyklit na slovech mimo jazyk); \emph{rekurzivní} (rozhodnutelný), pokud TS vždy zastaví. \emph{Problém zastavení} (zastaví TS $M$ na vstupu $w$?) je nerozhodnutelný: kdyby existoval rozhodovač $H$, sestrojíme $D(x)$, který se zacyklí, právě když $H$ řekne ,,$x(x)$ zastaví''; pak $D(D)$ vede ke sporu (diagonalizace).
\end{reseni}

\kalibrace{Bezkontextové jazyky a Chomského hierarchie jsou jádro Q1; Turinguv stroj + nerozhodnutelnost jsou vzácnější, ale legální tah z téhož okruhu - tady je držíme na definiční úrovni jako pojistku proti nule.}
\past{,,Bezkontextové'' $=$ PDA, ne konečný automat. Nerozhodnutelnost se dokazuje diagonalizací/redukcí, ne ,,je to těžké''.}

\ulohahead{společná informatika}{Složitost, třídy P a NP, NP-úplnost}
\begin{zadani}
\begin{enumerate}[label=(\alph*)]
  \item \bod{1} Definujte časovou složitost a asymptotickou notaci ($O,\Omega,\Theta$). Definujte třídy P a NP.
  \item \bod{2} Definujte polynomiální převoditelnost, NP-těžkost a NP-úplnost. Uveďte aspoň tři NP-úplné problémy.
  \item \bod{3} Jak se dokazuje, že je problém NP-úplný? Co tvrdí Cookova-Levinova věta a co znamená otevřená otázka $\mathrm{P}\overset{?}{=}\mathrm{NP}$?
\end{enumerate}
\end{zadani}

\begin{reseni}
\textbf{(a)} \emph{Časová složitost} = počet elementárních kroku jako funkce velikosti vstupu (nejhorší případ). $f=O(g)$: $\exists c,n_0:\,f(n)\le c\,g(n)$ pro $n\ge n_0$; $\Omega$ je dolní mez, $\Theta=O\cap\Omega$. \emph{P} = problémy řešitelné deterministicky v polynomiálním čase. \emph{NP} = problémy, jejichž ,,ano'' instance mají v polynomiálním čase ověřitelný certifikát (ekvivalentně řešitelné nedeterministicky v polynomiálním čase).

\textbf{(b)} \emph{Polynomiální převod} $A\le_p B$: funkce vyčíslitelná v polynomiálním čase, která instance $A$ převádí na instance $B$ se zachováním odpovědi. $B$ je \emph{NP-těžký}, pokud $A\le_p B$ pro každé $A\in\mathrm{NP}$; \emph{NP-úplný}, je-li navíc $B\in\mathrm{NP}$. Příklady NP-úplných: SAT, 3-SAT, klika, vrcholové pokrytí, barevnost grafu, Hamiltonovská kružnice, batoh (rozhodovací).

\textbf{(c)} NP-úplnost se dokazuje ve dvou krocích: (1) ukázat $B\in\mathrm{NP}$ (certifikát + ověření); (2) zvolit známý NP-úplný problém $A$ a sestrojit polynomiální redukci $A\le_p B$. \emph{Cookova-Levinova věta:} SAT je NP-úplný (první NP-úplný problém), což dává startovní bod pro všechny další redukce. $\mathrm{P}=\mathrm{NP}$? je otevřená otázka, zda lze vše ověřitelné v polynomiálním čase i v polynomiálním čase \emph{vyřešit}; obecně se věří, že $\mathrm{P}\neq\mathrm{NP}$.
\end{reseni}

\kalibrace{Algoritmy a složitost (P/NP) jsou opakovaný okruh; čistě definiční úloha s vysokou návratností bodu. Redukce z (b)/(c) přidá třetí bod.}
\past{NP \emph{není} ,,ne-polynomiální''. NP-těžkost nevyžaduje $B\in$ NP; NP-úplnost ano. Směr redukce: ze \emph{známého} těžkého na \emph{nový} problém.}

\ulohahead{společná informatika}{Rozděl a panuj, Master theorem, třídění}
\begin{zadani}
\begin{enumerate}[label=(\alph*)]
  \item \bod{1} Popište princip metody rozděl a panuj. Zformulujte Master theorem (kuchařkovou větu) a dolní odhad složitosti porovnávacího třídění.
  \item \bod{2} Popište Mergesort (pseudokód) a odvoďte jeho složitost z rekurentní rovnice. Uveďte princip a nejhorší/průměrný případ Quicksortu.
  \item \bod{3} Pomocí Master theorem určete asymptotiku pro $T(n)=2T(n/2)+n$ a $T(n)=3T(n/2)+n$. Co z toho plyne pro Karatsubovo násobení?
\end{enumerate}
\end{zadani}

\begin{reseni}
\textbf{(a)} \emph{Rozděl a panuj}: rozděl úlohu na podproblémy, vyřeš rekurzivně, slož výsledky. \emph{Master theorem:} pro $T(n)=aT(n/b)+f(n)$ ($a\ge1,b>1$) porovnej $f(n)$ s $n^{\log_b a}$:
\begin{itemize}\setlength\itemsep{0pt}
  \item $f=O(n^{\log_b a-\varepsilon})\Rightarrow T=\Theta(n^{\log_b a})$;
  \item $f=\Theta(n^{\log_b a})\Rightarrow T=\Theta(n^{\log_b a}\log n)$;
  \item $f=\Omega(n^{\log_b a+\varepsilon})$ a regularita $\Rightarrow T=\Theta(f(n))$.
\end{itemize}
\emph{Dolní odhad} porovnávacího třídění: $\Omega(n\log n)$ (rozhodovací strom má $n!$ listu, hloubka $\ge\log_2 n!=\Theta(n\log n)$).

\textbf{(b)} Mergesort: rozděl pole na poloviny, setřiď rekurzivně, slij ($\textsc{Merge}$) v $O(n)$. Rekurence $T(n)=2T(n/2)+O(n)$, dle Master (případ 2) $T(n)=\Theta(n\log n)$; stabilní, $O(n)$ paměť navíc. Quicksort: zvol pivota, rozděl na menší/větší, rekurzivně; průměrně $O(n\log n)$, nejhůře $O(n^2)$ (špatný pivot), in-place.

\textbf{(c)} $T(n)=2T(n/2)+n$: $a=b=2$, $n^{\log_b a}=n$, $f=\Theta(n)$, případ 2 $\Rightarrow \Theta(n\log n)$. $T(n)=3T(n/2)+n$: $\log_2 3\approx1{,}585$, $n^{\log_2 3}$ roste rychleji než $f=n$, případ 1 $\Rightarrow \Theta(n^{\log_2 3})$. \emph{Karatsuba} násobí $n$-ciferná čísla rekurzí $T(n)=3T(n/2)+O(n)$, tedy $\Theta(n^{\log_2 3})\approx n^{1{,}585}$ - lepší než školní $\Theta(n^2)$.
\end{reseni}

\kalibrace{Třídění a rozděl-a-panuj jsou opakovaný algoritmický okruh; Master theorem + Mergesort je vděčný typový výpočet na 2-3 body.}
\past{Master theorem porovnává $f(n)$ s $n^{\log_b a}$, ne s $n$. Quicksort má nejhorší případ $O(n^2)$ - nezaměňuj s průměrným.}

\ulohahead{společná informatika}{Grafové algoritmy: prohledávání, nejkratší cesty, kostry, toky}
\begin{zadani}
\begin{enumerate}[label=(\alph*)]
  \item \bod{1} Popište prohledávání do šířky (BFS) a do hloubky (DFS), jejich složitost a reprezentaci grafu.
  \item \bod{2} Popište topologické třídění DAG a algoritmy pro nejkratší cesty: Dijkstra a Bellman-Ford (kdy který, složitost).
  \item \bod{3} Popište minimální kostru a algoritmy Jarník (Prim) a Borůvka (řezové lemma). Uveďte princip hledání maximálního toku (Ford-Fulkerson).
\end{enumerate}
\end{zadani}

\begin{reseni}
\textbf{(a)} \emph{BFS}: z počátku prohledává po vrstvách pomocí fronty; dává nejkratší cesty v neohodnoceném grafu. \emph{DFS}: jde do hloubky pomocí zásobníku/rekurze; klasifikuje hrany (stromové, zpětné, dopředné, příčné), detekuje cykly. Obě $O(V+E)$ při reprezentaci \emph{seznamy sousedu}.

\textbf{(b)} \emph{Topologické třídění} (jen DAG): lineární uspořádání vrcholu tak, že každá hrana vede ,,vpřed''; spočítá se opakovaným odebíráním vrcholu se vstupním stupněm 0 nebo přes časy dokončení DFS, $O(V+E)$. \emph{Dijkstra}: nejkratší cesty z jednoho zdroje pro \emph{nezáporné} váhy, hladově s prioritní frontou, $O((V+E)\log V)$. \emph{Bellman-Ford}: zvládá \emph{záporné} hrany, $V-1$ relaxací všech hran, $O(VE)$, navíc detekuje záporný cyklus.

\textbf{(c)} \emph{Minimální kostra} (MST) = kostra s minimálním součtem vah. \emph{Jarník/Prim}: roste jeden strom, vždy přidá nejlevnější hranu ven, $O((V+E)\log V)$. \emph{Borůvka}: každá komponenta paralelně přidá svou nejlevnější hranu, $O(E\log V)$. Korektnost plyne z \emph{řezového lemmatu}: nejlevnější hrana přes libovolný řez patří do nějaké MST. \emph{Ford-Fulkerson}: opakovaně hledej zlepšující cestu ze $s$ do $t$ v reziduální síti a posílej po ní tok, dokud existuje. Pro celočíselné kapacity vždy skončí; když už zlepšující cesta neexistuje, je tok maximální a jeho velikost se rovná kapacitě minimálního řezu.
\end{reseni}

\kalibrace{Grafové algoritmy jsou opakovaný informatický okruh (BFS/DFS, topo, MST, toky). Definice + jeden algoritmus s složitostí = 2 body; tady jde o breadth, ne hloubku důkazu.}
\past{Dijkstra selže na záporných hranách - tam Bellman-Ford. Řezové lemma je klíč ke korektnosti MST; znej ho aspoň slovně.}

\ulohahead{společná informatika}{Datové struktury: vyhledávací stromy a haldy}
\begin{zadani}
\begin{enumerate}[label=(\alph*)]
  \item \bod{1} Definujte binární vyhledávací strom (BVS) a jeho operace (find/insert/delete). Co je AVL strom a jakou má výšku?
  \item \bod{2} Popište binární (min-)haldu v poli a operace \textsc{Insert} a \textsc{ExtractMin}. Jak funguje Heapsort a jaká je jeho složitost?
  \item \bod{3} Popište hledání následníka (successor) v BVS. Porovnejte složitost operací u nevyváženého a vyváženého stromu. Proč jde halda postavit v $O(n)$?
\end{enumerate}
\end{zadani}

\begin{reseni}
\textbf{(a)} \emph{BVS}: binární strom, kde pro každý uzel platí klíče vlevo $<$ klíč uzlu $<$ klíče vpravo. \textsc{Find}/\textsc{Insert} jdou shora dle porovnání; \textsc{Delete} listu/uzlu s jedním synem je přímé, uzel se dvěma syny se nahradí svým následníkem. Vše $O(h)$ (výška). \emph{AVL} je BVS udržující v každém uzlu rozdíl výšek podstromu $\le1$ (rotacemi), takže $h=O(\log n)$.

\textbf{(b)} \emph{Min-halda} = úplný binární strom v poli (uzel $i$ má syny $2i{+}1,2i{+}2$), kde rodič $\le$ synové. \textsc{Insert}: přidej na konec a ,,probublej nahoru'' ($O(\log n)$). \textsc{ExtractMin}: vrať kořen, dej poslední prvek na vrchol a ,,probublej dolu'' ($O(\log n)$). \emph{Heapsort}: postav haldu a $n$-krát vyber minimum/maximum; $O(n\log n)$, in-place.

\textbf{(c)} \emph{Successor} (nejbližší větší): má-li uzel pravý podstrom, je to jeho nejlevnější uzel; jinak jdi nahoru, dokud nejsi levým synem rodiče (ten je následník). U \emph{nevyváženého} BVS je $h$ až $O(n)$ (degenerace na seznam), takže operace $O(n)$; u \emph{vyváženého} (AVL) $O(\log n)$. Haldu lze postavit \emph{zdola nahoru} (\textsc{Heapify} od posledního vnitřního uzlu): součet prací $\sum_h \tfrac{n}{2^{h+1}}\cdot O(h)=O(n)$.
\end{reseni}

\kalibrace{Stromy a haldy jsou opakovaný okruh datových struktur (BVS, AVL, Heapsort se objevily v zadáních). Definice + operace s složitostí = 2 body.}
\past{Mazání uzlu se dvěma syny v BVS jde přes následníka, ne ledabyle. Build-heap je $O(n)$, nikoli $O(n\log n)$ - klasická chytací otázka.}

\ulohahead{společná informatika}{Reprezentace dat a čtení binárního souboru (C\#)}
\begin{zadani}
Binární soubor má tento formát: 4 bajty ASCII ,,\texttt{REC1}'' (magic), poté \texttt{uint32} \texttt{count} (little-endian), poté \texttt{count} záznamů, kde každý záznam je \texttt{int32 id} (little-endian) následovaný \texttt{float32 value} (IEEE\,754, little-endian).
\begin{enumerate}[label=(\alph*)]
  \item \bod{1} Vysvětlete pojmy \emph{little-endian} vs.\ \emph{big-endian} a \emph{zarovnání} (alignment). Jak je v paměti uloženo \texttt{int32} s hodnotou \texttt{0x01020304} v little-endian?
  \item \bod{2} Napište v C\# metodu, která soubor načte do pole struktur \texttt{(int Id, float Value)}.
  \item \bod{3} Co se stane na big-endian stroji a jak to ošetřit? Jak ověřit magic a obránit se proti poškozenému/zkrácenému souboru?
\end{enumerate}
\end{zadani}

\begin{reseni}
\textbf{(a)} \emph{Endianita} je pořadí bajtu vícebajtové hodnoty v paměti. \emph{Little-endian} ukládá nejméně významný bajt na nejnižší adresu. Hodnota \texttt{0x01020304} se tedy v paměti jeví jako bajty \texttt{04 03 02 01} (od nízké adresy). Big-endian by uložil \texttt{01 02 03 04}. \emph{Zarovnání} znamená, že hodnota typu o velikosti $s$ má typicky začínat na adrese dělitelné $s$ (např.\ \texttt{int32} na násobku 4); překladač proto vkládá výplň (padding) mezi položky struktury. V binárním \emph{formátu souboru} se ale spoléháme na přesné pořadí bajtu a žádné implicitní zarovnání nepředpokládáme.

\textbf{(b)} \texttt{BinaryReader} v .NET čte v little-endian, což sedí na zadání:
\begin{lstlisting}
public readonly record struct Rec(int Id, float Value);

public static Rec[] Load(string path)
{
    using var fs = File.OpenRead(path);
    using var br = new BinaryReader(fs);

    // magic "REC1"
    var magic = br.ReadBytes(4);
    if (magic.Length != 4 || magic[0] != (byte)'R' || magic[1] != (byte)'E'
        || magic[2] != (byte)'C' || magic[3] != (byte)'1')
        throw new InvalidDataException("spatny magic");

    uint count = br.ReadUInt32();          // little-endian
    var recs = new Rec[count];
    for (uint i = 0; i < count; i++)
    {
        int id = br.ReadInt32();           // 4 bajty LE
        float val = br.ReadSingle();       // 4 bajty IEEE754 LE
        recs[i] = new Rec(id, val);
    }
    return recs;
}
\end{lstlisting}
Ruční složení int32 z bajtu (ukázka principu): \texttt{id = b0 | b1<<8 | b2<<16 | b3<<24}.

\textbf{(c)} Pozor na endianitu CPU: \texttt{BinaryReader} i explicitní složení \texttt{b0 | b1<<8 | b2<<16 | b3<<24} dávají \emph{vždy} little-endian výsledek nezávisle na procesoru (to je správně a přenositelné). Co přenositelné \emph{není}, je reinterpretace syrových bajtu v nativním pořadí CPU (např.\ \texttt{BitConverter.ToInt32} nebo přetypování ukazatele) - ta by na big-endian stroji vrátila prohozené bajty. Robustní je proto skládat bajty s pevným pořadím podle \emph{formátu} (LE): v .NET \texttt{BinaryPrimitives.ReadInt32LittleEndian(span)} a \texttt{BinaryPrimitives.ReadSingleLittleEndian(span)}, případně podle \texttt{BitConverter.IsLittleEndian} bajty otočit. \emph{Obrana proti poškození}: ověř magic; po načtení \texttt{count} zkontroluj, že zbývající délka souboru je $=8\cdot\texttt{count}$ bajtu (jinak vyhoď výjimku) místo slepého čtení do EOF; \texttt{ReadBytes}/\texttt{ReadInt32} při zkrácení vyhodí \texttt{EndOfStreamException}, kterou je nutno ošetřit.
\end{reseni}

\kalibrace{Implementační otázky ze společné informatiky (binární soubor, alokátor, semafor, OO návrh) jsou hlavní příčina propadnutí ,,naoko připraveného'' studenta - na nule z této otázky se podlaha $\ge 5/12$ hroutí. Cíl: NIKDY nedostat 0. Část (a) je čistá definice (1 bod), kostra metody v (b) přidá druhý bod i bez doladění (c).}
\past{\texttt{BinaryReader} čte vždy little-endian bez ohledu na CPU - na zkoušce to klidně řekni, ale dodej, že pro přenositelnost a big-endian CPU je správně \texttt{BinaryPrimitives}. Nezapomeň na ošetření konce souboru.}

\ulohahead{společná informatika}{Architektura a OS: režimy, procesy, virtuální paměť}
\begin{zadani}
\begin{enumerate}[label=(\alph*)]
  \item \bod{1} Vysvětlete adresu a adresový prostor, privilegovaný vs neprivilegovaný režim procesoru a roli jádra OS. Co je proces, vlákno, kontext a přepínání kontextu? Uveďte stavy vlákna.
  \item \bod{2} Vysvětlete rozdíl virtuální vs fyzická adresa a komponenty stránkování (číslo stránky, číslo rámce, offset). Pro stránku 4\,KiB přeložte virtuální adresu \texttt{0x00003ABC}, je-li stránka 3 namapována na rámec 7.
  \item \bod{3} Jakou roli hrají virtuální adresové prostory v ochraně paměti? Jak procesor realizuje konstrukce vyššího jazyka (přiřazení, podmínka, cyklus) instrukcemi?
\end{enumerate}
\end{zadani}

\begin{reseni}
\textbf{(a)} \emph{Adresa} identifikuje místo v paměti; \emph{adresový prostor} = množina platných adres (každý proces má vlastní). V \emph{privilegovaném} (jádrovém) režimu smí procesor vše (I/O, správa paměti), v \emph{neprivilegovaném} (uživatelském) jen omezeně; přechod přes přerušení/syscall. \emph{Jádro OS} spravuje prostředky a izolaci. \emph{Proces} = běžící program s vlastním adresovým prostorem; \emph{vlákno} = tok provádění uvnitř procesu (sdílí paměť). \emph{Kontext} = registry + čítač instrukcí + stav; \emph{přepnutí kontextu} je uloží/obnoví. Stavy vlákna: \emph{běží} (running), \emph{připraveno} (ready), \emph{čeká/blokováno} (waiting); plánovač přepíná preemptivně nebo kooperativně.

\textbf{(b)} \emph{Virtuální} adresu používá program; \emph{fyzická} ukazuje do RAM; překlad zajišťuje MMU přes stránkovací tabulku. Adresa $=$ (číslo stránky $\Vert$ offset). Stránka 4\,KiB $\Rightarrow$ offset 12 bitu. \texttt{0x00003ABC}: offset $=$ \texttt{0xABC}, číslo stránky $=$ \texttt{0x00003ABC}$\gg12=$ \texttt{0x3} $=3$. Tabulka: stránka $3\to$ rámec $7$. Fyzická adresa $=(7\ll12)\,\vert\,\texttt{0xABC}=\texttt{0x7ABC}$.

\textbf{(c)} Každý proces má vlastní mapování, takže nevidí cizí paměť (\emph{ochrana/izolace}); přístup mimo mapování vyvolá výpadek (page fault). Konstrukce vyššího jazyka: \emph{přiřazení} = load/store mezi registrem a pamětí; \emph{podmínka} = porovnání (\texttt{cmp}) + podmíněný skok; \emph{cyklus} = podmíněný skok zpět; \emph{volání funkce} = uložení návratové adresy a rámce na zásobník (\texttt{call}/\texttt{ret}).
\end{reseni}

\kalibrace{Architektura a OS jsou druhý nejčastější informatický okruh (~60\,\%); virtuální paměť a překlad adres se v zadáních opakují jako konkrétní výpočet. Tabulka adres je vděčné 2-3 body.}
\past{Offset má tolik bitu, kolik je $\log_2$(velikost stránky); nepleť číslo stránky (virtuální) s číslem rámce (fyzické). Stavy vlákna nejsou jen ,,běží/neběží''.}

\ulohahead{společná informatika}{Synchronizace: kritická sekce, semafory, producent-konzument (C\#)}
\begin{zadani}
\begin{enumerate}[label=(\alph*)]
  \item \bod{1} Definujte časově závislou chybu (race condition), kritickou sekci a vzájemné vyloučení. Co je zámek a jaký je rozdíl mezi aktivním a pasivním čekáním?
  \item \bod{2} Co je semafor (operace \textsc{wait}/P a \textsc{signal}/V)? Napište v C\# řešení producent-konzument s omezeným bufferem pomocí semaforu.
  \item \bod{3} Uveďte čtyři podmínky uváznutí (deadlock). Porovnejte spinlock a mutex a vysvětlete atomickou operaci compare-and-swap (CAS).
\end{enumerate}
\end{zadani}

\begin{reseni}
\textbf{(a)} \emph{Race condition}: výsledek závisí na nedeterministickém prokládání vláken nad sdíleným stavem. \emph{Kritická sekce}: úsek kódu, který smí běžet nejvýše jedno vlákno najednou. \emph{Vzájemné vyloučení} (mutual exclusion) tuto vlastnost zajišťuje. \emph{Zámek} (lock) chrání kritickou sekci; \emph{aktivní čekání} (spin) cyklí a pálí CPU, \emph{pasivní} vlákno uspí a plánovač ho probudí (lepší při delším čekání).

\textbf{(b)} \emph{Semafor} = čítač s atomickými operacemi: \textsc{wait}/P sníží čítač (a zablokuje při 0), \textsc{signal}/V zvýší (a probudí čekající). Omezený buffer kapacity $N$:
\begin{lstlisting}
// empty: kolik je volnych mist, full: kolik je obsazenych
var empty = new SemaphoreSlim(N, N);
var full  = new SemaphoreSlim(0, N);
var mutex = new object();
var buffer = new Queue<int>();

void Producer(int item) {
    empty.Wait();                 // pockej na volne misto (P)
    lock (mutex) { buffer.Enqueue(item); }
    full.Release();               // pribyl prvek (V)
}

int Consumer() {
    full.Wait();                  // pockej na prvek (P)
    int item;
    lock (mutex) { item = buffer.Dequeue(); }
    empty.Release();              // uvolnil se slot (V)
    return item;
}
\end{lstlisting}

\textbf{(c)} \emph{Deadlock} nastane při současném splnění: (1) vzájemné vyloučení, (2) drž a čekej, (3) bez odebrání (no preemption), (4) kruhové čekání. \emph{Spinlock} aktivně čeká (vhodný pro velmi krátké sekce na více jádrech), \emph{mutex} uspí vlákno (vhodný pro delší). \emph{CAS}\,(\texttt{CompareAndSwap}(addr, expected, new)): atomicky porovná a případně zapíše; základ neblokujících (lock-free) struktur.
\end{reseni}

\kalibrace{Synchronizace (semafory, kritická sekce, race condition) je nejčastější OS podtéma a často chce \emph{kód}. Definice (a) = jistý bod, kostra semaforu (b) = druhý; tím se vyhneš nule na implementační otázce.}
\past{Pořadí semaforu u producent-konzument je klíčové: čekej na zdroj (\textsc{wait}) PŘED zámkem, jinak hrozí deadlock. \texttt{empty} a \texttt{full} se nesmí prohodit.}

\ulohahead{společná informatika}{Objektové programování a návrh (C\#)}
\begin{zadani}
\begin{enumerate}[label=(\alph*)]
  \item \bod{1} Vysvětlete abstrakci, zapouzdření a polymorfismus a související konstrukty (třída, rozhraní, dědičnost, viditelnost). Rozdíl mezi statickým a dynamickým typováním a mezi virtuální a nevirtuální metodou.
  \item \bod{2} Navrhněte v C\# malou hierarchii geometrických tvaru s rozhraním a polymorfním výpočtem obsahu; ukažte ošetření chyby výjimkou.
  \item \bod{3} Jak je implementován dynamický polymorfismus (tabulka virtuálních metod)? K čemu jsou generické typy a jaký je rozdíl mezi hodnotovými a referenčními typy v C\#?
\end{enumerate}
\end{zadani}

\begin{reseni}
\textbf{(a)} \emph{Abstrakce} skrývá detaily za rozhraním; \emph{zapouzdření} drží data + operace pohromadě a omezuje přístup (viditelnost \texttt{private}/\texttt{public}/\texttt{protected}); \emph{polymorfismus} = jednotné volání nad různými typy. \emph{Třída} definuje typ, \emph{rozhraní} (interface) jen kontrakt, \emph{dědičnost} přebírá a rozšiřuje. \emph{Statické} typování kontroluje typy při překladu, \emph{dynamické} za běhu. \emph{Virtuální} metoda se vybírá podle skutečného typu objektu za běhu (override), \emph{nevirtuální} podle deklarovaného typu.

\textbf{(b)}
\begin{lstlisting}
public interface IShape { double Area(); }

public sealed class Circle : IShape {
    private readonly double r;
    public Circle(double r) {
        if (r < 0) throw new ArgumentException("zaporny polomer");
        this.r = r;
    }
    public double Area() => Math.PI * r * r;
}

public sealed class Rectangle : IShape {
    private readonly double w, h;
    public Rectangle(double w, double h) { this.w = w; this.h = h; }
    public double Area() => w * h;
}

double TotalArea(IEnumerable<IShape> shapes) {
    double sum = 0;
    foreach (var s in shapes) sum += s.Area(); // polymorfni volani
    return sum;
}
\end{lstlisting}

\textbf{(c)} \emph{Dynamický polymorfismus} se realizuje \emph{tabulkou virtuálních metod} (vtable): každý objekt drží ukazatel na tabulku své třídy, volání virtuální metody = skok přes položku tabulky (vybere se override potomka). \emph{Generické typy} umožní psát kód nezávislý na konkrétním typu (\texttt{List<T>}) bez ztráty typové kontroly a bez boxingu u hodnotových typu. \emph{Hodnotové} typy (\texttt{struct}, \texttt{int}) se kopírují podle hodnoty (typicky na zásobníku), \emph{referenční} (\texttt{class}) žijí na haldě a předává se reference; spravuje je garbage collector.
\end{reseni}

\kalibrace{Programovací jazyky / OO návrh je v moderním formátu vzácnější, ale když padne, bývá implementační. Definice (a) + funkční kostra hierarchie (b) tě udrží nad nulou; jazyk je C\# (volba majitele).}
\past{Rozhraní není abstraktní třída (nemá stav/implementaci, lze jich implementovat víc). Virtuální dispatch jde podle \emph{skutečného} typu - jinak by polymorfismus nefungoval.}

\ulohahead{společná informatika \textnormal{\footnotesize[úroveň 2]}}{Jednoduchý správce paměti (alokátor haldy)}
\begin{zadani}
Implementujte jednoduchý alokátor nad souvislým polem bajtu (haldou).
\begin{enumerate}[label=(\alph*)]
  \item \bod{1} Popište reprezentaci: blok s hlavičkou (velikost + příznak volný/obsazený) a volný seznam.
  \item \bod{2} Napište \texttt{Alloc} strategií first-fit (včetně rozdělení bloku) a \texttt{Free} (včetně slučování sousedu).
  \item \bod{3} Vysvětlete vnitřní a vnější fragmentaci a roli zarovnání (alignment).
\end{enumerate}
\end{zadani}

\begin{reseni}
\textbf{(a)} Halda je posloupnost bloku; každý má \emph{hlavičku} s velikostí a příznakem, zda je volný. Volné bloky lze provázat do \emph{volného seznamu}. Z ukazatele na data se k hlavičce dostaneme odečtením její velikosti.

\textbf{(b)}
\begin{lstlisting}
// hlavicka: int Size; bool Free; (ulozena pred daty)
IntPtr Alloc(int n) {
    n = Align(n, 8);                       // zarovnani pozadavku
    for (var b = first; b != null; b = b.Next)
        if (b.Free && b.Size >= n) {
            if (b.Size >= n + MinBlock)    // rozdel velky blok
                Split(b, n);
            b.Free = false;
            return b.Data;
        }
    return IntPtr.Zero;                     // nenalezeno
}

void Free(Block b) {
    b.Free = true;
    if (b.Next != null && b.Next.Free) Coalesce(b, b.Next); // sluc vpravo
    if (b.Prev != null && b.Prev.Free) Coalesce(b.Prev, b); // sluc vlevo
}
\end{lstlisting}

\textbf{(c)} \emph{Vnitřní fragmentace}: blok je větší než požadavek (kvuli zarovnání / minimální velikosti), nevyužitý zbytek leží uvnitř bloku. \emph{Vnější fragmentace}: volná paměť je roztříštěná do malých kousku, takže velký souvislý požadavek selže, ač je celkem dost místa. \emph{Zarovnání}: data daného typu musí začínat na adrese dělitelné jejich velikostí (rychlejší/korektní přístup CPU), proto se velikosti zaokrouhlují nahoru.
\end{reseni}

\begin{jadro}
Blok = hlavička (velikost + volný?) + data. First-fit: první dost velký volný blok, případně rozděl. Free: označ volný a slučuj se sousedy. Fragmentace: vnitřní (zbytek v bloku) vs vnější (roztříštěné volno).
\end{jadro}

\kalibrace{Implementační otázka ze společné informatiky (alokátor) padla 2025-léto; je to typ, na kterém ,,naoko připravený'' propadá. Definice + kostra Alloc/Free = vyhnutí se nule.}
\past{Po \texttt{Free} nezapomeň slučovat sousední volné bloky (jinak vnější fragmentace roste). Hlavička zabírá místo - počítej s ní u malých alokací.}

\ulohahead{společná informatika \textnormal{\footnotesize[úroveň 2]}}{Souborový systém typu FAT}
\begin{zadani}
Souborový systém ukládá soubor jako řetěz bloku spojený tabulkou FAT (každá položka udává číslo \emph{dalšího} bloku, nebo značku konce \texttt{EOF}, nebo \texttt{FREE}).
\begin{enumerate}[label=(\alph*)]
  \item \bod{1} Popište, jak FAT reprezentuje soubor a jak se pozná konec a volné bloky.
  \item \bod{2} Soubor začíná blokem 2 a tabulka je: $\text{FAT}[2]{=}5$, $\text{FAT}[5]{=}7$, $\text{FAT}[7]{=}\texttt{EOF}$. Vypište bloky souboru v pořadí. Jak se přidá blok na konec?
  \item \bod{3} Jaké jsou nevýhody FAT (fragmentace, náhodný přístup) a jak je řeší inody/extenty?
\end{enumerate}
\end{zadani}

\begin{reseni}
\textbf{(a)} Adresář drží u souboru číslo \emph{prvního} bloku. Položka $\text{FAT}[i]$ udává číslo bloku následujícího po bloku $i$; speciální hodnota \texttt{EOF} značí poslední blok, \texttt{FREE} značí volný blok. Řetězením od prvního bloku se přečte celý soubor.

\textbf{(b)} Od bloku 2: $2\to\text{FAT}[2]{=}5\to\text{FAT}[5]{=}7\to\text{FAT}[7]{=}\texttt{EOF}$. Bloky souboru: \textbf{2, 5, 7}. Přidání bloku: najdi volný blok (např.\ $k$ s $\text{FAT}[k]{=}\texttt{FREE}$), nastav $\text{FAT}[7]{=}k$ a $\text{FAT}[k]{=}\texttt{EOF}$.

\textbf{(c)} \emph{Nevýhody}: soubor bývá rozházený po disku (vnější fragmentace, pomalé čtení), a \emph{náhodný přístup} na $n$-tý blok vyžaduje projít řetěz od začátku ($O(n)$). \emph{Inody} drží pole přímých (a nepřímých) odkazu na bloky, takže přístup k libovolnému bloku je rychlý; \emph{extenty} popisují souvislé úseky (start + délka), což šetří metadata a zrychluje sekvenční čtení.
\end{reseni}

\begin{jadro}
FAT[$i$] = další blok řetězu (nebo EOF/FREE). Čtení = řetěz od prvního bloku z adresáře. Slabina: náhodný přístup $O(n)$ a fragmentace; inody/extenty to řeší přímými odkazy/souvislými úseky.
\end{jadro}

\kalibrace{FAT / interní struktura souborového systému padla 2020-podzim. Sledování řetězu bloku + nevýhody = 2 body z implementační otázky.}
\past{FAT položka ukazuje na \emph{další} blok, není to bitmapa volných (to je jiná struktura). Náhodný přístup je u FAT lineární, ne konstantní.}

\ulohahead{společná informatika \textnormal{\footnotesize[úroveň 2]}}{Dvouúrovňové stránkování a překlad adresy}
\begin{zadani}
16bitová virtuální adresa, velikost stránky 256\,B. Číslo stránky je rozděleno na dvě úrovně po 4 bitech. Virtuální adresa je \texttt{0x1234}.
\begin{enumerate}[label=(\alph*)]
  \item \bod{1} Rozdělte adresu na index L1, index L2 a offset.
  \item \bod{2} L1 tabulka v položce \texttt{0x1} ukazuje na L2 tabulku, jejíž položka \texttt{0x2} udává rámec \texttt{0x7}. Spočtěte fyzickou adresu.
  \item \bod{3} Proč se používá víceúrovňová tabulka místo jedné ploché? K čemu je TLB?
\end{enumerate}
\end{zadani}

\begin{reseni}
\textbf{(a)} Stránka 256\,B $\Rightarrow$ offset má $\log_2 256=8$ bitu. Zbývá 8 bitu čísla stránky, rozdělené na L1 (4 bity) a L2 (4 bity). \texttt{0x1234} $=$ \texttt{0001\,0010\,0011\,0100}$_2$: offset (dolních 8 bitu) $=$ \texttt{0x34}, číslo stránky (horních 8) $=$ \texttt{0x12}, tj.\ L1 $=$ \texttt{0x1}, L2 $=$ \texttt{0x2}.

\textbf{(b)} L1[\texttt{0x1}] $\to$ L2 tabulka; L2[\texttt{0x2}] $=$ rámec \texttt{0x7}. Fyzická adresa $=$ (rámec $\ll 8$) $\,|\,$ offset $=$ \texttt{0x7}\,$\ll8\ |\ $\texttt{0x34} $=$ \texttt{0x0734}.

\textbf{(c)} Jedna plochá tabulka by musela mít položku pro \emph{každou} stránku adresového prostoru (i nepoužité), což je u velkých prostoru obrovské. \emph{Víceúrovňová} tabulka alokuje L2 podtabulky jen pro skutečně použité oblasti (řídký adresový prostor) - šetří paměť. \emph{TLB} (translation lookaside buffer) je malá rychlá cache nedávných překladu stránka$\to$rámec, aby se nemusela při každém přístupu procházet tabulka v paměti.
\end{reseni}

\begin{jadro}
Adresa = (číslo stránky $\Vert$ offset); offset má $\log_2$(velikost stránky) bitu. Číslo stránky se u víceúrovňové tabulky dělí na indexy úrovní. Fyzická $=$ rámec $\ll$ (bity offsetu) $|$ offset. Víceúrovňová = šetří paměť u řídkého prostoru; TLB = cache překladu.
\end{jadro}

\kalibrace{Překlad adres je opakovaný výpočet (architektura/OS ~60\,\%); dvouúrovňová varianta je těžší než v úrovni 1. Rozdělení adresy + výpočet rámce = 2-3 body.}
\past{Offset má tolik bitu, kolik je $\log_2$(velikost stránky). Fyzická adresa skládá \emph{rámec} (ne číslo stránky) s offsetem. Bity čísla stránky se dělí mezi úrovně.}

\ulohahead{společná informatika \textnormal{\footnotesize[úroveň 2]}}{Překlad konstrukcí vyššího jazyka do instrukcí}
\begin{zadani}
\begin{enumerate}[label=(\alph*)]
  \item \bod{1} Jak se na úrovni instrukcí procesoru realizuje přiřazení, podmínka a cyklus?
  \item \bod{2} Přeložte do pseudo-assembleru cyklus \texttt{i=0; while (i<n) \{ sum += a[i]; i++; \}}.
  \item \bod{3} Jak se realizuje volání funkce (zásobníkový rámec, \texttt{call}/\texttt{ret}, návratová hodnota)?
\end{enumerate}
\end{zadani}

\begin{reseni}
\textbf{(a)} \emph{Přiřazení} = výpočet do registru a \texttt{store} do paměti (nebo \texttt{load} při čtení). \emph{Podmínka} = porovnání (\texttt{cmp}) následované podmíněným skokem (\texttt{jl}, \texttt{je}, \dots). \emph{Cyklus} = test podmínky + podmíněný skok zpět na začátek těla.

\textbf{(b)}
\begin{lstlisting}
        mov   i, 0            ; i = 0
        mov   sum, 0          ; sum = 0
loop:   cmp   i, n            ; porovnej i a n
        jge   end             ; if (i >= n) goto end
        mov   r1, a[i]        ; r1 = a[i]   (adresa = base + i*velikost)
        add   sum, r1         ; sum += r1
        add   i, 1            ; i++
        jmp   loop            ; zpet na test
end:    ...
\end{lstlisting}

\textbf{(c)} \emph{Volání funkce}: volající uloží argumenty (do registru / na zásobník) dle volací konvence a provede \texttt{call} (uloží návratovou adresu na zásobník a skočí). Volaná funkce si vytvoří \emph{zásobníkový rámec} (uloží starý base pointer, vyhradí místo pro lokální proměnné), vykoná tělo, výsledek dá do dohodnutého registru (např.\ \texttt{eax}), obnoví rámec a \texttt{ret} (vyzvedne návratovou adresu a skočí zpět).
\end{reseni}

\begin{jadro}
Přiřazení = load/store; podmínka = \texttt{cmp} + podmíněný skok; cyklus = skok zpět na test. Volání = \texttt{call} (uloží návratovou adresu) + rámec na zásobníku + \texttt{ret}; výsledek v dohodnutém registru.
\end{jadro}

\kalibrace{Překlad konstrukcí na instrukce padl 2022-léto. Cyklus/podmínka -> skoky a princip volání = 2 body z implementační/architektonické otázky.}
\past{Podmínka cyklu se testuje na \emph{začátku} (u \texttt{while}); skok je \emph{podmíněný} ven a \emph{nepodmíněný} zpět. \texttt{call} ukládá návratovou adresu, \texttt{ret} ji vyzvedává.}

\ulohahead{společná informatika \textnormal{\footnotesize[úroveň 2]}}{Pumping lemma pro bezkontextové jazyky a Master theorem}
\begin{zadani}
\begin{enumerate}[label=(\alph*)]
  \item \bod{1} Zformulujte pumping lemma pro bezkontextové jazyky (CFL).
  \item \bod{2} Dokažte, že $L=\{a^n b^n c^n\mid n\ge0\}$ není bezkontextový.
  \item \bod{3} Vysvětlete hlavní myšlenku Master theoremu (rekurzní strom) a jeho tři případy. \emph{(Požadavek uvádí Master theorem ,,bez důkazu'' - formální důkaz se nevyžaduje, jen pochopení.)}
\end{enumerate}
\end{zadani}

\begin{reseni}
\textbf{(a)} Pro každý CFL existuje $p$ tak, že každé slovo $w\in L$ s $|w|\ge p$ lze psát $w=uvxyz$, kde $|vxy|\le p$, $|vy|\ge1$ a pro každé $i\ge0$ je $uv^ixy^iz\in L$. (Pumpují se \emph{dva} úseky najednou.)

\textbf{(b)} Sporem: nechť $L$ je CFL s konstantou $p$ a $w=a^pb^pc^p$. Rozklad $w=uvxyz$ s $|vxy|\le p$ znamená, že $vxy$ zasahuje nejvýše \emph{dva} ze tří bloku ($a,b,c$) - nemůže obsahovat $a$ i $c$, neboť mezi nimi je $p$ symbolu $b$. Napumpováním ($i=2$) se zvýší počet symbolu jen v jednom či dvou blocích, takže přestanou být všechny tři stejně početné: $uv^2xy^2z\notin L$. Spor, $L$ není CFL.

\textbf{(c)} \emph{Master theorem} ($T(n)=aT(n/b)+f(n)$) přes \emph{rekurzní strom}: na úrovni $i$ je $a^i$ podproblému velikosti $n/b^i$, práce mimo rekurzi na úrovni $i$ je $a^i f(n/b^i)$; listu je $a^{\log_b n}=n^{\log_b a}$. Sečtením úrovní rozhoduje, zda dominují listy ($n^{\log_b a}$), kořen ($f(n)$), nebo jsou úrovně vyrovnané (pak faktor $\log n$). Konkrétně pro $f(n)=\Theta(n^c)$ ($a\ge1,b>1$) porovnej $c$ s $\log_b a$: $c<\log_b a\Rightarrow\Theta(n^{\log_b a})$ (listy); $c=\log_b a\Rightarrow\Theta(n^{c}\log n)$; $c>\log_b a\Rightarrow\Theta(n^{c})$ (kořen, za podmínky regularity). To jsou tři případy věty.
\end{reseni}

\begin{jadro}
CFL pumping: $w=uvxyz$, $|vxy|\le p$, $|vy|\ge1$, $uv^ixy^iz\in L$. $\{a^nb^nc^n\}$: $vxy$ pokryje max 2 bloky $\Rightarrow$ pumpování rozváží třetí. Master theorem: porovnej $n^{\log_b a}$ (listy) s $f(n)$.
\end{jadro}

\kalibrace{Důkaz nebezkontextovosti (přes pumping) podpírá zařazení do Chomského hierarchie - to je v požadavku. Master theorem se dle požadavku umí \emph{použít} (bez důkazu); intuice rekurzního stromu je jen pro pochopení.}
\past{CFL pumping pumpuje \emph{dva} úseky ($v$ a $y$) zároveň, na rozdíl od regulárního (jeden). Klíč u $a^nb^nc^n$ je $|vxy|\le p$ (nepokryje krajní bloky současně).}

\ulohahead{společná informatika \textnormal{\footnotesize[úroveň 2]}}{NP-úplnost: redukce 3-SAT na kliku}
\begin{zadani}
\begin{enumerate}[label=(\alph*)]
  \item \bod{1} Co je potřeba dokázat, aby byl problém NP-úplný? Co je polynomiální redukce?
  \item \bod{2} Popište redukci 3-SAT $\le_p$ KLIKA (konstrukci grafu).
  \item \bod{3} Zdůvodněte korektnost (formule splnitelná $\iff$ graf má kliku dané velikosti) a polynomialitu.
\end{enumerate}
\end{zadani}

\begin{reseni}
\textbf{(a)} Problém $B$ je NP-úplný, pokud (1) $B\in\mathrm{NP}$ (existuje polynomiálně ověřitelný certifikát) a (2) $B$ je \emph{NP-těžký}: každý problém z NP se na $B$ polynomiálně redukuje ($\forall K\in\mathrm{NP}:K\le_p B$). V praxi (díky tranzitivitě $\le_p$) stačí zredukovat jeden \emph{známý} NP-úplný problém na $B$. \emph{Polynomiální redukce} $A\le_p B$ je v polynomiálním čase vyčíslitelná funkce převádějící instance $A$ na instance $B$ se zachováním odpovědi ano/ne.

\textbf{(b)} Dáno 3-CNF s $m$ klauzulemi. Pro každou klauzuli vytvoř \emph{trojici vrcholu} (po jednom za každý literál). Hranou spoj dva vrcholy z \emph{různých} klauzulí, pokud si jejich literály \emph{neodporují} (nejsou to $x$ a $\neg x$). Uvnitř téže klauzule hrany nevedou. Ptáme se na kliku velikosti $m$.

\textbf{(c)} \emph{Korektnost}: klika velikosti $m$ musí obsahovat právě jeden vrchol z každé klauzule (uvnitř klauzule hrany nejsou) a vybrané literály se navzájem neodporují - lze je tedy všechny nastavit na \emph{pravda} konzistentním ohodnocením, které splní každou klauzuli. Naopak ze splňujícího ohodnocení vyber v každé klauzuli jeden pravdivý literál; tyto vrcholy jsou po dvou spojené (neodporují si) a tvoří kliku velikosti $m$. \emph{Polynomialita}: graf má $3m$ vrcholu a $O(m^2)$ hran, konstrukce běží v polynomiálním čase.
\end{reseni}

\begin{jadro}
NP-úplný = v NP $+$ NP-těžký (redukce ze známého NP-úplného). 3-SAT$\le_p$KLIKA: vrchol za literál v klauzuli, hrana mezi neodporujícími literály různých klauzulí, hledej kliku velikosti $=$ počet klauzulí.
\end{jadro}

\kalibrace{Konkrétní redukce (důkaz NP-úplnosti) je legální 3.\ bod u otázky o složitosti. Úroveň 2 přidává jeden vzorový převod.}
\past{Směr redukce: ze \emph{známého} těžkého (3-SAT) na \emph{nový} (KLIKA). Hrany vedou mezi \emph{neodporujícími} literály \emph{různých} klauzulí; klika má velikost počtu klauzulí.}

\ulohahead{společná informatika \textnormal{\footnotesize[úroveň 2]}}{Dolní odhad porovnávacího třídění}
\begin{zadani}
\begin{enumerate}[label=(\alph*)]
  \item \bod{1} Popište model rozhodovacího stromu pro porovnávací třídicí algoritmy.
  \item \bod{2} Dokažte dolní odhad $\Omega(n\log n)$ na počet porovnání.
  \item \bod{3} Proč platí dolní mez jen pro \emph{porovnávací} model? \emph{(Nesrovnávací algoritmy jako counting/radix sort ji obcházejí - stačí zmínit.)}
\end{enumerate}
\end{zadani}

\begin{reseni}
\textbf{(a)} Porovnávací algoritmus lze popsat \emph{rozhodovacím stromem}: vnitřní uzel je porovnání dvou prvku ($a_i\lessgtr a_j$), větve odpovídají výsledku, list odpovídá jedné konkrétní permutaci (výslednému uspořádání). Běh algoritmu = cesta od kořene k listu; počet porovnání v nejhorším případě $=$ výška stromu.

\textbf{(b)} Pro $n$ prvku existuje $n!$ možných uspořádání, a protože algoritmus musí umět rozlišit každé, má strom aspoň $n!$ listu. Binární strom s $L$ listy má výšku $\ge\log_2 L$, tedy výška $\ge\log_2(n!)$. Stirlingovým vzorcem $\log_2(n!)=\Theta(n\log n)$. Výška = nejhorší počet porovnání, takže každý porovnávací algoritmus potřebuje $\Omega(n\log n)$ porovnání.

\textbf{(c)} \emph{Nesrovnávací} algoritmy nepoužívají porovnání dvojic, ale strukturu klíču: \emph{counting sort} pro celá čísla z malého rozsahu $k$ spočítá četnosti a umístí prvky, $O(n+k)$; \emph{radix sort} třídí po číslicích pomocí stabilního counting sortu, $O(d\,(n+k))$. Tím obcházejí dolní mez $\Omega(n\log n)$, která platí jen pro \emph{porovnávací} model.
\end{reseni}

\begin{jadro}
Rozhodovací strom má $\ge n!$ listu $\Rightarrow$ výška $\ge\log_2(n!)=\Theta(n\log n)$ $\Rightarrow$ porovnávací třídění je $\Omega(n\log n)$. Counting/radix sort tuto mez obchází (nepoužívají porovnání), za cenu předpokladu o rozsahu klíču.
\end{jadro}

\kalibrace{Dolní odhad třídění je opakované téma (třídění) a klasický důkaz na 3.\ bod. Úroveň 2 přidává jeho vypracování.}
\past{Mez $\Omega(n\log n)$ platí jen pro \emph{porovnávací} model; counting/radix ji nelámou ,,kouzlem'', ale jiným modelem (předpoklad o klíčích). $\log_2(n!)=\Theta(n\log n)$, ne $\Theta(n)$.}

\ulohahead{společná informatika \textnormal{\footnotesize[úroveň 3]}}{Nativní vs interpretovaný běh, JIT/AOT, sestavení programu}
\begin{zadani}
\begin{enumerate}[label=(\alph*)]
  \item \bod{1} Vysvětlete rozdíl nativní vs interpretovaný běh, co je bytecode a interpret.
  \item \bod{2} Vysvětlete JIT a AOT překlad a proces sestavení (oddělený překlad, linkování, staticky vs dynamicky linkované knihovny).
  \item \bod{3} Stručně: běhové prostředí procesu, volací konvence a relokace (position-independent code).
\end{enumerate}
\end{zadani}

\begin{reseni}
\textbf{(a)} \emph{Nativní} běh: program je přeložen do strojového kódu procesoru a CPU ho vykonává přímo (rychlé, závislé na platformě). \emph{Interpretovaný} běh: \emph{interpret} čte a vykonává program (zdroj nebo \emph{bytecode} = přenositelný mezikód virtuálního stroje, např.\ CIL/JVM) instrukci po instrukci - přenositelné, ale pomalejší.

\textbf{(b)} \emph{JIT} (just-in-time): bytecode se překládá do nativního kódu \emph{za běhu} těsně před prvním spuštěním metody (kombinuje přenositelnost s rychlostí, může optimalizovat dle běhových dat). \emph{AOT} (ahead-of-time): překlad do nativního kódu \emph{před} během (rychlý start, bez JIT režie). \emph{Sestavení}: zdroje se \emph{odděleně} přeloží na objektové soubory a \emph{linker} je spojí (vyřeší odkazy mezi moduly) do spustitelného souboru. \emph{Staticky} linkované knihovny se vkopírují do binárky (větší, soběstačná); \emph{dynamické} (.dll/.so) se zavádějí za běhu (sdílené, menší binárka, závislost na verzi).

\textbf{(c)} \emph{Běhové prostředí procesu}: adresový prostor (kód, data, halda, zásobník), vazba na OS přes systémová volání. \emph{Volací konvence} určuje, jak se předávají argumenty (registry/zásobník), kdo uklízí zásobník a kde je návratová hodnota. \emph{Relokace} upravuje adresy při zavedení modulu na jinou adresu; \emph{position-independent code} (PIC) používá relativní adresování, takže kód lze zavést na různé adresy s žádnými/menšími relokacemi v sekci kódu (dynamické relokace symbolu přes GOT/PLT mohou zustat); typické pro sdílené knihovny.
\end{reseni}

\begin{jadro}
Nativní = přímo CPU; interpret = vykonává bytecode. JIT = překlad za běhu, AOT = předem. Sestavení: oddělený překlad $\to$ objektové soubory $\to$ linker. Statické knihovny v binárce, dynamické (.dll/.so) zaváděné za běhu.
\end{jadro}

\kalibrace{Řízení překladu a sestavení (JIT/AOT, linkování) je v požadavku programovacích jazyku, ale v moderních zkouškách vzácné - úroveň 3 to pokrývá pro úplnost.}
\past{JIT překládá \emph{za běhu}, ne před ním (to je AOT). Dynamické knihovny šetří paměť sdílením, ale přinášejí závislost na verzi (,,DLL hell'').}

\ulohahead{společná informatika \textnormal{\footnotesize[úroveň 3]}}{Správa zdrojů a ošetření chyb (C\#)}
\begin{zadani}
\begin{enumerate}[label=(\alph*)]
  \item \bod{1} Co je správa životního cyklu zdroje (soubor, zámek, spojení) a jak souvisí s výjimkami?
  \item \bod{2} Vysvětlete \texttt{IDisposable} a \texttt{using} v C\# (a obdobu RAII / try-with-resources). Co je deterministické uvolnění?
  \item \bod{3} Vysvětlete reference counting a finalizéry: kdy se volají a jaká jsou jejich úskalí.
\end{enumerate}
\end{zadani}

\begin{reseni}
\textbf{(a)} Zdroj (soubor, síťové spojení, zámek) je nutné po použití \emph{uvolnit}, i když nastane chyba. Výjimka může běžný tok přerušit, takže uvolnění nesmí být jen za ,,šťastnou'' cestou - musí proběhnout i při propagaci výjimky (jinak hrozí únik zdroje / zaseknutý zámek).

\textbf{(b)} V C\# typ vlastnící zdroj implementuje \texttt{IDisposable} s metodou \texttt{Dispose()}. Konstrukce \texttt{using} zaručí volání \texttt{Dispose()} při opuštění bloku (i přes výjimku), což je \emph{deterministické} uvolnění (přesně v daném bodě), na rozdíl od nedeterministického garbage collectoru.
\begin{lstlisting}
using (var f = File.OpenRead(path)) {
    // prace se souborem
}   // zde se vola f.Dispose() i kdyz vyletela vyjimka
\end{lstlisting}
Obdoby: \emph{RAII} v C++ (destruktor uvolní zdroj při opuštění rozsahu), \emph{try-with-resources} v Javě; bez nich slouží \texttt{try/finally}.

\textbf{(c)} \emph{Reference counting} drží u objektu počet odkazu; při poklesu na 0 se objekt uvolní (deterministicky), ale neumí uvolnit \emph{cykly}. \emph{Finalizér} (\texttt{\textasciitilde Trida} / \texttt{Finalize}) volá GC \emph{nedeterministicky} před uvolněním objektu jako záchrana pro nespravované zdroje - není zaručeno kdy (ani zda) se zavolá, proto se na něj nespoléhá; správně je \texttt{Dispose()} (vzor dispose + finalizér jako pojistka).
\end{reseni}

\begin{jadro}
Zdroj uvolni i při výjimce: C\# \texttt{using}+\texttt{IDisposable} (deterministicky), C++ RAII (destruktor), jinak \texttt{try/finally}. Reference counting neuvolní cykly; finalizér běží nedeterministicky přes GC - nespoléhat, použít \texttt{Dispose}.
\end{jadro}

\kalibrace{Správa zdroju a výjimky jsou v požadavku programovacích jazyku; jazyk je C\# (volba majitele). Úroveň 3 doplňuje tuto tail oblast.}
\past{Finalizér $\neq$ destruktor v C++ - běží nedeterministicky a není zaručen. Deterministické uvolnění zajistí \texttt{using}/\texttt{Dispose}, ne GC.}

\ulohahead{společná informatika \textnormal{\footnotesize[úroveň 3]}}{Generické typy a funkcionální prvky (C\#)}
\begin{zadani}
\begin{enumerate}[label=(\alph*)]
  \item \bod{1} Co jsou generické typy a proč se používají (typová bezpečnost, výkon)?
  \item \bod{2} Co jsou typy reprezentující funkce: delegáti, \texttt{Func}/\texttt{Action}, lambda funkce?
  \item \bod{3} Napište generickou metodu v C\# a vysvětlete rozdíl mezi generikami (C\#) a šablonami (C++).
\end{enumerate}
\end{zadani}

\begin{reseni}
\textbf{(a)} \emph{Generické typy} parametrizují třídu/metodu typem (\texttt{List<T>}, \texttt{Dictionary<K,V>}), takže lze psát kód nezávislý na konkrétním typu \emph{bez ztráty typové kontroly} (vše se ověří při překladu) a \emph{bez boxingu} hodnotových typu (na rozdíl od kontejneru \texttt{object}). Výsledek: méně duplicity, bezpečnější a rychlejší kód.

\textbf{(b)} \emph{Delegát} je typ reprezentující odkaz na metodu (typový ukazatel na funkci). \texttt{Func<...,TResult>} a \texttt{Action<...>} jsou předdefinované generické delegáty (s/bez návratové hodnoty). \emph{Lambda} je stručný zápis anonymní funkce, např.\ \texttt{x => x*x}, který lze přiřadit do delegátu a předat jako argument (funkce vyššího řádu).

\textbf{(c)}
\begin{lstlisting}
T Max<T>(T a, T b) where T : IComparable<T>
    => a.CompareTo(b) >= 0 ? a : b;
// pouziti: Max(3, 7) -> 7;  Max("a","b") -> "b"
\end{lstlisting}
\emph{Generika (C\#)} jsou \emph{reifikovaná} (typové parametry existují v metadatech/IL i za běhu, lze je omezit klauzulí \texttt{where}); JIT typicky sdílí kód pro referenční typy a specializuje pro hodnotové. \emph{Šablony (C++)} jsou \emph{instanciovány při překladu} - pro každý použitý typ vznikne samostatný kód (mocnější/statický polymorfismus, ale delší překlad a horší chybové hlášky).
\end{reseni}

\begin{jadro}
Generika = parametrizace typem, typová bezpečnost bez boxingu. Delegát = typ funkce; \texttt{Func}/\texttt{Action}; lambda \texttt{x=>...}. C\# generika: jeden typ + běhová podpora + \texttt{where}; C++ šablony: instanciace při překladu.
\end{jadro}

\kalibrace{Generické typy a funkcionální prvky jsou v požadavku programovacích jazyku; úroveň 3 doplňuje tuto oblast v C\#.}
\past{Generika C\# nejsou ,,smazání typu'' jako v Javě - typové parametry jsou dostupné i za běhu. Lambda je výraz, delegát je typ, který ji drží.}

\ulohahead{společná informatika \textnormal{\footnotesize[úroveň 3]}}{Garbage collection}
\begin{zadani}
\begin{enumerate}[label=(\alph*)]
  \item \bod{1} Co je garbage collection a princip dosažitelnosti (kořeny, živé objekty)? Popište mark-and-sweep.
  \item \bod{2} Co je generační hypotéza a jak ji využívá generační GC?
  \item \bod{3} Porovnejte GC s reference countingem (problém cyklu) a vysvětlete stop-the-world a kompakci.
\end{enumerate}
\end{zadani}

\begin{reseni}
\textbf{(a)} \emph{Garbage collector} automaticky uvolňuje objekty, které už nejsou potřeba. \emph{Dosažitelnost}: objekt je \emph{živý}, je-li dosažitelný z \emph{kořenu} (lokální proměnné, registry, statická pole) přes řetěz referencí; nedosažitelné objekty jsou odpad. \emph{Mark-and-sweep}: (1) \emph{mark} - projdi z kořenu a označ všechny dosažitelné objekty; (2) \emph{sweep} - uvolni neoznačené.

\textbf{(b)} \emph{Generační hypotéza}: většina objektu umírá mladá (krátká životnost). \emph{Generační GC} proto dělí haldu na generace; \emph{mladou} generaci sbírá často a rychle (málo přeživších se povýší do starší), \emph{starou} jen občas. Tím se zlevní typický průběh (nesbírá se pokaždé celá halda).

\textbf{(c)} \emph{Reference counting} uvolňuje deterministicky při poklesu počtu odkazu na 0, ale \emph{neuvolní cykly} (vzájemně se odkazující objekty mají počet $>0$, ač jsou nedosažitelné) - sledovací GC (mark-sweep) cykly zvládá, protože jde od kořenu. \emph{Stop-the-world}: po dobu (části) sběru se pozastaví běh aplikace, aby se graf objektu neměnil. \emph{Kompakce}: po sweepu se živé objekty posunou k sobě, čímž se odstraní fragmentace haldy a zrychlí alokace.
\end{reseni}

\begin{jadro}
GC uvolní nedosažitelné objekty (od kořenu). Mark-and-sweep: označ dosažitelné, ukliď zbytek. Generační GC těží z ,,objekty umírají mladé''. Reference counting neuvolní cykly; GC ano. Stop-the-world pozastaví aplikaci, kompakce odstraní fragmentaci.
\end{jadro}

\kalibrace{Garbage collection je v požadavku (správa paměti) a objevila se v dřívějších zkouškách; úroveň 3 ji pokrývá do hloubky.}
\past{Reference counting selhává na cyklech - to je klíčový rozdíl proti sledovacímu GC. Generační GC nesbírá pokaždé celou haldu.}


\section{Specializace: Strojové učení (Téma 3)}
\ulohahead{specializace UI-SU}{Lineární a logistická regrese, regularizace}
\begin{zadani}
\begin{enumerate}[label=(\alph*)]
  \item \bod{1} Definujte model lineární regrese a chybovou funkci (MSE). Napište \emph{normální rovnice} a analytické řešení metodou nejmenších čtverců.
  \item \bod{2} Co je L2 (ridge) regularizace? Jak změní normální rovnice a jaký má vliv na koeficienty? Proč pomáhá proti přeučení a kdy je nutná (singularita $X^{\!\top}X$)?
  \item \bod{3} Logistická regrese pro binární klasifikaci: napište predikci (sigmoid), ztrátovou funkci (NLL / křížová entropie) a gradient pro jeden krok SGD. Nemusíte nic počítat numericky.
\end{enumerate}
\end{zadani}

\begin{reseni}
\textbf{(a) Lineární regrese.} Data $X\in\mathbb{R}^{n\times d}$ (řádky $x_i$), cíle $y\in\mathbb{R}^n$, model $\hat y=Xw$ (sloupec jedniček v $X$ dává bias). Chyba:
\[
\mathrm{MSE}(w)=\tfrac1n\sum_{i=1}^n (x_i^{\!\top}w-y_i)^2=\tfrac1n\lVert Xw-y\rVert^2.
\]
Minimum: $\nabla_w=\tfrac2n X^{\!\top}(Xw-y)=0$, tj.\ \emph{normální rovnice}
\[
X^{\!\top}X\,w = X^{\!\top}y \quad\Rightarrow\quad w=(X^{\!\top}X)^{-1}X^{\!\top}y \ \text{(pokud je $X^{\!\top}X$ regulární).}
\]

\textbf{(b) L2 / ridge.} Minimalizujeme $\lVert Xw-y\rVert^2+\lambda\lVert w\rVert^2$, $\lambda>0$. Normální rovnice se změní na
\[
(X^{\!\top}X+\lambda I)\,w=X^{\!\top}y,\qquad w=(X^{\!\top}X+\lambda I)^{-1}X^{\!\top}y.
\]
Matice $X^{\!\top}X+\lambda I$ je vždy regulární (pozitivně definitní), takže řešení existuje i při kolineárních příznacích nebo $d>n$, kdy je $X^{\!\top}X$ singulární. Regularizace \emph{smršťuje} koeficienty směrem k nule (snižuje rozptyl modelu za cenu malého vychýlení), čímž omezuje přeučení. Větší $\lambda$ = silnější smrštění; $\lambda$ se volí křížovou validací.

\textbf{(c) Logistická regrese.} Predikce pravděpodobnosti třídy 1:
\[
p_i=\sigma(w^{\!\top}x_i)=\frac{1}{1+e^{-w^{\!\top}x_i}}.
\]
Ztráta (negativní log-věrohodnost / binární křížová entropie):
\[
L(w)=-\sum_{i=1}^n\big[y_i\ln p_i+(1-y_i)\ln(1-p_i)\big].
\]
Gradient má jednoduchý tvar $\nabla_w L=\sum_i (p_i-y_i)\,x_i$. Krok SGD na jednom příkladu $(x_i,y_i)$:
\[
w\leftarrow w-\eta\,(p_i-y_i)\,x_i,
\]
kde $\eta$ je rychlost učení. (Pro více tříd: softmax místo sigmoidu, ztráta = kategoriální křížová entropie, gradient $(p_i-\mathbb{1}_{y_i})x_i$.)
\end{reseni}

\begin{jadro}
Lineární: normální rovnice $X^{\!\top}X\,w=X^{\!\top}y$. Ridge: $(X^{\!\top}X+\lambda I)\,w=X^{\!\top}y$ (smršťuje, vždy řešitelné). Logistická: $p=\sigma(w^{\!\top}x)$, ztráta NLL, gradient $\sum_i(p_i-y_i)x_i$, SGD krok $w\leftarrow w-\eta(p_i-y_i)x_i$.
\end{jadro}

\kalibrace{Regrese je nejčastější ML okruh (~67\,\% zkoušek) a sytí kritickou podlahu specializace ($\ge 7/12$, kde potřebuješ tři ze čtyř otázek na $\ge 2$ body). Sigmoid, NLL a normální rovnice umět zpaměti = spolehlivé 2-3 body.}
\past{,,Metoda nejmenších čtverců'' = analytické řešení přes normální rovnice; SGD je iterativní alternativa. Pleteš-li si je, ztrácíš bod. U ridge nezapomeň, že bias ($w_0$) se obvykle neregularizuje.}

\ulohahead{specializace UI-SU}{Shlukování: k-means a hierarchické shlukování}
\begin{zadani}
\begin{enumerate}[label=(\alph*)]
  \item \bod{1} Definujte úlohu shlukování (učení bez učitele) a algoritmus k-means (krok přiřazení a krok aktualizace) a jeho účelovou funkci.
  \item \bod{2} Co je inicializace k-means++? Proč ztráta v k-means monotónně klesá a proč algoritmus konverguje (do lokálního minima)? Popište aglomerativní hierarchické shlukování a dendrogram.
  \item \bod{3} Jak zvolit počet shluku $k$? Proč je k-means citlivý na škálování příznaku? Proveďte jeden krok k-means pro body $1,2,\;10,11$ na přímce s centry $c_1=0,c_2=20$.
\end{enumerate}
\end{zadani}

\begin{reseni}
\textbf{(a)} \emph{Shlukování}: rozděl neoznačená data do $k$ skupin podle podobnosti. \emph{k-means} minimalizuje vnitroshlukový součet čtvercu $J=\sum_{i}\lVert x_i-\mu_{c(i)}\rVert^2$ střídáním: (1) \emph{přiřazení} každého bodu k nejbližšímu centru; (2) \emph{aktualizace} center na průměr přiřazených bodu.

\textbf{(b)} \emph{k-means++}: počáteční centra volí postupně s pravděpodobností úměrnou $D(x)^2$ (vzdálenost k nejbližšímu už zvolenému centru), což snižuje riziko špatného startu. Oba kroky $J$ nezvyšují (přiřazení k nejbližšímu i průměr jsou optimální dílčí kroky), $J$ je zdola omezené, konfigurací je konečně mnoho $\Rightarrow$ konvergence, ale jen do \emph{lokálního} minima (proto víc restartu). \emph{Aglomerativní} shlukování začíná s každým bodem zvlášť a opakovaně spojuje dva nejbližší shluky (linkage: single/complete/average); historii zachytí \emph{dendrogram}, který se ,,uřízne'' na požadovaný počet shluku.

\textbf{(c)} $k$ se volí např.\ metodou lokte (zlom v poklesu $J$) nebo silhouette skórem. \emph{Škálování}: euklidovská vzdálenost míchá jednotky, takže příznak s velkým rozsahem dominuje; proto se data standardizují. Krok pro body $\{1,2,10,11\}$, $c_1=0,c_2=20$: přiřazení: $1,2,10,11$ jsou blíž $c_1$ než $c_2$? $|10-0|=|10-20|=10$ (remíza), $|11-0|=11>|11-20|=9$, takže 11 jde k $c_2$. Tedy $\{1,2,10\}\to c_1$, $\{11\}\to c_2$ (remízu u 10 dáme $c_1$). Aktualizace: $c_1=\tfrac{1+2+10}{3}=4{,}33$, $c_2=11$. (Další iterace už rozdělí na $\{1,2\}$ a $\{10,11\}$.)
\end{reseni}

\kalibrace{Shlukování je druhý nejčastější ML okruh (~44\,\%, k-means se opakuje). Definice + jeden krok výpočtu = 2 body do kritické podlahy specializace.}
\past{k-means najde jen \emph{lokální} minimum a předpokládá zhruba kulové shluky; není deterministický (závisí na inicializaci). Před během škáluj příznaky.}

\ulohahead{specializace UI-SU}{Evaluační metriky binárního klasifikátoru}
\begin{zadani}
\begin{enumerate}[label=(\alph*)]
  \item \bod{1} Nakreslete matici záměn a definujte accuracy, precision, recall a F1.
  \item \bod{2} Proč je accuracy zavádějící u nevyvážených tříd? Kdy upřednostnit precision a kdy recall? Z hodnot $\mathrm{TPR}$, $\mathrm{FPR}$ a počtu pozitivních/negativních rekonstruujte $\mathrm{TP,FP,TN,FN}$.
  \item \bod{3} Co zobrazuje ROC křivka a co je AUC? Spočtěte precision a recall pro $\mathrm{TP}=20,\ \mathrm{FP}=30,\ \mathrm{FN}=5$.
\end{enumerate}
\end{zadani}

\begin{reseni}
\textbf{(a)} Matice záměn (řádky skutečnost, sloupce predikce):
\[
\begin{array}{c|cc}
 & \hat{+} & \hat{-}\\\hline
+ & \mathrm{TP} & \mathrm{FN}\\
- & \mathrm{FP} & \mathrm{TN}
\end{array}
\]
$\text{accuracy}=\frac{\mathrm{TP+TN}}{\text{vše}}$, $\text{precision}=\frac{\mathrm{TP}}{\mathrm{TP+FP}}$, $\text{recall (TPR)}=\frac{\mathrm{TP}}{\mathrm{TP+FN}}$, $F_1=\frac{2\,PR}{P+R}$ (harmonický průměr).

\textbf{(b)} Při nevyvážených třídách (např.\ 1\,\% pozitivních) dá triviální klasifikátor ,,vše negativní'' accuracy 99\,\%, ač nic nenajde - proto se sleduje precision/recall. \emph{Recall} je důležitý, když je drahé minout pozitivní (nemoc, podvod); \emph{precision}, když je drahý falešný poplach (spam označený jako důležitý). Rekonstrukce z $\mathrm{TPR},\mathrm{FPR}$ a počtu $P$ (pozitivních) a $N$ (negativních): $\mathrm{TP}=\mathrm{TPR}\cdot P$, $\mathrm{FN}=P-\mathrm{TP}$, $\mathrm{FP}=\mathrm{FPR}\cdot N$, $\mathrm{TN}=N-\mathrm{FP}$.

\textbf{(c)} \emph{ROC} vynáší $\mathrm{TPR}$ proti $\mathrm{FPR}$ při měnícím se prahu klasifikace; \emph{AUC} (plocha pod ní) je pravděpodobnost, že náhodný pozitivní dostane vyšší skóre než náhodný negativní (1 = ideál, 0{,}5 = náhoda). Pro zadání: $\text{precision}=\frac{20}{20+30}=0{,}4$, $\text{recall}=\frac{20}{20+5}=0{,}8$.
\end{reseni}

\begin{jadro}
$\text{precision}=\frac{TP}{TP+FP}$, $\text{recall}=\frac{TP}{TP+FN}$, $F_1=\frac{2PR}{P+R}$ (harmonický). U nevyvážených tříd accuracy klame. Rekonstrukce: $TP=\mathrm{TPR}\cdot P$, $FP=\mathrm{FPR}\cdot N$. ROC = TPR vs FPR.
\end{jadro}

\kalibrace{Evaluační metriky jsou velmi častý ML okruh (~44\,\%) a opakovaně chtějí rekonstrukci počtu z metrik. Vzorce P/R/F1 a převod TPR/FPR umět zpaměti = spolehlivé body.}
\past{Precision a recall si nepleť (jmenovatel: u precision predikované pozitivní, u recall skutečně pozitivní). F1 je harmonický, ne aritmetický průměr.}

\ulohahead{specializace UI-SU}{Rozhodovací stromy: kritéria větvení a prořezávání}
\begin{zadani}
\begin{enumerate}[label=(\alph*)]
  \item \bod{1} Definujte rozhodovací strom a kritéria větvení: entropii a informační zisk, Gini index.
  \item \bod{2} Popište hladovou konstrukci stromu. Spočtěte informační zisk štěpení, které z uzlu se 4 kladnými a 4 zápornými příklady udělá dva čisté listy (4/0 a 0/4).
  \item \bod{3} Co je přeučení u stromu a jak mu brání prořezávání (pre-pruning vs post-pruning)?
\end{enumerate}
\end{zadani}

\begin{reseni}
\textbf{(a)} \emph{Rozhodovací strom}: vnitřní uzly testují příznak, listy přiřazují třídu (či hodnotu). \emph{Entropie} množiny $S$ s podíly tříd $p_k$: $H(S)=-\sum_k p_k\log_2 p_k$. \emph{Informační zisk} štěpení $S$ na části $S_j$: $\mathrm{IG}=H(S)-\sum_j\frac{|S_j|}{|S|}H(S_j)$. \emph{Gini}: $G(S)=1-\sum_k p_k^2$ (alternativní míra nečistoty).

\textbf{(b)} \emph{Konstrukce} (ID3/CART): v každém uzlu zvol příznak s největším informačním ziskem (resp.\ nejmenší váženou nečistotou), rozděl data a rekurzivně pokračuj, dokud nejsou listy čisté nebo nenastane zastavení. Výpočet: rodič $4{+}/4{-}$ má $H=-\tfrac12\log_2\tfrac12-\tfrac12\log_2\tfrac12=1$ bit. Děti $4/0$ a $0/4$ jsou čisté: $H=0$. $\mathrm{IG}=1-\big(\tfrac12\cdot0+\tfrac12\cdot0\big)=1$ bit (maximální možný zisk - dokonalé rozdělení).

\textbf{(c)} \emph{Přeučení}: strom roste do hloubky, učí se šum a na trénovacích datech má nulovou chybu, ale špatně generalizuje. \emph{Pre-pruning} (předčasné zastavení): omezení hloubky, minimální počet vzorku v listu, minimální zisk. \emph{Post-pruning}: nech strom dorůst a pak ořezávej podstromy, které nezlepšují chybu na validační množině (nahrazení listem).
\end{reseni}

\kalibrace{Rozhodovací stromy jsou častý okruh (~33\,\%) a vděčí za body za entropii/Gini a jeden výpočet IG. Patří do plně zvládnutého jádra specializace.}
\past{Entropie se počítá s $\log_2$ (bity). Informační zisk je \emph{vážený} podle velikostí dětí, ne prostý rozdíl entropií.}

\ulohahead{specializace UI-SU}{Kombinace modelu: bagging, boosting, náhodné lesy}
\begin{zadani}
\begin{enumerate}[label=(\alph*)]
  \item \bod{1} Co je kombinace více modelu (ensemble)? Vysvětlete rozdíl mezi baggingem a boostingem.
  \item \bod{2} Popište náhodný les (random forest) a vysvětlete, proč zlepšuje přesnost oproti jednomu stromu.
  \item \bod{3} Popište princip AdaBoostu. Za jakých podmínek ensemble pomáhá nejvíc a kdy vůbec ne (majoritní hlasování)?
\end{enumerate}
\end{zadani}

\begin{reseni}
\textbf{(a)} \emph{Ensemble} kombinuje predikce více modelu (hlasování/průměr) do silnějšího. \emph{Bagging} trénuje modely \emph{nezávisle} na bootstrap vzorcích a průměruje (snižuje hlavně \emph{rozptyl}). \emph{Boosting} trénuje modely \emph{postupně}, každý se zaměří na chyby předchozích (snižuje hlavně \emph{vychýlení}).

\textbf{(b)} \emph{Random forest} = bagging rozhodovacích stromu, kde se navíc v každém štěpení uvažuje jen náhodná podmnožina příznaku. To stromy \emph{dekoreluje}; protože průměr nekorelovaných chyb má menší rozptyl, les generalizuje lépe než jeden hluboký (přeučený) strom, a přitom téměř nepotřebuje ladění.

\textbf{(c)} \emph{AdaBoost}: udržuje váhy příkladu, v každém kole natrénuje slabý klasifikátor, zvýší váhy špatně klasifikovaných a přidá klasifikátor s váhou podle jeho přesnosti; finální predikce = vážené hlasování. Ensemble pomáhá nejvíc, jsou-li dílčí modely \emph{lepší než náhoda} a dělají \emph{různé/nekorelované} chyby. Naopak nepomůže, dělají-li všechny stejné chyby (dokonale korelované) - pak majoritní hlasování dopadne stejně jako jeden model; nejlepší případ je, když se chyby nepřekrývají a většina vždy ,,přehlasuje'' menšinu chybujících.
\end{reseni}

\begin{jadro}
Bagging = trénuj nezávisle na bootstrap vzorcích + průměruj (snižuje rozptyl). Boosting = sekvenčně, každý model na chybách předchozích (snižuje vychýlení). Random forest = bagging stromu + náhodná podmnožina příznaku (dekorelace). Pomáhá při nekorelovaných chybách.
\end{jadro}

\kalibrace{Kombinace modelu (~22\,\%) je vděčná konceptuální otázka (žádný výpočet). Bagging vs boosting + random forest = 2-3 body za vysvětlení.}
\past{Bagging $\neq$ boosting: bagging paralelně/nezávisle (rozptyl), boosting sekvenčně na chybách (vychýlení). Les pomáhá díky \emph{dekorelaci} stromu, ne jen jejich počtu.}

\ulohahead{specializace UI-SU}{Metoda podpůrných vektoru (SVM)}
\begin{zadani}
\begin{enumerate}[label=(\alph*)]
  \item \bod{1} Popište SVM pro lineárně separabilní třídy: co je rozdělující nadrovina, margin a podpůrné vektory (support vectors)?
  \item \bod{2} Co je měkký margin a parametr $C$? Jaký je kompromis mezi šířkou marginu a počtem chyb?
  \item \bod{3} K čemu slouží jádrové funkce (kernel trick) a jak se chová RBF jádro? Jak SVM klasifikuje do více tříd?
\end{enumerate}
\end{zadani}

\begin{reseni}
\textbf{(a)} SVM hledá \emph{rozdělující nadrovinu} $w^{\!\top}x+b=0$ s \emph{maximálním marginem} (pásem) mezi třídami. Margin je vzdálenost nadroviny k nejbližším bodum; ty body leží na okraji a nazývají se \emph{podpůrné vektory} - jen ony určují řešení (ostatní lze odebrat beze změny). Maximalizace marginu $=$ minimalizace $\tfrac12\lVert w\rVert^2$ za podmínek $y_i(w^{\!\top}x_i+b)\ge1$.

\textbf{(b)} Když data nejsou separabilní, zavede se \emph{měkký margin} s prom.\ uvolnění $\xi_i\ge0$: minimalizujeme $\tfrac12\lVert w\rVert^2+C\sum_i\xi_i$. Parametr $C$ váží chyby proti šířce marginu: \emph{velké} $C$ = málo chyb, úzký margin (riziko přeučení); \emph{malé} $C$ = širší margin, toleruje víc chyb (lepší generalizace).

\textbf{(c)} \emph{Jádrový trik}: nahradí skalární součin $x_i^{\!\top}x_j$ jádrem $K(x_i,x_j)$, čímž SVM počítá v (implicitně vysokorozměrném) prostoru rysu, aniž ho explicitně sestrojí - umožní nelineární hranice. \emph{RBF} $K(x,z)=\exp(-\gamma\lVert x-z\rVert^2)$ měří podobnost lokálně; malé $\gamma$ = hladká hranice, velké $\gamma$ = klikatá (přeučení). \emph{Více tříd}: kombinací binárních klasifikátoru one-vs-rest nebo one-vs-one.
\end{reseni}

\kalibrace{SVM (~12\,\%) je vzácnější, ale když padne, je konceptuální (,,popište, nakreslete hranici, vliv bodu, kdy RBF''). Definice marginu + role $C$/jádra = 2-bodová záloha.}
\past{Hranici určují jen podpůrné vektory. Velké $C$ = méně tolerance chyb (užší margin), ne naopak; RBF s velkým $\gamma$ se přeučuje.}

\ulohahead{specializace UI-SU}{Analýza hlavních komponent (PCA)}
\begin{zadani}
\begin{enumerate}[label=(\alph*)]
  \item \bod{1} Co je PCA a k čemu slouží? Co jsou hlavní komponenty a jak souvisí s kovarianční maticí?
  \item \bod{2} Popište postup PCA (centrování, kovariance, vlastní rozklad / SVD, projekce na top-$k$). Co je vysvětlený rozptyl?
  \item \bod{3} Jaká je souvislost PCA a SVD? Proč je nutné data před PCA škálovat a co PCA optimalizuje (rekonstrukční chyba)?
\end{enumerate}
\end{zadani}

\begin{reseni}
\textbf{(a)} \emph{PCA} je lineární redukce dimenze: hledá nové ortogonální osy (\emph{hlavní komponenty}), podél nichž mají data největší rozptyl, a projektuje na prvních $k$ z nich. Hlavní komponenty jsou \emph{vlastní vektory kovarianční matice} dat; příslušná vlastní čísla udávají rozptyl podél nich.

\textbf{(b)} Postup: (1) \emph{centruj} data (odečti průměr), případně standardizuj; (2) spočti \emph{kovarianční matici} $C=\tfrac1n X^{\!\top}X$; (3) najdi její \emph{vlastní čísla a vektory} (nebo SVD matice $X$); (4) seřaď komponenty podle vlastních čísel a \emph{projektuj} na prvních $k$. \emph{Vysvětlený rozptyl} komponenty = $\lambda_i/\sum_j\lambda_j$; volíme $k$ tak, aby pokryl např.\ 95\,\% rozptylu.

\textbf{(c)} Pro centrovaná data $X=U\Sigma V^{\!\top}$ (SVD) jsou hlavní komponenty sloupce $V$ a singulární hodnoty $\sigma_i$ souvisí s vlastními čísly kovariance přes $\lambda_i=\sigma_i^2/n$; SVD je numericky stabilnější cesta. \emph{Škálování} je nutné, protože PCA reaguje na rozptyl: příznak s velkou jednotkou (rozsahem) by jinak uměle dominoval první komponentě. PCA minimalizuje \emph{rekonstrukční chybu} (součet čtvercu kolmých vzdáleností k podprostoru), což je ekvivalentní maximalizaci zachyceného rozptylu.
\end{reseni}

\kalibrace{PCA (~12\,\%) je vzácnější, ale opakované jako konceptuální + vizuální otázka. Definice (komponenty = vlastní vektory kovariance) + postup = 2-bodová záloha.}
\past{PCA hledá směry maximálního \emph{rozptylu}, ne směry oddělující třídy (to je LDA). Bez centrování/škálování vyjdou komponenty zkreslené.}

\ulohahead{specializace UI-SU}{k-nejbližších sousedu (kNN)}
\begin{zadani}
\begin{enumerate}[label=(\alph*)]
  \item \bod{1} Popište algoritmus k-nejbližších sousedu pro klasifikaci i regresi. Proč se mu říká ,,líné'' učení (lazy)?
  \item \bod{2} Jak se volí $k$ (křížová validace) a jak malé vs velké $k$ ovlivní výsledek (vychýlení vs rozptyl)? Proč je nutné škálovat příznaky?
  \item \bod{3} Co je prokletí dimenzionality a jak postihuje kNN? Jaká je výpočetní složitost predikce?
\end{enumerate}
\end{zadani}

\begin{reseni}
\textbf{(a)} kNN si \emph{pamatuje} trénovací data; pro nový bod najde $k$ nejbližších (dle metriky, typicky euklidovské) a \emph{klasifikace} = většinové hlasování jejich tříd, \emph{regrese} = průměr jejich hodnot. Je \emph{líné}, protože netrénuje žádný model - veškerá práce je až při predikci.

\textbf{(b)} $k$ se volí \emph{křížovou validací}. \emph{Malé} $k$ (např.\ 1) dává nízké vychýlení, ale vysoký rozptyl (citlivé na šum, přeučení); \emph{velké} $k$ hranici vyhladí (vyšší vychýlení, nižší rozptyl). \emph{Škálování} je nutné, protože vzdálenost jinak ovládne příznak s největším rozsahem.

\textbf{(c)} \emph{Prokletí dimenzionality}: ve vysoké dimenzi se vzdálenosti mezi body vyrovnávají (poměr nejbližší/nejvzdálenější $\to1$), takže pojem ,,soused'' ztrácí smysl a kNN selhává; navíc je třeba exponenciálně víc dat. \emph{Složitost} naivní predikce je $O(nd)$ na jeden dotaz (projít $n$ bodu v dimenzi $d$); urychlení přes k-d stromy či přibližné hledání.
\end{reseni}

\begin{jadro}
kNN = ulož data, pro nový bod najdi $k$ nejbližších; klasifikace = hlasování, regrese = průměr. Líné učení. $k$ přes křížovou validaci (malé $k$ = rozptyl, velké = vychýlení). Škáluj příznaky; trpí prokletím dimenzionality.
\end{jadro}

\kalibrace{kNN (~12\,\%) je vzácnější, ale jednoduchý a vděčný. Definice + role $k$ a škálování = spolehlivá 2-bodová záloha; navíc nese pojem prokletí dimenzionality.}
\past{kNN netrénuje (lazy) - ,,trénink'' je jen uložení dat. Bez škálování dominuje příznak s velkou jednotkou; pozor i na liché $k$ u dvou tříd (remízy).}

\ulohahead{specializace UI-SU}{Statistické testy: t-test a chí-kvadrát}
\begin{zadani}
\begin{enumerate}[label=(\alph*)]
  \item \bod{1} Co je statistický test, nulová hypotéza, hladina významnosti a p-hodnota? Popište jednovýběrový a dvouvýběrový t-test.
  \item \bod{2} Popište chí-kvadrát test dobré shody: testovou statistiku, stupně volnosti a rozhodovací pravidlo.
  \item \bod{3} Co je párový t-test a kdy ho použít? Jak souvisí rozhodnutí testu s chybami 1.\ a 2.\ druhu?
\end{enumerate}
\end{zadani}

\begin{reseni}
\textbf{(a)} \emph{Test} rozhoduje mezi \emph{nulovou hypotézou} $H_0$ a alternativou na základě dat. \emph{Hladina významnosti} $\alpha$ = přípustná pravděpodobnost chyby 1.\ druhu; \emph{p-hodnota} = pravděpodobnost dat aspoň tak extrémních jako pozorovaná, platí-li $H_0$; zamítáme, je-li $p<\alpha$. \emph{Jednovýběrový} t-test: $H_0:\mu=\mu_0$, statistika $t=\frac{\bar x-\mu_0}{s/\sqrt n}$, která má za platnosti $H_0$ Studentovo $t$-rozdělení s $n-1$ stupni volnosti; zamítáme, padne-li $t$ do kritického oboru (resp.\ $p<\alpha$). \emph{Dvouvýběrový}: $H_0:\mu_1=\mu_2$, statistika $t=\frac{\bar x-\bar y}{\sqrt{s_X^2/n+s_Y^2/m}}$ (Welchova varianta při různých rozptylech).

\textbf{(b)} \emph{Chí-kvadrát test dobré shody} porovná pozorované četnosti $O_i$ s očekávanými $E_i$ podle modelu: $\chi^2=\sum_i\frac{(O_i-E_i)^2}{E_i}$. \emph{Stupně volnosti} $=$ (počet kategorií $-1-$ počet odhadovaných parametru). $H_0$ (data odpovídají modelu) zamítneme, je-li $\chi^2$ větší než kritická hodnota $\chi^2_{1-\alpha}$ pro dané df.

\textbf{(c)} \emph{Párový} t-test se použije, jsou-li data \emph{spárovaná} (např.\ měření před/po na týchž jedincích); testuje, zda je průměr rozdílu nulový. \emph{Chyba 1.\ druhu} = zamítnout platné $H_0$ (řízena $\alpha$); \emph{chyba 2.\ druhu} = nezamítnout neplatné $H_0$ (řízena silou testu $1-\beta$, roste s velikostí vzorku a efektu).
\end{reseni}

\begin{jadro}
Test: $H_0$, testová statistika, p-hodnota; zamítni při $p<\alpha$. Jednovýběrový t-test $t=\frac{\bar x-\mu_0}{s/\sqrt n}$. Chí-kvadrát dobré shody $\chi^2=\sum_i\frac{(O_i-E_i)^2}{E_i}$, df $=$ kategorie $-1-$ parametry.
\end{jadro}

\kalibrace{Statistické testy (~22\,\% dohromady) chtějí převážně definice + dosazení do statistiky. Stačí 2-bodová záloha: hypotéza, statistika, rozhodnutí.}
\past{p-hodnota není pravděpodobnost, že $H_0$ platí. Stupně volnosti u chí-kvadrát nejsou počet kategorií, ale po odečtení vazeb/parametru.}

\ulohahead{specializace UI-SU}{Vyhodnocení modelu: validace, přeučení, regularizace}
\begin{zadani}
\begin{enumerate}[label=(\alph*)]
  \item \bod{1} Vysvětlete rozdělení dat na train/dev/test a křížovou validaci ($k$-fold). Co je princip maximální věrohodnosti?
  \item \bod{2} Co je přeučení (overfitting) a generalizační chyba? Jak proti přeučení působí regularizace (L1 vs L2) a včasné zastavení trénování?
  \item \bod{3} Vysvětlete kompromis vychýlení-rozptyl (bias-variance) a prokletí dimenzionality při výběru modelu.
\end{enumerate}
\end{zadani}

\begin{reseni}
\textbf{(a)} \emph{Train} slouží k učení, \emph{dev/validace} k ladění hyperparametru a výběru modelu, \emph{test} k jedinému závěrečnému odhadu výkonu (nesmí se na něm ladit). \emph{$k$-fold křížová validace}: data rozděl na $k$ částí, $k$-krát trénuj na $k{-}1$ a validuj na zbylé, výsledky zprůměruj (lepší využití dat). \emph{Maximální věrohodnost}: zvol parametry maximalizující pravděpodobnost (věrohodnost) pozorovaných dat, $\hat\theta=\arg\max_\theta P(\text{data}\mid\theta)$.

\textbf{(b)} \emph{Přeučení}: model se naučí šum trénovacích dat, má malou trénovací, ale velkou \emph{generalizační} (testovací) chybu. \emph{Regularizace} přidá penaltu za velikost vah: \emph{L2} ($\lambda\lVert w\rVert_2^2$) hladce smršťuje, \emph{L1} ($\lambda\lVert w\rVert_1$) tlačí váhy přesně na nulu (řídké řešení, výběr příznaku). \emph{Včasné zastavení}: trénink se ukončí, když začne růst validační chyba.

\textbf{(c)} \emph{Bias-variance}: příliš jednoduchý model má velké \emph{vychýlení} (podučení), příliš složitý velký \emph{rozptyl} (přeučení); cílem je minimum jejich součtu. \emph{Prokletí dimenzionality}: s rostoucím počtem příznaku roste potřeba dat exponenciálně a vzdálenosti se vyrovnávají, takže složitější modely snáz přeučí - proto redukce dimenze/regularizace.
\end{reseni}

\begin{jadro}
Train (učení) / dev (ladění) / test (jediný závěrečný odhad); $k$-fold CV průměruje. Přeučení = malá train, velká test chyba. L2 hladce smršťuje, L1 dělá řídké váhy (výběr příznaku). Bias-variance: minimalizuj součet.
\end{jadro}

\kalibrace{Toto je ,,zastřešující'' ML téma (evaluace, regularizace) prolínající mnoho otázek; opakovaně se ptá na křížovou validaci a přeučení. Solidní definice = body napříč specializací.}
\past{Na testovacích datech se \emph{neladí}. L1 dělá řídké váhy (výběr příznaku), L2 jen smršťuje - nezaměňuj jejich efekt.}

\ulohahead{specializace UI-SU}{Predikce z tabulárních dat: předzpracování a metriky}
\begin{zadani}
\begin{enumerate}[label=(\alph*)]
  \item \bod{1} Popište předzpracování tabulárních dat: kódování kategoriálních příznaku (one-hot), škálování spojitých příznaku a ošetření chybějících hodnot (imputace).
  \item \bod{2} Definujte regresní metriky MAE, RMSE a $R^2$. Kdy je vhodnější MAE a kdy RMSE?
  \item \bod{3} Jaký model zvolit pro tabulární regresi (lineární / stromové / náhodný les / boosting) a proč stromové metody nepotřebují škálování?
\end{enumerate}
\end{zadani}

\begin{reseni}
\textbf{(a)} \emph{Kategoriální} příznaky se převedou na čísla: \emph{one-hot} kódování (každá hodnota = binární sloupec) pro neuspořádané kategorie; ordinální kódování pro uspořádané. \emph{Spojité} příznaky se \emph{škálují} (standardizace na nulový průměr a jednotkový rozptyl, nebo min-max) - nutné pro vzdálenostní a gradientové modely. \emph{Chybějící hodnoty}: imputace (průměr/medián/nejčastější hodnota), případně příznak ,,bylo chybějící''.

\textbf{(b)} Pro predikce $\hat y_i$ a cíle $y_i$: $\mathrm{MAE}=\tfrac1n\sum|y_i-\hat y_i|$; $\mathrm{RMSE}=\sqrt{\tfrac1n\sum(y_i-\hat y_i)^2}$; $R^2=1-\frac{\sum(y_i-\hat y_i)^2}{\sum(y_i-\bar y)^2}$ (podíl vysvětleného rozptylu). \emph{MAE} je odolnější vuči odlehlým hodnotám; \emph{RMSE} velké chyby trestá kvadraticky (preferuj, když jsou velké chyby zvlášť nežádoucí).

\textbf{(c)} \emph{Lineární regrese} pro zhruba lineární vztahy (rychlá, interpretovatelná); \emph{stromové metody} (rozhodovací strom, \emph{náhodný les}, \emph{gradient boosting}) pro nelineární vztahy a smíšené typy příznaku - na tabulárních datech bývají nejlepší a robustní. Stromy dělí podle \emph{prahu} jednoho příznaku, takže jsou \emph{invariantní k monotónní transformaci} (škálování nezmění pořadí hodnot, tedy ani štěpení) - proto nepotřebují normalizaci.
\end{reseni}

\kalibrace{Tabulární workflow (~10\,\%) se objevil nedávno a chtěl předzpracování + metriky (ne jen definici modelu). MAE/RMSE a one-hot/škálování/imputace = 2-bodová záloha proti ,,end-to-end'' otázce.}
\past{One-hot pro neuspořádané kategorie (jinak modelu podsouváš falešné pořadí). RMSE je v jednotkách cíle (odmocnina), MSE ne; stromy škálování nepotřebují, lineární/kNN/SVM ano.}

\ulohahead{specializace UI-SU \textnormal{\footnotesize[úroveň 2]}}{k-means ručně: dvě iterace a hodnota ztráty}
\begin{zadani}
Body v rovině: $P_1=(1,1)$, $P_2=(2,1)$, $P_3=(4,3)$, $P_4=(5,4)$. Spusť k-means s $k=2$ a počátečními centry $c_1=(1,1)$, $c_2=(5,4)$.
\begin{enumerate}[label=(\alph*)]
  \item \bod{1} Proveďte krok přiřazení (euklidovská vzdálenost).
  \item \bod{2} Přepočítejte centra a proveďte druhou iteraci; rozhodněte, zda algoritmus zkonvergoval.
  \item \bod{3} Spočtěte hodnotu účelové funkce $J=\sum_i\lVert x_i-\mu_{c(i)}\rVert^2$ ve výsledném rozdělení.
\end{enumerate}
\end{zadani}

\begin{reseni}
\textbf{(a)} Vzdálenosti k $c_1=(1,1)$ a $c_2=(5,4)$:
$P_1$: $0$ vs $5$; $P_2$: $1$ vs $\sqrt{18}\approx4{,}24$; $P_3$: $\sqrt{13}\approx3{,}61$ vs $\sqrt2\approx1{,}41$; $P_4$: $5$ vs $0$.
Přiřazení: $\{P_1,P_2\}\to c_1$, $\{P_3,P_4\}\to c_2$.

\textbf{(b)} Nová centra: $c_1=\big(\tfrac{1+2}{2},\tfrac{1+1}{2}\big)=(1{,}5,\,1)$, $c_2=\big(\tfrac{4+5}{2},\tfrac{3+4}{2}\big)=(4{,}5,\,3{,}5)$. Druhá iterace: $P_1,P_2$ jsou jasně blíž $c_1$ (vzdálenosti $0{,}5$) než $c_2$; $P_3,P_4$ blíž $c_2$. Přiřazení se nezměnilo $\Rightarrow$ \emph{konvergence}.

\textbf{(c)} Čtverce vzdáleností k přiřazenému centru:
$c_1$: $\lVert P_1-c_1\rVert^2=0{,}25$, $\lVert P_2-c_1\rVert^2=0{,}25$;
$c_2$: $\lVert P_3-c_2\rVert^2=0{,}25+0{,}25=0{,}5$, $\lVert P_4-c_2\rVert^2=0{,}25+0{,}25=0{,}5$.
$J=0{,}25+0{,}25+0{,}5+0{,}5=1{,}5$.
\end{reseni}

\begin{jadro}
Iterace = (1) přiřaď ke knejbližšímu centru, (2) přepočti centra na průměr přiřazených bodu. $J=$ součet čtvercu vzdáleností k centrum, po každém kroku neroste. Konec, když se přiřazení nezmění.
\end{jadro}

\kalibrace{Numerický k-means je opakovaný typ (shlukování ~44\,\%). Umět dvě iterace ručně = spolehlivé 2-3 body; tato úloha snižuje šanci na skóre $\le1$ u shlukovací otázky.}
\past{Vzdálenosti můžeš porovnávat přes čtverce (není třeba odmocňovat). Centrum je průměr, ne medián; remízy v přiřazení vyřeš pevným pravidlem.}

\ulohahead{specializace UI-SU \textnormal{\footnotesize[úroveň 2]}}{Rozhodovací strom: výpočet informačního zisku}
\begin{zadani}
Uzel obsahuje 6 příkladu, z toho 3 kladné a 3 záporné. Uvažte dva atributy:
$A$ rozdělí data na $A{=}0$: (2 kladné, 1 záporný) a $A{=}1$: (1 kladný, 2 záporné);
$B$ rozdělí data na $B{=}0$: (3 kladné, 0 záporných) a $B{=}1$: (0 kladných, 3 záporné).
\begin{enumerate}[label=(\alph*)]
  \item \bod{1} Spočtěte entropii kořene.
  \item \bod{2} Spočtěte informační zisk atributu $A$ a atributu $B$.
  \item \bod{3} Který atribut zvolíte a proč? Co by udělalo prořezávání?
\end{enumerate}
\end{zadani}

\begin{reseni}
\textbf{(a)} $H(\text{kořen})=-\tfrac12\log_2\tfrac12-\tfrac12\log_2\tfrac12=1$ bit.

\textbf{(b)} Atribut $A$: obě větve mají rozdělení $(2,1)$, resp.\ $(1,2)$, tedy
\[
H=-\tfrac23\log_2\tfrac23-\tfrac13\log_2\tfrac13=\tfrac23\cdot0{,}585+\tfrac13\cdot1{,}585=0{,}918\ \text{bitu}.
\]
Vážená entropie $=\tfrac36\cdot0{,}918+\tfrac36\cdot0{,}918=0{,}918$, takže $\mathrm{IG}(A)=1-0{,}918=0{,}082$.
Atribut $B$: obě větve jsou čisté, $H=0$, vážená entropie $0$, takže $\mathrm{IG}(B)=1-0=1$ bit.

\textbf{(c)} Zvolíme $B$ (maximální informační zisk, dokonalé rozdělení). \emph{Prořezávání}: kdyby další štěpení (např.\ podle $A$ uvnitř větve) zlepšovalo jen trénovací chybu, ne validační, ponecháme list (post-pruning), aby strom negeneralizoval šum.
\end{reseni}

\begin{jadro}
$H=-\sum_k p_k\log_2 p_k$; pro $(2,1)$ vyjde $0{,}918$. $\mathrm{IG}=H(\text{rodič})-\sum_j\frac{|S_j|}{|S|}H(S_j)$. Čisté listy $\Rightarrow H=0\Rightarrow$ maximální zisk. Vyber atribut s největším IG.
\end{jadro}

\kalibrace{Výpočet IG je opakovaný (rozhodovací stromy ~33\,\%). Hodnoty $\log_2\tfrac13=-1{,}585$, $\log_2\tfrac23=-0{,}585$ je vhodné znát; tím se vyhneš skóre $\le1$.}
\past{Informační zisk je \emph{vážený} velikostmi větví. $\log_2$, ne $\ln$. Záporný argument logaritmu nehrozí (jen podíly v $(0,1]$).}

\ulohahead{specializace UI-SU \textnormal{\footnotesize[úroveň 2]}}{Rekonstrukce matice záměn u nevyvážených tříd}
\begin{zadani}
Testovací množina má 1000 vzorku, z toho 5\,\% pozitivních. Klasifikátor má recall $0{,}8$ a precision $0{,}5$.
\begin{enumerate}[label=(\alph*)]
  \item \bod{1} Určete počet pozitivních a negativních a spočtěte $\mathrm{TP}$ a $\mathrm{FN}$.
  \item \bod{2} Z precision dopočítejte $\mathrm{FP}$ a poté $\mathrm{TN}$. Sestavte matici záměn.
  \item \bod{3} Spočtěte accuracy a $F_1$ a vysvětlete, proč je zde accuracy zavádějící.
\end{enumerate}
\end{zadani}

\begin{reseni}
\textbf{(a)} Pozitivních $P=0{,}05\cdot1000=50$, negativních $N=950$. $\mathrm{TP}=\text{recall}\cdot P=0{,}8\cdot50=40$, $\mathrm{FN}=P-\mathrm{TP}=10$.

\textbf{(b)} Z $\text{precision}=\frac{\mathrm{TP}}{\mathrm{TP}+\mathrm{FP}}=0{,}5$ plyne $\mathrm{TP}+\mathrm{FP}=80$, tedy $\mathrm{FP}=40$. $\mathrm{TN}=N-\mathrm{FP}=950-40=910$. Matice záměn:
\[
\begin{array}{c|cc}
 & \hat{+} & \hat{-}\\\hline
+ & 40 & 10\\
- & 40 & 910
\end{array}
\]

\textbf{(c)} $\text{accuracy}=\frac{\mathrm{TP}+\mathrm{TN}}{1000}=\frac{950}{1000}=0{,}95$. $F_1=\frac{2\cdot0{,}5\cdot0{,}8}{0{,}5+0{,}8}=\frac{0{,}8}{1{,}3}\approx0{,}615$. Accuracy $0{,}95$ vypadá skvěle, ale triviální klasifikátor ,,vše negativní'' by měl accuracy $0{,}95$ taky (950/1000) a nenašel by nic - proto se u nevyvážených tříd díváme na precision/recall/$F_1$.
\end{reseni}

\begin{jadro}
$\mathrm{TP}=\text{recall}\cdot P$, $\mathrm{FN}=P-\mathrm{TP}$; z precision $\mathrm{TP}+\mathrm{FP}=\mathrm{TP}/\text{precision}$, $\mathrm{TN}=N-\mathrm{FP}$. Accuracy u nevyvážených tříd klame (porovnej s ,,vše majoritní'').
\end{jadro}

\kalibrace{Tento přesný typ (rekonstrukce počtu z metrik) padl opakovaně (2023-léto, 2025-podzim). Nacvičit = jistá 2-3 z evaluační otázky.}
\past{Jmenovatel precision jsou \emph{predikované} pozitivní ($\mathrm{TP}+\mathrm{FP}$), recall \emph{skutečné} pozitivní ($\mathrm{TP}+\mathrm{FN}$). Nepleť $P$ (skutečné pozitivní) s predikcemi.}

\ulohahead{specializace UI-SU \textnormal{\footnotesize[úroveň 2]}}{Lineární regrese ručně a vliv L2 regularizace}
\begin{zadani}
Data jednorozměrné regrese: $(x,y)\in\{(1,1),(2,2),(3,2)\}$, model $\hat y=w_0+w_1x$.
\begin{enumerate}[label=(\alph*)]
  \item \bod{1} Napište normální rovnice v maticovém tvaru.
  \item \bod{2} Spočtěte řešení nejmenších čtverců $w_0,w_1$.
  \item \bod{3} Jak by se řešení změnilo s L2 regularizací (ridge) a proč? Uveďte upravené normální rovnice.
\end{enumerate}
\end{zadani}

\begin{reseni}
\textbf{(a)} S maticí $X=\begin{psmallmatrix}1&1\\1&2\\1&3\end{psmallmatrix}$ a $y=(1,2,2)^{\!\top}$ jsou normální rovnice $X^{\!\top}X\,w=X^{\!\top}y$, kde
$X^{\!\top}X=\begin{psmallmatrix}3&6\\6&14\end{psmallmatrix}$, $X^{\!\top}y=\begin{psmallmatrix}5\\11\end{psmallmatrix}$.

\textbf{(b)} Snazší přes vzorce jednoduché regrese: $\bar x=2$, $\bar y=\tfrac53$.
$w_1=\dfrac{\sum(x-\bar x)(y-\bar y)}{\sum(x-\bar x)^2}=\dfrac{(-1)(-\tfrac23)+0+(1)(\tfrac13)}{1+0+1}=\dfrac{1}{2}=0{,}5$.
$w_0=\bar y-w_1\bar x=\tfrac53-0{,}5\cdot2=\tfrac23\approx0{,}667$. Tedy $\hat y=0{,}667+0{,}5\,x$.
(Kontrola maticově: $\det(X^{\!\top}X)=3\cdot14-36=6$, $w=\tfrac16\begin{psmallmatrix}14&-6\\-6&3\end{psmallmatrix}\begin{psmallmatrix}5\\11\end{psmallmatrix}=\tfrac16\begin{psmallmatrix}4\\3\end{psmallmatrix}=\begin{psmallmatrix}0{,}667\\0{,}5\end{psmallmatrix}$.)

\textbf{(c)} Ridge minimalizuje $\lVert Xw-y\rVert^2+\lambda\lVert w\rVert^2$, normální rovnice $(X^{\!\top}X+\lambda I)\,w=X^{\!\top}y$. S rostoucím $\lambda$ se $w_1$ (a $w_0$, neregularizuje-li se zvlášť) \emph{smršťuje k nule}: matice $X^{\!\top}X+\lambda I$ je vždy regulární a dobře podmíněná, takže řešení existuje i při kolineárních datech a model méně přeučuje (menší rozptyl).
\end{reseni}

\begin{jadro}
$w=(X^{\!\top}X)^{-1}X^{\!\top}y$; jednodušeji $w_1=\frac{\sum(x-\bar x)(y-\bar y)}{\sum(x-\bar x)^2}$, $w_0=\bar y-w_1\bar x$. Ridge: $(X^{\!\top}X+\lambda I)w=X^{\!\top}y$, smršťuje koeficienty.
\end{jadro}

\kalibrace{Ruční OLS / jeden ridge krok je opakovaný (regrese ~67\,\%). Tato úloha sytí kritickou podlahu specializace a snižuje šanci na skóre $\le1$.}
\past{$X$ musí mít sloupec jedniček pro bias. Vzorec pro $w_1$ má ve jmenovateli rozptyl $x$, ne kovarianci.}

\ulohahead{specializace UI-SU \textnormal{\footnotesize[úroveň 2]}}{Logistická regrese: jeden krok SGD}
\begin{zadani}
Binární logistická regrese s vahami $w=(w_0,w_1)=(0,0)$ (s biasem). Trénovací příklad $x=(1,2)$ (první složka je bias), skutečná třída $y=1$, rychlost učení $\eta=0{,}1$.
\begin{enumerate}[label=(\alph*)]
  \item \bod{1} Spočtěte predikci $p=\sigma(w^{\!\top}x)$.
  \item \bod{2} Spočtěte gradient ztráty na tomto příkladu a proveďte jeden krok SGD.
  \item \bod{3} Spočtěte novou predikci na stejném bodě a ověřte, že se posunula správným směrem.
\end{enumerate}
\end{zadani}

\begin{reseni}
\textbf{(a)} $w^{\!\top}x=0\cdot1+0\cdot2=0$, takže $p=\sigma(0)=\tfrac{1}{1+e^{0}}=0{,}5$.

\textbf{(b)} Gradient NLL na jednom příkladu: $\nabla_w=(p-y)\,x=(0{,}5-1)\,(1,2)=(-0{,}5,\,-1)$. Krok SGD:
\[
w\leftarrow w-\eta\,\nabla_w=(0,0)-0{,}1\cdot(-0{,}5,-1)=(0{,}05,\,0{,}1).
\]

\textbf{(c)} Nová predikce: $w^{\!\top}x=0{,}05\cdot1+0{,}1\cdot2=0{,}25$, $p=\sigma(0{,}25)=\tfrac{1}{1+e^{-0{,}25}}\approx\tfrac{1}{1+0{,}779}\approx0{,}562$. Pravděpodobnost třídy 1 vzrostla z $0{,}5$ na $0{,}562$, tedy směrem ke skutečné třídě $y=1$. Správně.
\end{reseni}

\begin{jadro}
$p=\sigma(w^{\!\top}x)=\frac{1}{1+e^{-w^{\!\top}x}}$. Gradient $(p-y)x$. SGD: $w\leftarrow w-\eta(p-y)x$. Při $y=1$ a $p<1$ je $(p-y)<0$, takže váhy ve směru $x$ rostou a $p$ se zvyšuje.
\end{jadro}

\kalibrace{Jeden SGD krok logistické regrese je opakovaný (logistická regrese se zkoušela 2024-podzim, 2025-jaro). Sigmoid + gradient $(p-y)x$ zpaměti = jistá 2-3.}
\past{Gradient je $(p-y)x$, ne $(y-p)x$ - pak by byl krok $w\leftarrow w+\eta(y-p)x$ (znaménko musí sedět se směrem minimalizace). $\sigma(0)=0{,}5$.}

\ulohahead{specializace UI-SU \textnormal{\footnotesize[úroveň 2]}}{PCA: kovariance, vlastní rozklad, projekce}
\begin{zadani}
Centrovaná data mají kovarianční matici $C=\begin{psmallmatrix}2&1\\1&2\end{psmallmatrix}$.
\begin{enumerate}[label=(\alph*)]
  \item \bod{1} Spočtěte vlastní čísla $C$ a první hlavní komponentu (vlastní vektor).
  \item \bod{2} Jaký podíl rozptylu vysvětluje první komponenta?
  \item \bod{3} Na kterou osu data projektujete při redukci na 1 dimenzi a jak projekci spočtete?
\end{enumerate}
\end{zadani}

\begin{reseni}
\textbf{(a)} $\det(C-\lambda I)=(2-\lambda)^2-1=\lambda^2-4\lambda+3=(\lambda-1)(\lambda-3)$, tedy $\lambda_1=3$, $\lambda_2=1$. K $\lambda_1=3$: $(C-3I)=\begin{psmallmatrix}-1&1\\1&-1\end{psmallmatrix}$, vlastní vektor $v_1=(1,1)^{\!\top}$, normovaně $\tfrac{1}{\sqrt2}(1,1)$. To je první hlavní komponenta.

\textbf{(b)} Vysvětlený rozptyl první komponenty $=\dfrac{\lambda_1}{\lambda_1+\lambda_2}=\dfrac{3}{4}=0{,}75$, tedy 75\,\%.

\textbf{(c)} Projektujeme na směr $v_1=\tfrac1{\sqrt2}(1,1)$ (největší vlastní číslo). Nová (1D) souřadnice bodu $x$ je skalární součin $z=v_1^{\!\top}x=\tfrac{1}{\sqrt2}(x_1+x_2)$. Tím zachováme maximum rozptylu (75\,\%) a zahodíme směr $(1,-1)$ s rozptylem 25\,\%.
\end{reseni}

\begin{jadro}
Hlavní komponenty = vlastní vektory kovariance, seřazené podle vlastních čísel (= rozptylu podél nich). Vysvětlený rozptyl $=\lambda_i/\sum\lambda_j$. Projekce na PC: $z=v^{\!\top}x$. Před PCA centruj (a obvykle škáluj).
\end{jadro}

\kalibrace{Numerická PCA na 2D (kovariance -> vlastní čísla -> projekce) je vděčný typ (PCA ~12\,\%, ale opakovaný jako vizuální/výpočetní). Sytí specializační podlahu.}
\past{Komponenta s \emph{největším} vlastním číslem nese nejvíc rozptylu. Vlastní vektory normuj. PCA hledá rozptyl, ne oddělení tříd.}

\ulohahead{specializace UI-SU \textnormal{\footnotesize[úroveň 2]}}{SVM: geometrie hranice a vliv bodu}
\begin{zadani}
Dvě lineárně separabilní třídy v rovině: kladné body $(3,3),(4,4)$, záporné body $(0,0),(1,1)$.
\begin{enumerate}[label=(\alph*)]
  \item \bod{1} Popište, kudy povede rozdělující nadrovina s maximálním marginem a které body jsou podpůrné vektory. Načrtněte (slovně) hranici.
  \item \bod{2} Co se stane s hranicí, přidáme-li bod $(2,2)$ s kladnou třídou? A co bod $(5,5)$ s kladnou třídou daleko za hranicí?
  \item \bod{3} Kdy lineární SVM selže a jak pomůže RBF jádro?
\end{enumerate}
\end{zadani}

\begin{reseni}
\textbf{(a)} Body leží na přímce $y=x$; kladné jsou ,,nahoře vpravo'', záporné ,,dole vlevo''. Maximálně marginová nadrovina je kolmá na spojnici nejbližších bodu protilehlých tříd, tedy přímka $x+y=c$ procházející středem mezi $(1,1)$ a $(3,3)$, tj.\ $x+y=4$ (kolmá na směr $(1,1)$). \emph{Podpůrné vektory} jsou nejbližší body $(1,1)$ a $(3,3)$; body $(0,0)$ a $(4,4)$ jsou dál a hranici neurčují.

\textbf{(b)} Bod $(2,2)$ leží přesně na hranici $x+y=4$; jako \emph{kladný} tedy leží uvnitř marginu (porušuje ho), stane se podpůrným vektorem a hranici \emph{posune/zúží} (řešení se změní). Bod $(5,5)$: $x+y=10$, leží daleko v kladné oblasti za marginem, je správně klasifikovaný a není mezi nejbližšími - hranici proto \emph{nezmění}. Obecně: body daleko a správně klasifikované (mimo margin) řešení neovlivní; mění ho jen body v marginu, resp.\ nejbližší body.

\textbf{(c)} Lineární SVM selže, nejsou-li třídy lineárně separabilní (např.\ XOR nebo soustředné kružnice). \emph{RBF jádro} $K(x,z)=e^{-\gamma\lVert x-z\rVert^2}$ implicitně mapuje do vysokorozměrného prostoru, kde hranice může být nelineární (lokální, podle vzdálenosti k podpůrným vektorum); malé $\gamma$ = hladká hranice, velké $\gamma$ = klikatá (přeučení).
\end{reseni}

\begin{jadro}
Maximálně marginová hranice je kolmá na spojnici nejbližších protilehlých bodu, vede jejich středem; ty body jsou podpůrné vektory. Body daleko a správně neovlivní řešení. Neseparabilní data -> měkký margin ($C$) nebo nelineární jádro (RBF).
\end{jadro}

\kalibrace{Geometrická SVM otázka (,,nakreslete hranici, vliv přidaného bodu, kdy RBF'') padla 2024-léto skoro doslova. Konceptuálně-geometrické zvládnutí = 2-3 body.}
\past{Hranici určují jen podpůrné vektory; přidání bodu daleko za marginem nic nezmění, přidání do marginu ano. Velké $\gamma$/velké $C$ = přeučení.}

\ulohahead{specializace UI-SU \textnormal{\footnotesize[úroveň 2]}}{ROC křivka a AUC z konkrétních skóre}
\begin{zadani}
Klasifikátor přiřadil skóre: pozitivní příklady $\{0{,}9,\,0{,}6,\,0{,}4\}$, negativní příklady $\{0{,}7,\,0{,}3,\,0{,}2\}$ (tedy $P=3$, $N=3$).
\begin{enumerate}[label=(\alph*)]
  \item \bod{1} Pro klesající práh spočtěte body ROC (TPR, FPR).
  \item \bod{2} Načrtněte (popište) tvar ROC křivky.
  \item \bod{3} Spočtěte AUC a interpretujte ji.
\end{enumerate}
\end{zadani}

\begin{reseni}
\textbf{(a)} Setříděno sestupně: $0{,}9(+),0{,}7(-),0{,}6(+),0{,}4(+),0{,}3(-),0{,}2(-)$. Predikuj ,,+'' při skóre $\ge$ práh; $\mathrm{TPR}=\mathrm{TP}/3$, $\mathrm{FPR}=\mathrm{FP}/3$:
\[
\begin{array}{c|cc|cc}
\text{práh} & \mathrm{TP} & \mathrm{FP} & \mathrm{TPR} & \mathrm{FPR}\\\hline
0{,}9 & 1 & 0 & 0{,}33 & 0\\
0{,}7 & 1 & 1 & 0{,}33 & 0{,}33\\
0{,}6 & 2 & 1 & 0{,}67 & 0{,}33\\
0{,}4 & 3 & 1 & 1{,}00 & 0{,}33\\
0{,}3 & 3 & 2 & 1{,}00 & 0{,}67\\
0{,}2 & 3 & 3 & 1{,}00 & 1{,}00
\end{array}
\]

\textbf{(b)} Křivka jde z $(0,0)$ schody nahoru a doprava do $(1,1)$; nejprve stoupá strmě (zachytí 2 pozitivní při FPR $0{,}33$), pak doroste TPR na 1 a nakonec dobírá falešně pozitivní. Leží nad diagonálou (lepší než náhoda).

\textbf{(c)} AUC $=$ pravděpodobnost, že náhodný pozitivní dostane vyšší skóre než náhodný negativní. Spočítáme přes dvojice: $0{,}9>$ všechny 3 negativní; $0{,}6>\{0{,}3,0{,}2\}$ (2); $0{,}4>\{0{,}3,0{,}2\}$ (2). Celkem $3+2+2=7$ z $9$ dvojic. $\text{AUC}=\tfrac{7}{9}\approx0{,}78$. Tedy klasifikátor řadí pozitivní nad negativní v 78\,\% případu (1 = ideál, 0{,}5 = náhoda).
\end{reseni}

\begin{jadro}
ROC = TPR vs FPR při klesajícím prahu; každý pozitivní zvýší TPR, každý negativní FPR. AUC $=$ podíl dvojic (pozitivní, negativní), kde pozitivní má vyšší skóre $=$ P(skóre$^+>$ skóre$^-$).
\end{jadro}

\kalibrace{ROC/AUC z konkrétních skóre je legální numerický typ evaluační otázky (~44\,\%). Doplňuje rekonstrukci matice záměn a snižuje šanci na skóre $\le1$.}
\past{ROC vynáší TPR proti \emph{FPR}, ne proti precision (to je PR křivka). AUC přes počítání dvojic je rychlejší než integrace.}

\ulohahead{specializace UI-SU \textnormal{\footnotesize[úroveň 3]}}{Klasifikace do více tříd (softmax) a křížová validace}
\begin{zadani}
\begin{enumerate}[label=(\alph*)]
  \item \bod{1} Definujte softmax pro klasifikaci do $K$ tříd a uveďte alternativu one-vs-rest.
  \item \bod{2} Napište ztrátovou funkci (kategoriální křížová entropie) a gradient pro jeden krok.
  \item \bod{3} Popište $k$-násobnou křížovou validaci a spočtěte průměrnou přesnost z foldu $\{0{,}80;0{,}85;0{,}75;0{,}90;0{,}80\}$.
\end{enumerate}
\end{zadani}

\begin{reseni}
\textbf{(a)} \emph{Softmax} převede skóre $z_k=w_k^{\!\top}x$ na pravděpodobnosti tříd:
\[
p_k=\frac{e^{z_k}}{\sum_{j=1}^K e^{z_j}},\qquad \sum_k p_k=1.
\]
Predikuje se třída s nejvyšším $p_k$. \emph{One-vs-rest}: natrénuj $K$ binárních klasifikátoru (každý ,,třída $k$ vs zbytek'') a vyber ten s nejvyšším skóre - jednodušší, ale skóre nejsou kalibrovaná pravděpodobnost.

\textbf{(b)} \emph{Kategoriální křížová entropie} pro správnou třídu $y$: $L=-\ln p_y=-\ln\frac{e^{z_y}}{\sum_j e^{z_j}}$. Gradient vuči skóre má stejně elegantní tvar jako u logistické regrese:
\[
\frac{\partial L}{\partial z_k}=p_k-\mathbb{1}[k=y],\qquad \nabla_{w_k}L=(p_k-\mathbb{1}[k=y])\,x.
\]
Krok SGD: $w_k\leftarrow w_k-\eta(p_k-\mathbb{1}[k=y])x$.

\textbf{(c)} \emph{$k$-fold CV}: rozděl data na $k$ částí, $k$-krát trénuj na $k-1$ z nich a vyhodnoť na zbylé; výsledky zprůměruj (robustnější odhad než jeden split, využije všechna data). Použije se i k volbě hyperparametru (vyber hodnotu s nejlepší průměrnou CV přesností). Průměr: $\frac{0{,}80+0{,}85+0{,}75+0{,}90+0{,}80}{5}=\frac{4{,}10}{5}=0{,}82$.
\end{reseni}

\begin{jadro}
Softmax $p_k=e^{z_k}/\sum_j e^{z_j}$; ztráta $-\ln p_y$; gradient $(p_k-\mathbb{1}[k{=}y])x$. $k$-fold CV: průměr přesností přes foldy; vyber hyperparametr s nejlepším průměrem.
\end{jadro}

\kalibrace{Klasifikace do více tříd (softmax) je v požadavku logistické regrese; CV je opakované evaluační téma. Úroveň 3 doplňuje multiclass variantu.}
\past{Softmax normuje přes \emph{všechny} třídy (jmenovatel je suma). CV se průměruje přes foldy; na testu se nevybírá hyperparametr (jen na dev/CV).}


\section{Specializace: Základy umělé inteligence (Téma 1)}
\ulohahead{specializace UI-SU}{Splňování podmínek (CSP), hranová konzistence, AC-3}
\begin{zadani}
\begin{enumerate}[label=(\alph*)]
  \item \bod{1} Definujte problém splňování podmínek (CSP) a pojem \emph{hranové konzistence} (arc consistency).
  \item \bod{2} Popište algoritmus AC-3 (fronta hran, revize domény, šíření) a uveďte jeho časovou složitost.
  \item \bod{3} Pro proměnné $X,Y,Z$ s doménami $\{1,2,3\}$ a omezeními $X<Y$, $Y<Z$ proveďte AC-3 a uveďte výsledné domény. Vysvětlete rozdíl mezi lokální (hranovou) a globální konzistencí a proč hranová konzistence nezaručí existenci řešení.
\end{enumerate}
\end{zadani}

\begin{reseni}
\textbf{(a) Definice.} CSP je trojice $(\mathcal{V},\mathcal{D},\mathcal{C})$: proměnné $\mathcal{V}=\{X_1,\dots,X_n\}$, jejich domény $\mathcal{D}=\{D_1,\dots,D_n\}$ a omezení $\mathcal{C}$ (každé omezuje přípustné kombinace hodnot podmnožiny proměnných). Řešení = přiřazení hodnot z domén splňující všechna omezení. Hrana $(X_i,X_j)$ (binární omezení) je \emph{hranově konzistentní}, pokud pro každou hodnotu $a\in D_i$ existuje hodnota $b\in D_j$ taková, že dvojice $(a,b)$ splňuje omezení. CSP je hranově konzistentní, jsou-li konzistentní všechny hrany.

\textbf{(b) AC-3.} Udržuj frontu všech hran (orientovaně, obě strany každého omezení). Opakovaně vyjmi hranu $(X_i,X_j)$ a proveď \textsc{Revise}: odstraň z $D_i$ každou hodnotu, pro kterou v $D_j$ neexistuje podpora. Pokud se $D_i$ zmenšila, vlož zpět do fronty všechny hrany $(X_k,X_i)$ pro sousedy $X_k$ (mohla se porušit jejich konzistence). Konec, když je fronta prázdná, nebo nějaká doména zůstane prázdná (pak CSP nemá řešení). Složitost: $O(e\,d^3)$, kde $e$ je počet hran (binárních omezení) a $d$ velikost domény (každá z $O(e)$ hran se zařadí nejvýše $d$-krát, jedna \textsc{Revise} je $O(d^2)$).

\textbf{(c) Průběh.} Z $X<Y$: $X$ nemůže být $3$ (nic většího v $Y$) $\Rightarrow D_X=\{1,2\}$; $Y$ nemůže být $1$ $\Rightarrow D_Y=\{2,3\}$. Z $Y<Z$: $Y$ nemůže být $3$ $\Rightarrow D_Y=\{2\}$; $Z$ nemůže být $1$ ani $2$ $\Rightarrow D_Z=\{3\}$. Zmenšení $D_Y$ vrátí do fronty hranu $(X,Y)$: $X$ nyní potřebuje $y>x$ s $y\in\{2\}$, takže $x=2$ vypadne $\Rightarrow D_X=\{1\}$. Výsledek: $D_X=\{1\},\,D_Y=\{2\},\,D_Z=\{3\}$ (zde jediné řešení).

\emph{Lokální vs.\ globální.} Hranová konzistence je \emph{lokální} vlastnost (po dvojicích). Nezaručuje globální řešení: např.\ tři proměnné s doménami $\{1,2\}$ a omezeními ,,navzájem různé'' jsou hranově konzistentní (každá hodnota má podporu u souseda), ale tři po dvou různé hodnoty z dvouprvkové domény neexistují, takže CSP je nesplnitelný. Proto se AC-3 kombinuje s prohledáváním (backtracking), typicky algoritmus MAC = po každém přiřazení znovu udělat hranovou konzistenci.
\end{reseni}

\begin{jadro}
CSP $=(V,D,C)$. Hranová konzistence: každá hodnota $a\in D_i$ má podporu $b\in D_j$ splňující omezení. AC-3 (fronta hran, revize, šíření při zúžení), $O(e\,d^3)$. Nezaručuje řešení -> nutný backtracking/MAC.
\end{jadro}

\kalibrace{Specializaci tvoří ~3 ML + 1 ,,Základy UI''. AI okruhy (CSP, minimax, Bayesovské sítě, A*, MDP, HMM) jsou jednotlivě vzácné, ale slot na jednu z nich je skoro vždy. Proto v úrovni 1 musí mít každá AI rodina aspoň 2-bodovou zálohu: přesná definice + jeden kanonický postup.}
\past{Hranová konzistence \emph{nedokazuje} řešitelnost - to je oblíbená chytací podotázka. Vždy zmiň protipříklad (např.\ all-different na malé doméně) a roli backtrackingu/MAC.}

\ulohahead{specializace UI-SU}{Prohledávání: neinformované a informované (A*)}
\begin{zadani}
\begin{enumerate}[label=(\alph*)]
  \item \bod{1} Popište formulaci úlohy prohledáváním (stav, operátory, cíl, cena). Rozdíl mezi stromovým a grafovým prohledáváním. Uveďte neinformované metody BFS, DFS a uniform-cost.
  \item \bod{2} Popište algoritmus A* a funkci $f(n)=g(n)+h(n)$. Co je přípustná a co konzistentní heuristika?
  \item \bod{3} Za jakých podmínek je A* optimální? Co znamená dominance heuristik a proč je lepší ,,větší'' přípustná heuristika?
\end{enumerate}
\end{zadani}

\begin{reseni}
\textbf{(a)} \emph{Formulace}: počáteční stav, množina \emph{operátoru} (akce s cenou) generujících následníky, \emph{cílový test}, cena cesty. \emph{Stromové} prohledávání neeviduje navštívené stavy (může opakovat), \emph{grafové} drží uzavřenou množinu (closed set). \emph{BFS}: po vrstvách (fronta), úplné, optimální při jednotkových cenách. \emph{DFS}: do hloubky (zásobník), paměťově levné, neúplné/nesuboptimální. \emph{Uniform-cost}: rozšiřuje uzel s nejmenší $g$ (Dijkstra), optimální pro nezáporné ceny.

\textbf{(b)} \emph{A*} rozšiřuje uzel s nejmenším $f(n)=g(n)+h(n)$, kde $g$ je cena z počátku a $h$ \emph{heuristický} odhad ceny do cíle. Heuristika je \emph{přípustná}, pokud nikdy nepřeceňuje ($0\le h(n)\le h^*(n)$); \emph{konzistentní} (monotónní), pokud $h(n)\le c(n,n')+h(n')$ pro každý přechod (trojúhelníková nerovnost).

\textbf{(c)} A* je optimální se \emph{přípustnou} heuristikou ve stromovém prohledávání a s \emph{konzistentní} heuristikou v grafovém (jinak je nutné znovuotevírání uzlu). \emph{Dominance}: $h_2$ dominuje $h_1$, je-li $h_2(n)\ge h_1(n)$ pro všechny $n$ (a obě přípustné); pak A* s $h_2$ rozšíří méně uzlu (efektivnější), protože vyšší přípustný odhad ořeže víc.
\end{reseni}

\kalibrace{Prohledávání (~15\,\%) je opakovaná ,,Základy UI'' otázka (A* se objevil). Definice A* + přípustnost/konzistence = 2-bodová záloha pro AI slot specializace.}
\past{Přípustná $\neq$ konzistentní (konzistence je silnější a implikuje přípustnost). V grafovém A* je pro optimalitu potřeba konzistence, ne jen přípustnost.}

\ulohahead{specializace UI-SU}{Hry: minimax, alfa-beta a teorie her}
\begin{zadani}
\begin{enumerate}[label=(\alph*)]
  \item \bod{1} Popište algoritmus Minimax pro hru dvou hráču s nulovým součtem. Co předpokládá o protihráči?
  \item \bod{2} Popište alfa-beta prořezávání: co a proč prořezává a jak pořadí tahu ovlivní úsporu.
  \item \bod{3} Vysvětlete základy teorie her: vězňovo dilema, dominantní strategie, Nashovo ekvilibrium a Paretovskou optimalitu. Co je mechanism design a jaké jsou typy aukcí?
\end{enumerate}
\end{zadani}

\begin{reseni}
\textbf{(a)} \emph{Minimax}: prohledá herní strom; hráč MAX volí tah maximalizující hodnotu, MIN minimalizující. Hodnota listu = užitek; hodnota vnitřního uzlu = max/min hodnot synu podle toho, kdo je na tahu. Předpokládá \emph{optimálně hrajícího} protihráče (nejhorší případ pro nás). Hloubku lze omezit a listy nahradit \emph{ohodnocovací funkcí}.

\textbf{(b)} \emph{Alfa-beta} udržuje $\alpha$ (nejlepší jistá hodnota pro MAX) a $\beta$ (pro MIN). Větev se \emph{ořízne}, jakmile $\alpha\ge\beta$, protože ji racionální hráč stejně nezvolí, takže nemá smysl ji dál prohledávat. Výsledek je identický s minimaxem. Při \emph{dobrém pořadí} tahu (nejlepší první) klesne počet prohledaných uzlu z $O(b^d)$ až na $O(b^{d/2})$, tj.\ dvojnásobná dosažitelná hloubka.

\textbf{(c)} \emph{Vězňovo dilema}: dva hráči, každý volí spolupráci/zradu; zrada je \emph{dominantní strategie} (lepší bez ohledu na soupeře), takže oba zradí, ač by oboustranná spolupráce byla lepší. \emph{Nashovo ekvilibrium}: profil strategií, kde nikdo nezíská jednostrannou změnou (vzájemně nejlepší odpovědi); ve vězňově dilematu je to (zrada, zrada). \emph{Paretovsky optimální} je výsledek, který nelze zlepšit jednomu bez zhoršení druhému - (spolupráce, spolupráce) je Pareto-lepší, ale není Nashovo ekvilibrium. \emph{Mechanism design} je ,,obrácená'' teorie her: navrhuje pravidla hry tak, aby racionální hráči svým chováním dosáhli žádoucího výsledku. Příklad jsou \emph{aukce}: anglická (rostoucí, otevřená), holandská (klesající), prvocenová obálková a \emph{Vickreyova} (druhá cena) - u Vickreyovy je dominantní strategií dražit svou pravou hodnotu (pravdomluvnost).
\end{reseni}

\begin{jadro}
Minimax: MAX maximalizuje, MIN minimalizuje hodnoty listu. Alfa-beta ořeže větev při $\alpha\ge\beta$ (až $O(b^{d/2})$). Nashovo ekvilibrium = profil, kde nikdo nezíská jednostrannou změnou (vězňovo dilema: (zrada, zrada)).
\end{jadro}

\kalibrace{Hry (~22\,\%, minimax + teorie her) jsou opakovaná ,,Základy UI'' otázka. Minimax + alfa-beta + Nash = 2-3 body za vysvětlení.}
\past{Alfa-beta nemění \emph{výsledek}, jen rychlost. Nashovo ekvilibrium nemusí být Paretovsky optimální (vězňovo dilema je důkaz).}

\ulohahead{specializace UI-SU}{Bayesovské sítě a pravděpodobnostní inference}
\begin{zadani}
\begin{enumerate}[label=(\alph*)]
  \item \bod{1} Definujte Bayesovskou síť a její vztah k úplné sdružené distribuci (faktorizace).
  \item \bod{2} Jak se síť konstruuje (podmíněná nezávislost) a jak se počítá exaktní inference enumerací?
  \item \bod{3} Popište eliminaci proměnných a aproximační inferenci (Monte Carlo: zamítání, vážení věrohodností).
\end{enumerate}
\end{zadani}

\begin{reseni}
\textbf{(a)} \emph{Bayesovská síť} = orientovaný acyklický graf, kde uzly jsou náhodné proměnné a hrany vyjadřují přímou závislost; každý uzel má tabulku podmíněných pravděpodobností $P(X_i\mid \mathrm{rodiče}(X_i))$. Kóduje \emph{úplnou sdruženou distribuci} faktorizací
\[
P(X_1,\dots,X_n)=\prod_{i=1}^n P\big(X_i\mid \mathrm{rodiče}(X_i)\big),
\]
což šetří parametry oproti plné tabulce (využívá podmíněné nezávislosti).

\textbf{(b)} \emph{Konstrukce}: seřaď proměnné (ideálně příčiny před následky) a pro každou $X_i$ vyber za rodiče $\mathrm{Pa}(X_i)$ takovou podmnožinu předchozích proměnných, aby byla $X_i$ podmíněně nezávislá na ostatních předchozích proměnných při daných $\mathrm{Pa}(X_i)$. \emph{Inference enumerací} dotazu $P(Q\mid e)$: $P(Q\mid e)\propto\sum_{\text{skryté } h}P(Q,e,h)$, kde sčítáme součiny faktoru z faktorizace přes všechny hodnoty skrytých proměnných a nakonec normujeme.

\textbf{(c)} \emph{Eliminace proměnných} sčítá zprava efektivněji: postupně ,,vysčítá'' skryté proměnné, mezivýsledky drží jako faktory a opakované součiny nepřepočítává - výrazně rychlejší než holá enumerace. \emph{Aproximace (Monte Carlo)}: \emph{zamítavé vzorkování} (rejection) generuje vzorky z priorů a zahodí ty neslučitelné s $e$ (neefektivní při vzácném $e$); \emph{vážení věrohodností} (likelihood weighting) fixuje pozorované proměnné a každý vzorek váží součinem jejich pravděpodobností - žádný vzorek nezahazuje.
\end{reseni}

\begin{jadro}
Bayesovská síť = DAG + tabulky $P(X_i\mid\text{rodiče}(X_i))$, kóduje $P(X_1,\dots,X_n)=\prod_i P(X_i\mid\text{rodiče}(X_i))$. Inference: $P(Q\mid e)\propto\sum_{\text{skryté }h}P(Q,e,h)$ (enumerace, efektivněji eliminace proměnných).
\end{jadro}

\kalibrace{Bayesovské sítě (~22\,\% s Bayesovským učením) jsou opakovaná ,,Základy UI'' otázka. Faktorizace + enumerace = 2-bodová záloha; tady je AI rodina, kde je vhodné mít i hlubší zvládnutí.}
\past{Faktorizace je součin $P(X_i\mid\text{rodiče})$ - ne marginálu. Likelihood weighting nezahazuje vzorky (na rozdíl od rejection), proto je efektivnější při vzácné evidenci.}

\ulohahead{specializace UI-SU}{Bayesovské učení a EM algoritmus}
\begin{zadani}
\begin{enumerate}[label=(\alph*)]
  \item \bod{1} Popište Bayesovské učení: prior, věrohodnost, posterior. Co je maximálně věrohodná (ML) a maximálně aposteriorní (MAP) hypotéza a co je Bayesovsky optimální predikce?
  \item \bod{2} Popište EM algoritmus (E-krok a M-krok) pro model se skrytou proměnnou.
  \item \bod{3} Napište vzorec pro responsibility (odpovědnost) v E-kroku a aktualizaci parametru v M-kroku pro skrytou kategoriální proměnnou (shluk).
\end{enumerate}
\end{zadani}

\begin{reseni}
\textbf{(a)} Bayesovo pravidlo: $P(h\mid D)=\frac{P(D\mid h)P(h)}{P(D)}$, kde $P(h)$ je \emph{prior}, $P(D\mid h)$ \emph{věrohodnost}, $P(h\mid D)$ \emph{posterior}. \emph{ML hypotéza} $h_{ML}=\arg\max_h P(D\mid h)$; \emph{MAP} $h_{MAP}=\arg\max_h P(D\mid h)P(h)$ (= ML při uniformním prioru). \emph{Bayesovsky optimální predikce} nepoužívá jedinou hypotézu, ale váží přes všechny: $P(c\mid D)=\sum_h P(c\mid h)P(h\mid D)$.

\textbf{(b)} \emph{EM} maximalizuje věrohodnost při skrytých proměnných střídáním: \emph{E-krok} spočítá očekávané rozdělení skrytých proměnných (responsibility) za daných aktuálních parametru; \emph{M-krok} přepočítá parametry tak, aby maximalizovaly očekávanou (úplnou) log-věrohodnost. Iteruje do konvergence; věrohodnost monotónně neklesá, konverguje do lokálního maxima.

\textbf{(c)} Pro skrytý shluk $z\in\{1,\dots,K\}$ s vahami (priory) $\pi_k$ a modelem $p(x\mid\theta_k)$ je \emph{responsibility} bodu $x_i$ vuči shluku $k$:
\[
r_{ik}=P(z_i=k\mid x_i)=\frac{\pi_k\,p(x_i\mid\theta_k)}{\sum_{j=1}^K \pi_j\,p(x_i\mid\theta_j)}.
\]
\emph{M-krok}: $\pi_k=\frac1n\sum_i r_{ik}$ a parametry $\theta_k$ se odhadnou váženě responsibilitami (např.\ vážený průměr $\mu_k=\frac{\sum_i r_{ik}x_i}{\sum_i r_{ik}}$); u kategoriálních dat = vážené relativní četnosti.
\end{reseni}

\begin{jadro}
$h_{MAP}=\arg\max_h P(D\mid h)P(h)$ ($=$ ML při uniformním prioru). EM: E-krok responsibility $r_{ik}=\frac{\pi_k\,p(x_i\mid\theta_k)}{\sum_j\pi_j\,p(x_i\mid\theta_j)}$; M-krok $\pi_k=\frac1n\sum_i r_{ik}$ a vážené přepočty parametru.
\end{jadro}

\kalibrace{,,Bayesovské učení'' a EM jsou opakovaný ,,Základy UI'' okruh; EM zkouška chtěla \emph{jeden výpočetní update} (responsibility), ne jen pojem. Vzorec responsibility umět zpaměti.}
\past{MAP $\neq$ ML, liší se priorem. EM najde jen \emph{lokální} maximum (závisí na inicializaci); responsibility se musí normovat přes všechny shluky.}

\ulohahead{specializace UI-SU}{Markovské modely a skryté Markovské modely (HMM)}
\begin{zadani}
\begin{enumerate}[label=(\alph*)]
  \item \bod{1} Co je Markovský řetězec a Markovský předpoklad? Definujte skrytý Markovský model (HMM): skryté stavy, přechody, emise.
  \item \bod{2} Popište filtraci, predikci a vyhlazování a uveďte rekurentní vzorec pro filtraci (forward).
  \item \bod{3} Co počítá Viterbiho algoritmus (nejpravděpodobnější průchod)? Jak HMM souvisí s dynamickými Bayesovskými sítěmi?
\end{enumerate}
\end{zadani}

\begin{reseni}
\textbf{(a)} \emph{Markovský řetězec}: posloupnost stavu, kde budoucnost závisí jen na současném stavu (\emph{Markovský předpoklad} $P(X_t\mid X_{1:t-1})=P(X_t\mid X_{t-1})$). \emph{HMM}: skryté stavy $X_t$ s maticí přechodu $P(X_t\mid X_{t-1})$ a pozorování $E_t$ s modelem emisí $P(E_t\mid X_t)$; vidíme jen $E_t$, stavy odhadujeme.

\textbf{(b)} \emph{Filtrace}: $P(X_t\mid e_{1:t})$ (odhad současného stavu z dosavadních pozorování). \emph{Predikce}: $P(X_{t+k}\mid e_{1:t})$ (budoucí stav). \emph{Vyhlazování}: $P(X_k\mid e_{1:t})$ pro $k<t$ (minulý stav s využitím pozdějších pozorování). Rekurence filtrace (forward):
\[
P(X_t\mid e_{1:t})\ \propto\ P(e_t\mid X_t)\sum_{x_{t-1}}P(X_t\mid x_{t-1})\,P(x_{t-1}\mid e_{1:t-1}),
\]
tj.\ krok predikce (přes přechod) a poté přenásobení emisí a normalizace.

\textbf{(c)} \emph{Viterbi} dynamickým programováním najde \emph{nejpravděpodobnější posloupnost skrytých stavu} pro dané pozorování (max místo součtu ve forward rekurzi + zpětný chod). HMM je speciální případ \emph{dynamické Bayesovské sítě} (řetězec s jednou skrytou proměnnou na krok); DBN dovolují víc proměnných a strukturovanějších závislostí na časový krok.
\end{reseni}

\kalibrace{HMM (~12\,\%) je opakované ,,Základy UI'' téma (filtrace se zkoušela numericky). Definice + forward rekurence = 2-bodová záloha; pro 3 body umět jeden krok filtrace dosadit.}
\past{Filtrace je jen ,,dopředu'', vyhlazování využívá i \emph{budoucí} pozorování. Viterbi hledá nejpravděpodobnější \emph{celou cestu}, ne stav po stavu.}

\ulohahead{specializace UI-SU}{Markovské rozhodovací procesy (MDP) a zpětnovazební učení}
\begin{zadani}
\begin{enumerate}[label=(\alph*)]
  \item \bod{1} Definujte MDP (stavy, akce, přechody, odměny, diskont) a pojmy strategie (policy) a užitek (hodnota) stavu.
  \item \bod{2} Napište Bellmanovu rovnici pro optimální hodnotu a popište iteraci hodnot (value iteration).
  \item \bod{3} Co je iterace strategií? Popište zpětnovazební učení: Q-učení (update) a kompromis explorace vs exploitace.
\end{enumerate}
\end{zadani}

\begin{reseni}
\textbf{(a)} \emph{MDP} $=(S,A,P,R,\gamma)$: stavy $S$, akce $A$, přechody $P(s'\mid s,a)$, odměny $R(s,a,s')$, diskont $\gamma\in[0,1)$. \emph{Strategie} $\pi:S\to A$ určuje akci v každém stavu. \emph{Hodnota} (užitek) stavu při $\pi$: očekávaná diskontovaná suma odměn $V^\pi(s)=\mathbb{E}[\sum_{t\ge0}\gamma^t R_t\mid s_0=s,\pi]$.

\textbf{(b)} \emph{Bellmanova rovnice} pro optimální hodnotu:
\[
V^*(s)=\max_{a}\sum_{s'}P(s'\mid s,a)\big[R(s,a,s')+\gamma V^*(s')\big].
\]
\emph{Iterace hodnot}: inicializuj $V$, opakovaně aplikuj pravou stranu jako přiřazení (Bellmanuv update) pro všechny stavy, dokud se $V$ nestabilizuje; pak optimální strategie $\pi^*(s)=\arg\max_a\sum_{s'}P[\,R+\gamma V^*\,]$.

\textbf{(c)} \emph{Iterace strategií}: střídá \emph{vyhodnocení} dané strategie (vyřeš $V^\pi$) a \emph{zlepšení} (greedy vuči $V^\pi$); konverguje v konečně krocích. \emph{Zpětnovazební učení} (neznámé $P,R$): \emph{Q-učení} aktualizuje akční hodnotu z pozorovaného přechodu
\[
Q(s,a)\leftarrow Q(s,a)+\alpha\big[r+\gamma\max_{a'}Q(s',a')-Q(s,a)\big]
\]
(off-policy). \emph{Explorace vs exploitace}: je třeba občas zkusit i nepreferované akce (např.\ $\varepsilon$-greedy), jinak agent uvázne v suboptimu; SARSA je on-policy obdoba (místo $\max$ použije skutečně zvolenou $a'$).
\end{reseni}

\kalibrace{MDP/RL (~15\,\%) je opakovaná ,,Základy UI'' otázka (Bellmanova rovnice se zkoušela). Definice MDP + Bellman + iterace hodnot = 2-bodová záloha.}
\past{Bellman pro optimum má $\max_a$ (u dané strategie je tam suma přes $\pi$). Q-učení je off-policy ($\max$), SARSA on-policy - nezaměňuj.}

\ulohahead{specializace UI-SU}{Logické uvažování: DPLL, řetězení, rezoluce}
\begin{zadani}
\begin{enumerate}[label=(\alph*)]
  \item \bod{1} Zopakujte CNF a popište algoritmus DPLL pro řešení splnitelnosti (SAT): jednotková propagace a čistý literál.
  \item \bod{2} Co jsou Hornovy klauzule a jak funguje dopředné a zpětné řetězení?
  \item \bod{3} Popište rezoluci jako rozhodovací/dokazovací metodu (důkaz sporem, prázdná klauzule). K čemu je v predikátové logice unifikace?
\end{enumerate}
\end{zadani}

\begin{reseni}
\textbf{(a)} \emph{CNF} = konjunkce klauzulí (disjunkcí literálu). \emph{DPLL} prohledává přiřazení s prořezáváním: (1) \emph{jednotková propagace} - klauzule s jediným literálem vynutí jeho hodnotu (a zjednoduší ostatní); (2) \emph{čistý literál} - proměnná vyskytující se jen v jedné polaritě se nastaví tak, aby tyto klauzule splnila; (3) jinak rozděl podle nějaké proměnné (větev true/false) a rekurzivně. Konec: prázdná množina klauzulí = splnitelné, prázdná klauzule = konflikt (backtrack).

\textbf{(b)} \emph{Hornova klauzule} má nejvýše jeden pozitivní literál (tj.\ tvar $a_1\wedge\dots\wedge a_k\Rightarrow b$). \emph{Dopředné řetězení}: z známých faktu opakovaně odvozuj hlavy pravidel, jejichž těla jsou splněna, dokud nepřibývají nové fakty (data-driven). \emph{Zpětné řetězení}: od dotazovaného cíle rekurzivně dokazuj těla pravidel, která ho mají v hlavě (goal-driven, základ Prologu). Pro Hornovy klauzule je to lineární/efektivní.

\textbf{(c)} \emph{Rezoluce}: dokazuje $T\models\varphi$ sporem - k $T$ přidá $\neg\varphi$, převede na CNF a opakovaně tvoří rezolventy; odvození \emph{prázdné klauzule} znamená nesplnitelnost, tedy platnost $\varphi$. V \emph{predikátové} logice je nutná \emph{unifikace}: najití nejobecnějšího substitučního přiřazení proměnných, které dva literály ztotožní, aby šly rezolvovat.
\end{reseni}

\begin{jadro}
DPLL = jednotková propagace + čistý literál + větvení s backtrackingem. Hornovy klauzule ($\le1$ pozitivní literál) = dopředné/zpětné řetězení. Rezoluce = důkaz sporem (přidej $\neg\varphi$, odvoď prázdnou klauzuli); v predikátové logice s unifikací.
\end{jadro}

\kalibrace{Logické uvažování (~10\,\%, DPLL se zkoušelo) je vzácnější ,,Základy UI'' okruh. CNF + DPLL kroky = 2-bodová záloha; navíc se prolíná s matematickou logikou.}
\past{Jednotková propagace a čistý literál nejsou totéž. Rezoluce dokazuje sporem (přidej negaci cíle); cílem je odvodit prázdnou klauzuli.}

\ulohahead{specializace UI-SU}{Automatické plánování a reprezentace znalostí}
\begin{zadani}
\begin{enumerate}[label=(\alph*)]
  \item \bod{1} Popište formulaci plánovacího problému a definici operátoru (předpoklady a efekty), počáteční stav a cíl.
  \item \bod{2} Vysvětlete dopředné (progression) a zpětné (regression) plánování a jejich rozdíl.
  \item \bod{3} Co je situační kalkulus a problém rámce (frame problem) a jak se řeší?
\end{enumerate}
\end{zadani}

\begin{reseni}
\textbf{(a)} \emph{Plánovací problém} (např.\ STRIPS): počáteční stav (množina platných faktu), množina \emph{operátoru} a cílová podmínka. \emph{Operátor} má \emph{předpoklady} (precondition - co musí platit, aby šel provést) a \emph{efekty} (add list = co začne platit, delete list = co přestane platit). Plán = posloupnost operátoru vedoucí z počátečního stavu do stavu splňujícího cíl.

\textbf{(b)} \emph{Dopředné plánování}: začni v počátečním stavu a aplikuj použitelné operátory, dokud nedosáhneš cíle (prohledávání stavového prostoru vpřed); velký větvící faktor. \emph{Zpětné plánování}: začni od cíle a hledej operátory, jejichž efekty cíl (částečně) splňují, a nahraď cíl jejich předpoklady (regrese); pracuje jen s relevantními operátory, ale stavy jsou částečné (množiny podcílu).

\textbf{(c)} \emph{Situační kalkulus} popisuje svět v predikátové logice: \emph{fluenty} (vlastnosti závislé na situaci, např.\ $At(x,s)$), akce a funkci $do(a,s)$ pro situaci po akci. \emph{Problém rámce}: je třeba vyjádřit nejen co se akcí \emph{změní}, ale i co zustane \emph{beze změny} - jinak by se počet axiomu rozrostl pro každou dvojici akce/fluent. Řeší se \emph{successor-state axiomy} (pro každý fluent jeden axiom popisující, kdy v nové situaci platí), což stručně zachytí setrvačnost (co akce neovlivní, zustává).
\end{reseni}

\begin{jadro}
Operátor $=$ (precondition, add efekty, delete efekty). \emph{Dopředné} plánování jde od počátku vpřed, \emph{zpětné} regreduje od cíle přes efekty operátoru. Problém rámce $=$ jak vyjádřit, co se akcí nemění (successor-state axiomy).
\end{jadro}

\kalibrace{Plánování a reprezentace znalostí jsou vzácnější ,,Základy UI'' rodiny, ale legální draw do skoro vždy přítomného AI slotu. Bez této zálohy hrozí na takové otázce skóre 0-1, což přímo ohrožuje podlahu specializace ($\ge7/12$).}
\past{Dopředné vs zpětné nepleť: regrese jde od cíle a nahrazuje cíl předpoklady operátoru. Problém rámce není o ,,výpočtu'', ale o stručném popisu neměnnosti.}

\ulohahead{specializace UI-SU \textnormal{\footnotesize[úroveň 2]}}{Bayesovská síť: inference enumerací}
\begin{zadani}
Síť: dvě nezávislé příčiny $A,B$ a společný následek $C$ (hrany $A\to C$, $B\to C$). $P(A)=0{,}2$, $P(B)=0{,}5$ a
$P(C{=}t\mid A,B)$: $A t B t{:}\,0{,}9$; $A t B f{:}\,0{,}8$; $A f B t{:}\,0{,}7$; $A f B f{:}\,0{,}1$.
\begin{enumerate}[label=(\alph*)]
  \item \bod{1} Napište faktorizaci sdružené distribuce.
  \item \bod{2} Spočtěte $P(A{=}t, C{=}t)$ a $P(A{=}f, C{=}t)$ (vysčítání přes $B$).
  \item \bod{3} Spočtěte posterior $P(A{=}t\mid C{=}t)$.
\end{enumerate}
\end{zadani}

\begin{reseni}
\textbf{(a)} $P(A,B,C)=P(A)\,P(B)\,P(C\mid A,B)$ ($A,B$ jsou kořeny bez rodičů, nezávislé).

\textbf{(b)} Vysčítáme přes $B\in\{t,f\}$:
\[
P(A{=}t,C{=}t)=P(A{=}t)\!\sum_B P(B)P(C{=}t\mid A{=}t,B)=0{,}2\,(0{,}5\cdot0{,}9+0{,}5\cdot0{,}8)=0{,}2\cdot0{,}85=0{,}17.
\]
\[
P(A{=}f,C{=}t)=0{,}8\,(0{,}5\cdot0{,}7+0{,}5\cdot0{,}1)=0{,}8\cdot0{,}40=0{,}32.
\]

\textbf{(c)} $P(C{=}t)=0{,}17+0{,}32=0{,}49$, takže
\[
P(A{=}t\mid C{=}t)=\frac{0{,}17}{0{,}49}\approx0{,}347.
\]
\end{reseni}

\begin{jadro}
$P(\cdot)=\prod_i P(X_i\mid\text{rodiče})$. Dotaz: $P(Q\mid e)\propto\sum_{\text{skryté}}\prod(\text{faktory})$; spočti $P(Q,e)$ pro každou hodnotu $Q$ a normuj. Pozorování $C{=}t$ ,,nahoru'' zvýší pravděpodobnost příčin.
\end{jadro}

\kalibrace{Numerická inference v Bayes síti (enumerace) padla 2024-léto. Vysčítání + normalizace zpaměti = 2-3 body z AI slotu specializace.}
\past{Nezapomeň normovat ($\div P(e)$). Sčítáš \emph{součiny} faktoru přes skryté proměnné, ne pravděpodobnosti zvlášť.}

\ulohahead{specializace UI-SU \textnormal{\footnotesize[úroveň 2]}}{HMM: dvoukroková filtrace}
\begin{zadani}
Skrytý stav ,,prší'' $R\in\{t,f\}$. Přechody: $P(R_t{=}t\mid R_{t-1}{=}t)=0{,}7$, $P(R_t{=}t\mid R_{t-1}{=}f)=0{,}3$. Emise (deštník): $P(U{=}t\mid R{=}t)=0{,}9$, $P(U{=}t\mid R{=}f)=0{,}2$. Prior $P(R_0{=}t)=0{,}5$. Pozorování $U_1{=}t$, $U_2{=}t$.
\begin{enumerate}[label=(\alph*)]
  \item \bod{1} Spočtěte filtraci $P(R_1\mid U_1{=}t)$ (krok predikce + korekce).
  \item \bod{2} Spočtěte $P(R_2\mid U_1{=}t,U_2{=}t)$.
  \item \bod{3} Napište obecný rekurentní vzorec, který jste použili.
\end{enumerate}
\end{zadani}

\begin{reseni}
\textbf{(a)} \emph{Predikce} z prioru: $P(R_1{=}t)=0{,}7\cdot0{,}5+0{,}3\cdot0{,}5=0{,}5$, $P(R_1{=}f)=0{,}5$.
\emph{Korekce} emisí $U_1{=}t$: $\propto(0{,}9\cdot0{,}5,\ 0{,}2\cdot0{,}5)=(0{,}45,\ 0{,}10)$. Normalizace ($/0{,}55$): $P(R_1{=}t\mid U_1)=0{,}818$, $P(R_1{=}f\mid U_1)=0{,}182$.

\textbf{(b)} \emph{Predikce}: $P(R_2{=}t)=0{,}7\cdot0{,}818+0{,}3\cdot0{,}182=0{,}573+0{,}055=0{,}627$, $P(R_2{=}f)=0{,}373$.
\emph{Korekce} $U_2{=}t$: $\propto(0{,}9\cdot0{,}627,\ 0{,}2\cdot0{,}373)=(0{,}564,\ 0{,}075)$. Normalizace ($/0{,}639$): $P(R_2{=}t\mid U_1,U_2)\approx0{,}883$, $P(R_2{=}f\mid\cdot)\approx0{,}117$.

\textbf{(c)} $P(R_t\mid e_{1:t})\propto P(e_t\mid R_t)\sum_{r_{t-1}}P(R_t\mid r_{t-1})\,P(r_{t-1}\mid e_{1:t-1})$ (predikce přes přechod, korekce emisí, normalizace).
\end{reseni}

\begin{jadro}
Filtrace = střídej predikci (násob přechodovou maticí) a korekci (násob emisí pozorování) a normuj. Dva deštníky táhnou ,,prší'' nahoru: $0{,}5\to0{,}818\to0{,}883$.
\end{jadro}

\kalibrace{Numerická filtrace HMM padla 2023-léto. Dvoukrokový výpočet zpaměti = 2-3 body z AI slotu; HMM ~12\,\%.}
\past{Pořadí: nejdřív predikce (přechod), pak korekce (emise), pak normalizace. Emisi nesmíš použít dřív než přechod.}

\ulohahead{specializace UI-SU \textnormal{\footnotesize[úroveň 2]}}{EM algoritmus: jedna iterace}
\begin{zadani}
Směs dvou shluku $Z\in\{1,2\}$ nad binárním příznakem $x\in\{0,1\}$. Aktuální parametry: $\pi=(0{,}5,\,0{,}5)$, $p(x{=}1\mid Z{=}1)=0{,}8$, $p(x{=}1\mid Z{=}2)=0{,}4$. Data: $x_1=1$, $x_2=0$.
\begin{enumerate}[label=(\alph*)]
  \item \bod{1} E-krok: spočtěte responsibility $r_{i1}=P(Z{=}1\mid x_i)$ pro oba body.
  \item \bod{2} M-krok: aktualizujte váhy $\pi_1,\pi_2$.
  \item \bod{3} M-krok: aktualizujte $p(x{=}1\mid Z{=}1)$ a $p(x{=}1\mid Z{=}2)$.
\end{enumerate}
\end{zadani}

\begin{reseni}
\textbf{(a)} $r_{i1}=\dfrac{\pi_1 p(x_i\mid 1)}{\pi_1 p(x_i\mid 1)+\pi_2 p(x_i\mid 2)}$.
Pro $x_1{=}1$: $\dfrac{0{,}5\cdot0{,}8}{0{,}5\cdot0{,}8+0{,}5\cdot0{,}4}=\dfrac{0{,}4}{0{,}6}=0{,}667$, takže $r_{11}=0{,}667$, $r_{12}=0{,}333$.
Pro $x_2{=}0$ (použij $p(x{=}0\mid1)=0{,}2$, $p(x{=}0\mid2)=0{,}6$): $\dfrac{0{,}5\cdot0{,}2}{0{,}5\cdot0{,}2+0{,}5\cdot0{,}6}=\dfrac{0{,}1}{0{,}4}=0{,}25$, takže $r_{21}=0{,}25$, $r_{22}=0{,}75$.

\textbf{(b)} $\pi_k=\tfrac1n\sum_i r_{ik}$: $\pi_1=\tfrac{0{,}667+0{,}25}{2}=0{,}458$, $\pi_2=0{,}542$.

\textbf{(c)} $p(x{=}1\mid Z{=}k)=\dfrac{\sum_i r_{ik}\,[x_i{=}1]}{\sum_i r_{ik}}$:
\[
p(x{=}1\mid1)=\frac{0{,}667\cdot1+0{,}25\cdot0}{0{,}667+0{,}25}=\frac{0{,}667}{0{,}917}\approx0{,}727,\quad
p(x{=}1\mid2)=\frac{0{,}333}{0{,}333+0{,}75}=\frac{0{,}333}{1{,}083}\approx0{,}308.
\]
Shluk 1 se přiostřil k $x{=}1$, shluk 2 k $x{=}0$.
\end{reseni}

\begin{jadro}
E-krok: $r_{ik}=\frac{\pi_k p(x_i\mid\theta_k)}{\sum_j\pi_j p(x_i\mid\theta_j)}$. M-krok: $\pi_k=\frac1n\sum_i r_{ik}$, parametry = responsibilitami vážené (relativní) četnosti / průměry.
\end{jadro}

\kalibrace{Jeden EM update (responsibility + M-krok) padl 2025-léto. Toto je přesný typ; nacvičit = 2-3 body z AI slotu.}
\past{Responsibility se normuje přes shluky. M-krok je \emph{vážený} responsibilitami, ne prosté četnosti. Pro $x{=}0$ použij $p(x{=}0\mid\cdot)=1-p(x{=}1\mid\cdot)$.}

\ulohahead{specializace UI-SU \textnormal{\footnotesize[úroveň 2]}}{Minimax, alfa-beta a A* na konkrétním příkladu}
\begin{zadani}
\begin{enumerate}[label=(\alph*)]
  \item \bod{1} Strom: kořen MAX má dva syny MIN. Levý MIN má listy $3,5$; pravý MIN má listy $2,9$. Spočtěte minimaxovou hodnotu.
  \item \bod{2} Ukažte, kterou hodnotu alfa-beta ořeže (prochází zleva doprava).
  \item \bod{3} V grafu s hranami $S{-}A{=}1,\ A{-}B{=}2,\ A{-}G{=}5,\ B{-}G{=}1$ a heuristikou $h(S){=}4,h(A){=}3,h(B){=}1,h(G){=}0$ najděte A* cestu z $S$ do $G$.
\end{enumerate}
\end{zadani}

\begin{reseni}
\textbf{(a)} Levý MIN $=\min(3,5)=3$, pravý MIN $=\min(2,9)=2$. Kořen MAX $=\max(3,2)=3$.

\textbf{(b)} Po vyhodnocení levého podstromu je $\alpha=3$. V pravém MIN se přečte první list $2$; protože hodnota MIN uzlu bude $\le2\le\alpha=3$, MAX si ho stejně nevybere - druhý list $9$ se \emph{ořízne} (alfa cutoff). Výsledek $3$ je stejný jako u minimaxu.

\textbf{(c)} $f=g+h$. Start $S$: $g{=}0,f{=}4$. Rozvoj $S$: $A$ $g{=}1,f{=}1{+}3{=}4$; $B$ přes $S$ neexistuje hrana, jen přes $A$. Rozvoj $A$ (nejmenší $f{=}4$): $B$ $g{=}1{+}2{=}3,f{=}3{+}1{=}4$; $G$ $g{=}1{+}5{=}6,f{=}6$. Rozvoj $B$ ($f{=}4$): $G$ $g{=}3{+}1{=}4,f{=}4$. Rozvoj $G$ ($f{=}4$) = cíl. Cesta $S\to A\to B\to G$ s cenou $4$ (lepší než $S\to A\to G$ s cenou $6$; přímá hrana $S{-}B$ v grafu není).
\end{reseni}

\begin{jadro}
Minimax: listy nahoru, MAX bere $\max$, MIN $\min$. Alfa-beta ořeže větev, jakmile $\alpha\ge\beta$. A* rozvíjí nejmenší $f{=}g{+}h$; s přípustnou/konzistentní $h$ je optimální.
\end{jadro}

\kalibrace{Ruční alfa-beta a A* jsou opakované AI typy (hry ~22\,\%, prohledávání ~15\,\%). Jeden vyřešený strom + graf = 2-3 body z AI slotu.}
\past{Alfa-beta nemění výsledek, jen ořezává. U A* musí být $h$ přípustná, jinak nemusí najít optimum; rozvíjej vždy uzel s nejmenším $f$.}

\ulohahead{specializace UI-SU \textnormal{\footnotesize[úroveň 2]}}{Bayesovské učení: posterior a Bayes-optimal predikce}
\begin{zadani}
Pytlík je jedné ze dvou hypotéz: $h_1$ s $P(\text{třešeň})=0{,}8$, nebo $h_2$ s $P(\text{třešeň})=0{,}4$. Prior $P(h_1)=P(h_2)=0{,}5$. Vytáhneme jeden bonbon a je to \emph{třešeň}.
\begin{enumerate}[label=(\alph*)]
  \item \bod{1} Určete ML a MAP hypotézu po tomto pozorování.
  \item \bod{2} Spočtěte posterior $P(h_1\mid \text{třešeň})$ a $P(h_2\mid \text{třešeň})$.
  \item \bod{3} Spočtěte Bayes-optimal predikci $P(\text{další} = \text{třešeň}\mid \text{data})$.
\end{enumerate}
\end{zadani}

\begin{reseni}
\textbf{(a)} Věrohodnost pozorování ,,třešeň'': $P(\text{tř}\mid h_1)=0{,}8 > P(\text{tř}\mid h_2)=0{,}4$, takže \emph{ML} hypotéza je $h_1$. Prior je uniformní, takže \emph{MAP} $=$ ML $=h_1$.

\textbf{(b)} $P(h_1\mid\text{tř})\propto0{,}5\cdot0{,}8=0{,}4$; $P(h_2\mid\text{tř})\propto0{,}5\cdot0{,}4=0{,}2$. Normalizace ($/0{,}6$): $P(h_1\mid\text{tř})=\tfrac{0{,}4}{0{,}6}=0{,}667$, $P(h_2\mid\text{tř})=0{,}333$.

\textbf{(c)} Bayes-optimal váží přes hypotézy:
\[
P(\text{další}{=}\text{tř}\mid\text{data})=\sum_h P(\text{tř}\mid h)P(h\mid\text{data})=0{,}8\cdot0{,}667+0{,}4\cdot0{,}333=0{,}533+0{,}133=0{,}667.
\]
(Pozn.: predikce přes MAP samotnou by dala $0{,}8$; Bayes-optimal je obezřetnější, protože $h_2$ ještě není vyloučena.)
\end{reseni}

\begin{jadro}
Posterior $\propto$ věrohodnost $\times$ prior, pak normuj. ML $=\arg\max$ věrohodnost, MAP $=\arg\max$ posterior. Bayes-optimal predikce $=\sum_h P(\cdot\mid h)P(h\mid\text{data})$ (ne jen přes jednu hypotézu).
\end{jadro}

\kalibrace{Posterior + Bayes-optimal predikce padly 2025-jaro. Přesný typ; nacvičit = 2-3 body z AI slotu.}
\past{Bayes-optimal $\neq$ predikce přes MAP hypotézu - váží přes \emph{všechny} hypotézy jejich posteriorem. Nezapomeň normovat posterior.}

\ulohahead{specializace UI-SU \textnormal{\footnotesize[úroveň 2]}}{AC-3 krok za krokem}
\begin{zadani}
Proměnné $X,Y,Z$ s doménami $\{1,2,3,4\}$ a omezeními $X<Y$ a $Y<Z$. Proveďte AC-3 a veďte frontu hran.
\begin{enumerate}[label=(\alph*)]
  \item \bod{1} Vypište počáteční frontu hran.
  \item \bod{2} Zpracujte hrany a u každé revize uveďte zúžení domény a nově přidané hrany.
  \item \bod{3} Uveďte výsledné domény a vysvětlete, proč musela být znovu zařazena hrana $(X,Y)$.
\end{enumerate}
\end{zadani}

\begin{reseni}
\textbf{(a)} Hrany (obě orientace každého omezení): $(X,Y),(Y,X),(Y,Z),(Z,Y)$.

\textbf{(b)}
\begin{itemize}\setlength\itemsep{1pt}
  \item $(X,Y)$ [$X<Y$]: $X{=}4$ nemá $Y>4$ $\Rightarrow$ $D_X=\{1,2,3\}$. (jediný soused $X$ je $Y$ = právě porovnávaná hrana, nic nepřidáváme)
  \item $(Y,X)$ [pro $Y$ existuje $X<Y$]: $Y{=}1$ nemá $X<1$ $\Rightarrow$ $D_Y=\{2,3,4\}$. (přidej $(Z,Y)$; soused $X$ je právě porovnávaný)
  \item $(Y,Z)$ [$Y<Z$]: $Y{=}4$ nemá $Z>4$ $\Rightarrow$ $D_Y=\{2,3\}$. (přidej $(X,Y)$; soused $Z$ je právě porovnávaný)
  \item $(Z,Y)$ [pro $Z$ existuje $Y<Z$]: $Z{=}1$ (žádné $Y<1$) a $Z{=}2$ (žádné $Y<2$ v $\{2,3\}$) vypadnou $\Rightarrow$ $D_Z=\{3,4\}$. (jediný soused $Z$ je $Y$ = právě porovnávaná hrana)
  \item $(X,Y)$ (znovu): $X{=}3$ nemá $Y>3$ v $\{2,3\}$ $\Rightarrow$ $D_X=\{1,2\}$.
\end{itemize}

\textbf{(c)} Výsledek: $D_X=\{1,2\}$, $D_Y=\{2,3\}$, $D_Z=\{3,4\}$. Hrana $(X,Y)$ se musela zařadit znovu, protože při revizi $(Y,Z)$ se zmenšila doména $Y$; tím mohly ztratit podporu hodnoty v $X$, které se na zúžené $Y$ odvolávaly (proto AC-3 vrací do fronty hrany $(X_k,Y)$ od sousedu $X_k$ změněné proměnné, kromě právě porovnávané).
\end{reseni}

\begin{jadro}
AC-3: fronta orientovaných hran; \textsc{Revise}$(X_i,X_j)$ smaže z $D_i$ hodnoty bez podpory v $D_j$. Zmenší-li se $D_i$, vrať do fronty hrany $(X_k,X_i)$ pro sousedy $X_k\ne X_j$. Konec při prázdné frontě (nebo prázdné doméně = bez řešení). Složitost $O(e\,d^3)$.
\end{jadro}

\kalibrace{Krokování AC-3 padlo 2026-jaro. Vedení fronty a šíření zpaměti = 2-3 body z AI slotu (CSP ~22\,\%).}
\past{Po zúžení domény se musí znovu prověřit \emph{sousední} hrany - na to se zapomíná. Hranová konzistence ani po doběhnutí nezaručuje řešení.}

\ulohahead{specializace UI-SU \textnormal{\footnotesize[úroveň 2]}}{MDP: iterace hodnot}
\begin{zadani}
Stavy $A,B$ a cíl $C$ (terminální, $V(C)=0$). Diskont $\gamma=0{,}9$. Z $A$ jsou dvě akce: ,,do B'' (deterministicky do $B$, odměna $-1$) a ,,přímo'' (deterministicky do $C$, odměna $-3$). Z $B$ akce ,,do C'' (do $C$, odměna $-1$).
\begin{enumerate}[label=(\alph*)]
  \item \bod{1} Napište Bellmanovu rovnici pro $V^*(A)$ a $V^*(B)$.
  \item \bod{2} Proveďte iteraci hodnot z $V_0\equiv0$ (dva sweepy).
  \item \bod{3} Uveďte optimální hodnoty a strategii.
\end{enumerate}
\end{zadani}

\begin{reseni}
\textbf{(a)} $V^*(B)=-1+\gamma V^*(C)=-1$. $V^*(A)=\max\{\,-1+\gamma V^*(B),\ \ -3+\gamma V^*(C)\,\}$.

\textbf{(b)} $V_0(A)=V_0(B)=0$.
Sweep 1: $V_1(B)=-1+0{,}9\cdot0=-1$; $V_1(A)=\max\{-1+0{,}9\cdot0,\ -3+0{,}9\cdot0\}=\max\{-1,-3\}=-1$ (akce ,,do B'').
Sweep 2: $V_2(B)=-1$; $V_2(A)=\max\{-1+0{,}9\cdot(-1),\ -3\}=\max\{-1{,}9,\ -3\}=-1{,}9$ (akce ,,do B'').
Sweep 3 by dal totéž ($V_3(A)=\max\{-1+0{,}9\cdot(-1),-3\}=-1{,}9$) $\Rightarrow$ konvergence.

\textbf{(c)} $V^*(B)=-1$, $V^*(A)=-1{,}9$. Optimální strategie: z $A$ jdi ,,do B'', z $B$ jdi ,,do C'' (cesta $A\to B\to C$ s celkovou diskontovanou odměnou $-1{,}9$ je lepší než přímo $A\to C$ s $-3$).
\end{reseni}

\begin{jadro}
Bellman: $V^*(s)=\max_a\sum_{s'}P(s'\mid s,a)[R+\gamma V^*(s')]$. Iterace hodnot = opakuj tento update pro všechny stavy do konvergence; strategie = $\arg\max_a$.
\end{jadro}

\kalibrace{Bellmanova rovnice / iterace hodnot padla 2025-léto. Jeden-dva sweepy ručně = 2-3 body z AI slotu (MDP ~15\,\%).}
\past{Diskont $\gamma$ násobí \emph{hodnotu následníka}, ne odměnu kroku. U optima je $\max_a$ (u dané strategie jen ta jedna akce).}

\ulohahead{specializace UI-SU \textnormal{\footnotesize[úroveň 3]}}{POMDP a aktualizace beliefu}
\begin{zadani}
Částečně pozorovatelný MDP se dvěma stavy $s_1,s_2$. Pro akci $a$: $T(s_1\!\to\!s_1)=0{,}8$, $T(s_1\!\to\!s_2)=0{,}2$, $T(s_2\!\to\!s_1)=0{,}3$, $T(s_2\!\to\!s_2)=0{,}7$. Model pozorování: $O(o\mid s_1)=0{,}9$, $O(o\mid s_2)=0{,}2$. Počáteční belief $b=(0{,}5,\,0{,}5)$.
\begin{enumerate}[label=(\alph*)]
  \item \bod{1} Definujte POMDP a pojem belief (víra o stavu). \emph{(Toto je rozsah požadavku: ,,POMDP - základní definice''.)}
  \item Napište vzorec aktualizace beliefu po akci $a$ a pozorování $o$. \emph{(volitelně, nad rámec požadavku)}
  \item Spočtěte nový belief po provedení $a$ a pozorování $o$. \emph{(volitelně, nad rámec požadavku)}
\end{enumerate}
\end{zadani}

\begin{reseni}
\textbf{(a)} \emph{POMDP} $=(S,A,T,R,\Omega,O,\gamma)$: jako MDP, ale agent stav přímo \emph{nevidí}; po akci dostane jen \emph{pozorování} $o\in\Omega$ s pravděpodobností $O(o\mid s',a)$ (obecně $O:\Omega\times S\times A\to[0,1]$ - pozorování smí záviset i na akci). \emph{Belief} $b$ je rozdělení pravděpodobnosti nad stavy (co agent ví); rozhoduje se podle něj (POMDP lze chápat jako MDP nad spojitým prostorem beliefu).

\textbf{(b)} Aktualizace (předpověď přes přechod + korekce pozorováním, Bayes):
\[
b'(s')\ \propto\ O(o\mid s',a)\sum_{s} T(s'\mid s,a)\,b(s),
\]
poté normalizace, aby $\sum_{s'}b'(s')=1$.

\textbf{(c)} (V tomto příkladu $O$ na akci nezávisí, tj.\ $O(o\mid s')$.) \emph{Predikce}: $\tilde b(s_1)=0{,}8\cdot0{,}5+0{,}3\cdot0{,}5=0{,}55$, $\tilde b(s_2)=0{,}2\cdot0{,}5+0{,}7\cdot0{,}5=0{,}45$.
\emph{Korekce} pozorováním $o$: $\propto(0{,}9\cdot0{,}55,\ 0{,}2\cdot0{,}45)=(0{,}495,\ 0{,}09)$. Normalizace ($/0{,}585$): $b'(s_1)\approx0{,}846$, $b'(s_2)\approx0{,}154$.
\end{reseni}

\begin{jadro}
POMDP = MDP + skrytý stav + pozorování $O(o\mid s',a)$; agent drží belief $b$ (rozdělení nad stavy). Update: $b'(s')\propto O(o\mid s')\sum_s T(s'\mid s,a)b(s)$, pak normuj. (Stejná logika jako filtrace HMM s akcemi.)
\end{jadro}

\kalibrace{Požadavek chce u POMDP \emph{jen základní definici} (část a) - to je povinné jádro. Belief update (b, c) je nad rámec požadavku; uvádíme ho pro pochopení (je to táž logika jako filtrace HMM), ale na zkoušce se nevyžaduje.}
\past{Belief update = predikce (přechod) PAK korekce (pozorování) + normalizace, jako u HMM. Belief je rozdělení, ne jeden ,,nejpravděpodobnější'' stav.}

\ulohahead{specializace UI-SU \textnormal{\footnotesize[úroveň 3]}}{Rezoluce v predikátové logice a unifikace}
\begin{zadani}
\begin{enumerate}[label=(\alph*)]
  \item \bod{1} Co je unifikace a nejobecnější unifikátor (MGU)? Co je skolemizace?
  \item \bod{2} Z předpokladu ,,$\forall x\,(P(x)\to Q(x))$'' a ,,$P(a)$'' rezolucí odvoďte $Q(a)$.
  \item \bod{3} Vysvětlete, proč je rezoluce důkaz sporem a jakou roli hraje převod do klauzulární formy.
\end{enumerate}
\end{zadani}

\begin{reseni}
\textbf{(a)} \emph{Unifikace} hledá substituci proměnných, která dva literály (termy) ztotožní; \emph{nejobecnější unifikátor} (MGU) je ta nejméně specializující (každá jiná unifikace je její instancí). \emph{Skolemizace} odstraní existenční kvantifikátory při převodu do klauzulární formy: $\exists y$ se nahradí novou \emph{Skolemovou} konstantou (není-li pod $\forall$) nebo \emph{Skolemovou funkcí} závislých univerzálních proměnných. Skolemizace zachovává \emph{splnitelnost} (equisatisfiabilitu), ne obecně logickou ekvivalenci - rezoluci to stačí, protože dokazuje sporem (jde jí o nesplnitelnost).

\textbf{(b)} Klauzulární forma: $\forall x(P(x)\to Q(x))$ dá klauzuli $\{\neg P(x),\,Q(x)\}$; dále $\{P(a)\}$. Pro důkaz $Q(a)$ přidáme negaci cíle $\{\neg Q(a)\}$.
\begin{itemize}\setlength\itemsep{1pt}
  \item Rezolvuj $\{\neg P(x),Q(x)\}$ a $\{P(a)\}$ s MGU $\{x/a\}$: vznikne $\{Q(a)\}$.
  \item Rezolvuj $\{Q(a)\}$ a $\{\neg Q(a)\}$: vznikne prázdná klauzule $\square$.
\end{itemize}
Odvození prázdné klauzule dokazuje $Q(a)$.

\textbf{(c)} Rezoluce dokazuje $T\models\varphi$ \emph{sporem}: k teorii přidá $\neg\varphi$ a snaží se odvodit \emph{prázdnou klauzuli} (= nesplnitelnost $T\cup\{\neg\varphi\}$). Aby to šlo mechanicky, převede se vše do \emph{klauzulární formy} (PNF $\to$ skolemizace $\to$ CNF $\to$ množina klauzulí); rezoluce s unifikací je pak úplná zamítací procedura pro predikátovou logiku.
\end{reseni}

\begin{jadro}
Unifikace = substituce ztotožňující literály, MGU = nejobecnější. Skolemizace odstraní $\exists$. Rezoluce: přidej $\neg$cíl, převeď na klauzule, rezolvuj s MGU do prázdné klauzule (důkaz sporem).
\end{jadro}

\kalibrace{Rezoluce a Hornovy klauzule jsou v požadavku Základu UI i logiky; úroveň 3 doplňuje predikátovou rezoluci s unifikací (těžší než výroková).}
\past{Skolemizace ruší $\exists$, ne $\forall$. Rezoluce dokazuje sporem - bez přidání negace cíle nic nedokážeš.}

\ulohahead{specializace UI-SU \textnormal{\footnotesize[úroveň 3]}}{Plánování zpětnou regresí (STRIPS)}
\begin{zadani}
Operátor \texttt{Stack(A,B)}: předpoklady $\{Hold(A),\,Clear(B)\}$, přidává $\{On(A,B),\,Clear(A),\,HandEmpty\}$, ruší $\{Hold(A),\,Clear(B)\}$. Cíl je $On(A,B)$.
\begin{enumerate}[label=(\alph*)]
  \item \bod{1} Připomeňte princip zpětného (regresního) plánování.
  \item \bod{2} Proveďte jeden krok regrese cíle $On(A,B)$ přes operátor \texttt{Stack(A,B)}.
  \item \bod{3} Porovnejte dopředné a zpětné plánování (větvení, relevance operátoru).
\end{enumerate}
\end{zadani}

\begin{reseni}
\textbf{(a)} \emph{Zpětné (regresní) plánování} jde od cíle: vyber operátor, jehož \emph{přidávací} efekty obsahují nějaký cílový literál, a \emph{regreduj} cíl - nahraď splněné cílové literály předpoklady operátoru. Pracuje jen s \emph{relevantními} operátory (těmi, co k cíli přispívají).

\textbf{(b)} Cíl $G=\{On(A,B)\}$. Operátor \texttt{Stack(A,B)} přidává $On(A,B)$, je tedy relevantní a \emph{konzistentní} (neruší žádný jiný cílový literál). Regrese:
\[
G' = \big(G\setminus \text{add}(\texttt{Stack})\big)\,\cup\,\text{pre}(\texttt{Stack}) = \big(\{On(A,B)\}\setminus\{On(A,B),\dots\}\big)\cup\{Hold(A),Clear(B)\}.
\]
Nový podcíl $G'=\{Hold(A),\,Clear(B)\}$ - tj.\ abychom dosáhli $On(A,B)$ pomocí \texttt{Stack(A,B)}, musíme nejprve dosáhnout, že držíme $A$ a $B$ je volné.

\textbf{(c)} \emph{Dopředné} plánování jde od počátečního stavu a aplikuje použitelné operátory; má velký větvící faktor (mnoho použitelných akcí), ale konkrétní úplné stavy. \emph{Zpětné} jde od cíle a uvažuje jen \emph{relevantní} operátory (menší větvení k cíli), ale pracuje s \emph{částečnými} stavy (množinami podcílu) a musí hlídat konzistenci (operátor nesmí rušit jiný cílový literál).
\end{reseni}

\begin{jadro}
Regrese cíle $G$ přes operátor $o$ (jehož add pokrývá část $G$ a nic z $G$ neruší): $G'=(G\setminus\text{add}(o))\cup\text{pre}(o)$. Zpětné plánování = relevantní operátory od cíle; dopředné = použitelné operátory od počátku.
\end{jadro}

\kalibrace{Automatické plánování (dopředné/zpětné) je v požadavku Základu UI; úroveň 3 přidává vypracovaný regresní krok ke konceptu z úrovně 1.}
\past{Regrese odečítá z cíle \emph{add} efekty a přidává \emph{předpoklady}; operátor je použitelný jen, neruší-li jiný cílový literál (konzistence).}


\section{Zkoušky nanečisto a dril}
\subsection*{Zkouška nanečisto 1}
\addcontentsline{toc}{subsection}{Zkouška nanečisto 1}
\begin{infobox}[title={Pokyny}]
\footnotesize 8 otázek (2 matematika, 2 informatika, 4 specializace), každá 0-3 body. Cíl: společné $\ge5/12$ \textbf{a} specializace $\ge7/12$. Zakryj řešení, vyřeš na papír, pak porovnej.
\end{infobox}

\mockq{Otázka 1}{matematika}{Průběh funkce}
\begin{zadani}Najděte lokální extrémy funkce $f(x)=x^3-3x$ a spočtěte $\lim_{x\to0}\frac{1-\cos x}{x^2}$.\end{zadani}
\begin{reseni}$f'(x)=3x^2-3=0\Rightarrow x=\pm1$; $f''=6x$, takže $x{=}{-}1$ je maximum ($f={2}$), $x{=}1$ minimum ($f={-}2$). Limita: $1-\cos x=\frac{x^2}{2}+o(x^2)$, tedy $\lim=\frac12$.\end{reseni}

\mockq{Otázka 2}{matematika}{Vlastní čísla a diagonalizace}
\begin{zadani}Pro $A=\begin{psmallmatrix}1&2\\2&1\end{psmallmatrix}$ najděte vlastní čísla, vlastní vektory a rozhodněte o diagonalizovatelnosti.\end{zadani}
\begin{reseni}$\det(A-\lambda I)=(1-\lambda)^2-4=\lambda^2-2\lambda-3=(\lambda-3)(\lambda+1)$, takže $\lambda=3,-1$. Vektory: $\lambda{=}3\to(1,1)$, $\lambda{=}{-}1\to(1,-1)$. Dvě různá vlastní čísla $\Rightarrow$ diagonalizovatelná, $A=PDP^{-1}$, $P=\begin{psmallmatrix}1&1\\1&-1\end{psmallmatrix}$, $D=\operatorname{diag}(3,-1)$.\end{reseni}

\mockq{Otázka 3}{informatika}{Konečný automat}
\begin{zadani}Sestrojte DFA pro binární čísla (čtená odshora) dělitelná 3.\end{zadani}
\begin{reseni}Stavy = zbytek po dělení 3: $\{0,1,2\}$, počátek a jediný přijímající stav $0$. Přechod při bitu $b$: $q\to(2q+b)\bmod 3$. Tedy $0\xrightarrow{0}0,\;0\xrightarrow{1}1,\;1\xrightarrow{0}2,\;1\xrightarrow{1}0,\;2\xrightarrow{0}1,\;2\xrightarrow{1}2$. (Funguje, protože přečtení bitu odpovídá $n\mapsto 2n+b$.)\end{reseni}

\mockq{Otázka 4}{informatika}{Virtuální paměť}
\begin{zadani}Stránka 4\,KiB. Přeložte virtuální adresu \texttt{0x5678}, je-li stránka 5 namapována na rámec 2. Co je race condition?\end{zadani}
\begin{reseni}Offset 12 bitu: offset $=$\texttt{0x678}, číslo stránky $=$\texttt{0x5}. Stránka $5\to$ rámec $2$, fyzická adresa $=(2\ll12)\,|\,\texttt{0x678}=\texttt{0x2678}$. \emph{Race condition}: výsledek závisí na prokládání vláken nad sdíleným stavem; řeší se vzájemným vyloučením (zámek/kritická sekce).\end{reseni}

\mockq{Otázka 5}{specializace}{Logistická regrese}
\begin{zadani}Napište predikci, ztrátu a jeden krok SGD pro $w=(0,0)$, $x=(1,1)$, $y=1$, $\eta=0{,}5$.\end{zadani}
\begin{reseni}$p=\sigma(0)=0{,}5$; ztráta NLL $=-\ln p$. Gradient $(p-y)x=(-0{,}5,-0{,}5)$, krok $w\leftarrow(0,0)-0{,}5\cdot(-0{,}5,-0{,}5)=(0{,}25,\,0{,}25)$. Nová predikce $\sigma(0{,}5)\approx0{,}62$ (posun k 1).\end{reseni}

\mockq{Otázka 6}{specializace}{Rozhodovací strom}
\begin{zadani}Uzel se 4 kladnými a 4 zápornými příklady. Spočtěte entropii a informační zisk štěpení na čisté listy $(4,0)$ a $(0,4)$.\end{zadani}
\begin{reseni}$H=-\tfrac12\log_2\tfrac12-\tfrac12\log_2\tfrac12=1$ bit. Čisté listy mají $H=0$, vážená entropie $0$, takže $\mathrm{IG}=1-0=1$ bit (dokonalé rozdělení).\end{reseni}

\mockq{Otázka 7}{specializace}{Shlukování k-means}
\begin{zadani}Body $1,2,8,9$ na přímce, $k=2$, počáteční centra $1$ a $9$. Proveďte jednu iteraci.\end{zadani}
\begin{reseni}Přiřazení: $1,2\to c_1{=}1$ (blíž), $8,9\to c_2{=}9$. Nová centra: $c_1=\tfrac{1+2}{2}=1{,}5$, $c_2=\tfrac{8+9}{2}=8{,}5$. Další iterace už nic nezmění (konvergence), ztráta $J=0{,}25\cdot4=1{,}0$.\end{reseni}

\mockq{Otázka 8}{specializace}{Bayesovská inference}
\begin{zadani}$P(A)=0{,}3$, $P(B\mid A)=0{,}8$, $P(B\mid\neg A)=0{,}4$. Spočtěte $P(A\mid B)$.\end{zadani}
\begin{reseni}$P(A,B)=0{,}3\cdot0{,}8=0{,}24$, $P(\neg A,B)=0{,}7\cdot0{,}4=0{,}28$, $P(B)=0{,}52$. $P(A\mid B)=\tfrac{0{,}24}{0{,}52}\approx0{,}46$.\end{reseni}

\begin{infobox}[title={Bodování (orientačně)}]
\footnotesize Každá úplná odpověď $\approx3$, definice + hlavní postup $\approx2$. Na průchod stačí: u 4 společných (Q1-Q4) nasbírat $\ge5$ (např.\ $2{+}2{+}1{+}1$), u 4 specializačních (Q5-Q8) $\ge7$ (např.\ $2{+}2{+}2{+}1$). Pokud někde vyjde $\le1$, dožeň to jinde - hlídej obě podlahy zvlášť.
\end{infobox}

\subsection*{Zkouška nanečisto 2}
\addcontentsline{toc}{subsection}{Zkouška nanečisto 2}
\begin{infobox}[title={Pokyny}]
\footnotesize 8 otázek (2 matematika, 2 informatika, 4 specializace), 0-3 body. Cíl: společné $\ge5/12$ \textbf{a} specializace $\ge7/12$.
\end{infobox}

\mockq{Otázka 1}{matematika}{Integrály}
\begin{zadani}Spočtěte $\int_0^2 (x^2-1)\,dx$ a obsah mezi $y=x$ a $y=x^2$ na $[0,1]$.\end{zadani}
\begin{reseni}$\int_0^2(x^2-1)\,dx=[\tfrac{x^3}{3}-x]_0^2=\tfrac83-2=\tfrac23$. Obsah $=\int_0^1(x-x^2)\,dx=[\tfrac{x^2}{2}-\tfrac{x^3}{3}]_0^1=\tfrac12-\tfrac13=\tfrac16$.\end{reseni}

\mockq{Otázka 2}{matematika}{Teorie grafu}
\begin{zadani}Definujte barevnost. Je $K_5$ rovinný? Zdůvodněte přes $E\le 3V-6$.\end{zadani}
\begin{reseni}Barevnost $\chi(G)$ = nejmenší počet barev pro dobré obarvení (sousedé různě). $K_5$: $V=5$, $E=\binom52=10$; rovinný jednoduchý graf splňuje $E\le3V-6=9$. Protože $10>9$, $K_5$ rovinný \emph{není}. ($\chi(K_5)=5$.)\end{reseni}

\mockq{Otázka 3}{informatika}{Složitost}
\begin{zadani}Definujte třídy P a NP a NP-úplnost. Jak se dokazuje NP-úplnost?\end{zadani}
\begin{reseni}P = řešitelné v polynomiálním čase; NP = ano-instance mají polynomiálně ověřitelný certifikát. $B$ je NP-úplný, je-li v NP a NP-těžký ($\forall K\in$NP$:K\le_p B$). Důkaz: (1) ukázat $B\in$NP, (2) polynomiální redukce ze \emph{známého} NP-úplného problému na $B$.\end{reseni}

\mockq{Otázka 4}{informatika}{Reprezentace dat}
\begin{zadani}Jak je v little-endian uloženo \texttt{int32} \texttt{0x01020304}? Co je zarovnání?\end{zadani}
\begin{reseni}Bajty od nejnižší adresy: \texttt{04 03 02 01} (nejméně významný první). Zarovnání: hodnota velikosti $s$ začíná na adrese dělitelné $s$ (rychlejší přístup CPU); překladač vkládá výplň mezi položky struktury.\end{reseni}

\mockq{Otázka 5}{specializace}{Lineární regrese a L2}
\begin{zadani}Napište normální rovnice OLS a ridge. Jaký má L2 efekt na koeficienty?\end{zadani}
\begin{reseni}OLS: $X^{\!\top}Xw=X^{\!\top}y$, $w=(X^{\!\top}X)^{-1}X^{\!\top}y$ (je-li $X^{\!\top}X$ regulární; jinak pseudoinverze). Ridge: $(X^{\!\top}X+\lambda I)w=X^{\!\top}y$ (bias se obvykle nepenalizuje). L2 koeficienty \emph{smršťuje} k nule (menší rozptyl, řešitelné i při kolinearitě).\end{reseni}

\mockq{Otázka 6}{specializace}{Evaluace klasifikátoru}
\begin{zadani}Z 1000 vzorku je 5\,\% pozitivních; recall $0{,}8$, precision $0{,}5$. Spočtěte TP, FP, FN, TN a $F_1$.\end{zadani}
\begin{reseni}$P=50,N=950$. $TP=0{,}8\cdot50=40$, $FN=10$; z precision $TP+FP=80\Rightarrow FP=40$, $TN=910$. $F_1=\tfrac{2\cdot0{,}5\cdot0{,}8}{1{,}3}\approx0{,}615$. (Accuracy $0{,}95$ je zde zavádějící - nevyvážené třídy.)\end{reseni}

\mockq{Otázka 7}{specializace}{SVM a PCA}
\begin{zadani}Vysvětlete maximální margin a podpůrné vektory. Co PCA optimalizuje a jak souvisí s kovariancí?\end{zadani}
\begin{reseni}SVM hledá nadrovinu s největším marginem ($\min\tfrac12\|w\|^2$ s $y_i(w^{\!\top}x_i+b)\ge1$); určují ji nejbližší body = podpůrné vektory. PCA hledá směry maximálního rozptylu = vlastní vektory kovarianční matice (top-$k$ podle vlastních čísel); ekvivalentně minimalizuje rekonstrukční chybu.\end{reseni}

\mockq{Otázka 8}{specializace}{MDP a Bellman}
\begin{zadani}Napište Bellmanovu rovnici pro $V^*$ a popište iteraci hodnot.\end{zadani}
\begin{reseni}$V^*(s)=\max_a\sum_{s'}P(s'\mid s,a)[R(s,a,s')+\gamma V^*(s')]$. Iterace hodnot: inicializuj $V$, opakovaně aplikuj pravou stranu jako update pro všechny stavy do konvergence; pak $\pi^*(s)=\arg\max_a(\dots)$.\end{reseni}

\begin{infobox}[title={Bodování (orientačně)}]
\footnotesize Společné Q1-Q4: cíl $\ge5$. Specializace Q5-Q8: cíl $\ge7$ (tři ze čtyř na $\ge2$). Hlídej obě podlahy nezávisle.
\end{infobox}

\subsection*{Tvrdý specializační dril}
\addcontentsline{toc}{subsection}{Tvrdý specializační dril}
\begin{infobox}[title={Pokyny}]
\footnotesize 6 specializačních otázek za sebou (numerické). Vázající podlaha je specializace ($\ge7/12$, tři ze čtyř na $\ge2$). Trénuj rychlost a přesnost výpočtu.
\end{infobox}

\mockq{Otázka 1}{specializace}{EM: responsibility}
\begin{zadani}$\pi=(0{,}5,0{,}5)$, $p(x{=}1\mid1)=0{,}7$, $p(x{=}1\mid2)=0{,}3$. Spočtěte responsibility bodu $x{=}1$.\end{zadani}
\begin{reseni}$r_1=\dfrac{0{,}5\cdot0{,}7}{0{,}5\cdot0{,}7+0{,}5\cdot0{,}3}=\dfrac{0{,}35}{0{,}5}=0{,}7$, $r_2=0{,}3$.\end{reseni}

\mockq{Otázka 2}{specializace}{Matice záměn}
\begin{zadani}$TP{=}30,FP{=}10,FN{=}20,TN{=}40$. Spočtěte precision, recall, $F_1$ a accuracy.\end{zadani}
\begin{reseni}precision $=\tfrac{30}{40}=0{,}75$, recall $=\tfrac{30}{50}=0{,}6$, $F_1=\tfrac{2\cdot0{,}75\cdot0{,}6}{1{,}35}\approx0{,}667$, accuracy $=\tfrac{70}{100}=0{,}7$.\end{reseni}

\mockq{Otázka 3}{specializace}{HMM: filtrace}
\begin{zadani}Prior $(0{,}6,0{,}4)$; přechod $P(t\mid t)=0{,}7,P(t\mid f)=0{,}2$; emise $P(o\mid t)=0{,}8,P(o\mid f)=0{,}1$. Spočtěte $P(\cdot\mid o)$.\end{zadani}
\begin{reseni}Predikce: $P(t)=0{,}7\cdot0{,}6+0{,}2\cdot0{,}4=0{,}5$, $P(f)=0{,}5$. Korekce: $\propto(0{,}8\cdot0{,}5,\,0{,}1\cdot0{,}5)=(0{,}4,\,0{,}05)$, norm $/0{,}45\to(0{,}889,\,0{,}111)$.\end{reseni}

\mockq{Otázka 4}{specializace}{Bayes-optimal predikce}
\begin{zadani}$h_1{:}P(+)=0{,}9$ (prior $0{,}4$), $h_2{:}P(+)=0{,}5$ (prior $0{,}6$). Po pozorování ,,$+$'' spočtěte posterior a $P(\text{další}{=}+)$.\end{zadani}
\begin{reseni}Posterior $\propto(0{,}4\cdot0{,}9,\,0{,}6\cdot0{,}5)=(0{,}36,\,0{,}30)$, norm $/0{,}66\to(0{,}545,\,0{,}455)$. Predikce $=0{,}9\cdot0{,}545+0{,}5\cdot0{,}455\approx0{,}718$.\end{reseni}

\mockq{Otázka 5}{specializace}{Ensemble: majoritní hlasování}
\begin{zadani}Tři nezávislé klasifikátory, každý přesnost $0{,}7$, majoritní hlasování. Spočtěte přesnost a uveďte nejlepší/nejhorší případ.\end{zadani}
\begin{reseni}Při nezávislosti $P(\ge2\text{ správně})=\binom32 0{,}7^2\cdot0{,}3+0{,}7^3=3\cdot0{,}147+0{,}343=0{,}784$. Nejlepší případ (chyby se nepřekrývají) až $1$; jsou-li chyby dokonale korelované, ensemble nepomůže a zustává $0{,}7$; při nejnepříznivější závislosti (maximální překryv dvojic chyb) může klesnout až k $0{,}55$. Ensemble pomáhá při \emph{nekorelovaných} chybách.\end{reseni}

\mockq{Otázka 6}{specializace}{Iterace hodnot}
\begin{zadani}Stavy $A,B$, $\gamma=0{,}9$. Z $A$ akce do $B$ (odměna 0), z $B$ akce do cíle (odměna 10). Spočtěte $V^*(A),V^*(B)$.\end{zadani}
\begin{reseni}$V^*(B)=10+0{,}9\cdot0=10$. $V^*(A)=0+0{,}9\cdot V^*(B)=0{,}9\cdot10=9$. (Iterace: $V_1(B)=10$, $V_2(A)=9$.)\end{reseni}

\begin{infobox}[title={Vyhodnocení}]
\footnotesize Ber to jako dva specializační bloky po čtyřech. V každém čtyřbloku miř na $\ge7$: tři otázky vyřešené úplně ($3{\times}2$) plus jedna alespoň částečně ($1$) stačí. Pokud ti numerika padá pod 2 body, zopakuj příslušnou úlohu z úrovní 1-2.
\end{infobox}


\end{document}
